% Options for packages loaded elsewhere
% Options for packages loaded elsewhere
\PassOptionsToPackage{unicode}{hyperref}
\PassOptionsToPackage{hyphens}{url}
%
\documentclass[
  10pt,
  ignorenonframetext,
  aspectratio=169,
  spanish,
]{beamer}
\newif\ifbibliography
\usepackage{pgfpages}
\setbeamertemplate{caption}[numbered]
\setbeamertemplate{caption label separator}{: }
\setbeamercolor{caption name}{fg=normal text.fg}
\beamertemplatenavigationsymbolsempty
% remove section numbering
\setbeamertemplate{part page}{
  \centering
  \begin{beamercolorbox}[sep=16pt,center]{part title}
    \usebeamerfont{part title}\insertpart\par
  \end{beamercolorbox}
}
\setbeamertemplate{section page}{
  \centering
  \begin{beamercolorbox}[sep=12pt,center]{section title}
    \usebeamerfont{section title}\insertsection\par
  \end{beamercolorbox}
}
\setbeamertemplate{subsection page}{
  \centering
  \begin{beamercolorbox}[sep=8pt,center]{subsection title}
    \usebeamerfont{subsection title}\insertsubsection\par
  \end{beamercolorbox}
}
% Prevent slide breaks in the middle of a paragraph
\widowpenalties 1 10000
\raggedbottom
\AtBeginPart{
  \frame{\partpage}
}
\AtBeginSection{
  \ifbibliography
  \else
    \frame{\sectionpage}
  \fi
}
\AtBeginSubsection{
  \frame{\subsectionpage}
}
\usepackage{iftex}
\ifPDFTeX
  \usepackage[T1]{fontenc}
  \usepackage[utf8]{inputenc}
  \usepackage{textcomp} % provide euro and other symbols
\else % if luatex or xetex
  \usepackage{unicode-math} % this also loads fontspec
  \defaultfontfeatures{Scale=MatchLowercase}
  \defaultfontfeatures[\rmfamily]{Ligatures=TeX,Scale=1}
\fi
\usepackage{lmodern}

\usetheme[]{Madrid}
\usecolortheme[]{dolphin}
\usefonttheme[]{professionalfonts}
\usefonttheme{serif} % use mainfont rather than sansfont for slide text
\ifPDFTeX\else
  % xetex/luatex font selection
  \setmainfont[]{DejaVu Sans}
  \setsansfont[]{DejaVu Sans}
  \setmonofont[]{DejaVu Sans Mono}
\fi
% Use upquote if available, for straight quotes in verbatim environments
\IfFileExists{upquote.sty}{\usepackage{upquote}}{}
\IfFileExists{microtype.sty}{% use microtype if available
  \usepackage[]{microtype}
  \UseMicrotypeSet[protrusion]{basicmath} % disable protrusion for tt fonts
}{}
\makeatletter
\@ifundefined{KOMAClassName}{% if non-KOMA class
  \IfFileExists{parskip.sty}{%
    \usepackage{parskip}
  }{% else
    \setlength{\parindent}{0pt}
    \setlength{\parskip}{6pt plus 2pt minus 1pt}}
}{% if KOMA class
  \KOMAoptions{parskip=half}}
\makeatother


\usepackage{longtable,booktabs,array}
\usepackage{calc} % for calculating minipage widths
\usepackage{caption}
% Make caption package work with longtable
\makeatletter
\def\fnum@table{\tablename~\thetable}
\makeatother
\usepackage{graphicx}
\makeatletter
\newsavebox\pandoc@box
\newcommand*\pandocbounded[1]{% scales image to fit in text height/width
  \sbox\pandoc@box{#1}%
  \Gscale@div\@tempa{\textheight}{\dimexpr\ht\pandoc@box+\dp\pandoc@box\relax}%
  \Gscale@div\@tempb{\linewidth}{\wd\pandoc@box}%
  \ifdim\@tempb\p@<\@tempa\p@\let\@tempa\@tempb\fi% select the smaller of both
  \ifdim\@tempa\p@<\p@\scalebox{\@tempa}{\usebox\pandoc@box}%
  \else\usebox{\pandoc@box}%
  \fi%
}
% Set default figure placement to htbp
\def\fps@figure{htbp}
\makeatother



\ifLuaTeX
\usepackage[bidi=basic]{babel}
\else
\usepackage[bidi=default]{babel}
\fi
\ifPDFTeX
\else
\babelfont{rm}[]{DejaVu Sans}
\fi
% get rid of language-specific shorthands (see #6817):
\let\LanguageShortHands\languageshorthands
\def\languageshorthands#1{}


\setlength{\emergencystretch}{3em} % prevent overfull lines

\providecommand{\tightlist}{%
  \setlength{\itemsep}{0pt}\setlength{\parskip}{0pt}}



 


% Estilo separado del contenido. Compilación comprobada con XeLaTeX.
\definecolor{SYSBAzul}{HTML}{0B3954}
\definecolor{SYSBVerde}{HTML}{168C91}
\definecolor{SYSBTexto}{HTML}{203E50}
\setbeamercolor{structure}{fg=SYSBAzul}
\setbeamercolor{normal text}{fg=SYSBTexto,bg=white}
\setbeamercolor{frametitle}{fg=white,bg=SYSBAzul}
\setbeamercolor{block title}{fg=white,bg=SYSBAzul}
\setbeamercolor{block body}{fg=SYSBTexto,bg=SYSBAzul!4}
\setbeamerfont{frametitle}{size=\large,series=\bfseries}
\setbeamerfont{footline}{size=\tiny}
\setbeamersize{text margin left=6mm,text margin right=6mm}
\setbeamertemplate{navigation symbols}{}
\setbeamertemplate{footline}{\hfill\textcolor{SYSBAzul}{SYSB\quad\insertframenumber/\inserttotalframenumber}\hspace{6mm}\vspace{2mm}}
\setbeameroption{hide notes}
\setlength{\parskip}{2pt}
\setlength{\abovedisplayskip}{5pt}
\setlength{\belowdisplayskip}{5pt}
\hypersetup{pdftitle={Sistemas y Señales Biomédicos},pdfauthor={Pablo Eduardo Caicedo-Rodríguez}}

\makeatletter
\@ifpackageloaded{caption}{}{\usepackage{caption}}
\AtBeginDocument{%
\ifdefined\contentsname
  \renewcommand*\contentsname{Tabla de contenidos}
\else
  \newcommand\contentsname{Tabla de contenidos}
\fi
\ifdefined\listfigurename
  \renewcommand*\listfigurename{Listado de Figuras}
\else
  \newcommand\listfigurename{Listado de Figuras}
\fi
\ifdefined\listtablename
  \renewcommand*\listtablename{Listado de Tablas}
\else
  \newcommand\listtablename{Listado de Tablas}
\fi
\ifdefined\figurename
  \renewcommand*\figurename{Figura}
\else
  \newcommand\figurename{Figura}
\fi
\ifdefined\tablename
  \renewcommand*\tablename{Tabla}
\else
  \newcommand\tablename{Tabla}
\fi
}
\@ifpackageloaded{float}{}{\usepackage{float}}
\floatstyle{ruled}
\@ifundefined{c@chapter}{\newfloat{codelisting}{h}{lop}}{\newfloat{codelisting}{h}{lop}[chapter]}
\floatname{codelisting}{Listado}
\newcommand*\listoflistings{\listof{codelisting}{Listado de Listados}}
\makeatother
\makeatletter
\makeatother
\makeatletter
\@ifpackageloaded{caption}{}{\usepackage{caption}}
\@ifpackageloaded{subcaption}{}{\usepackage{subcaption}}
\makeatother

\usepackage{bookmark}
\IfFileExists{xurl.sty}{\usepackage{xurl}}{} % add URL line breaks if available
\urlstyle{same}
\hypersetup{
  pdftitle={Notas del docente · SYSB},
  pdflang={es},
  hidelinks,
  pdfcreator={LaTeX via pandoc}}


\title{Notas del docente · SYSB}
\author{}
\date{}

\begin{document}
\frame{\titlepage}


\begin{frame}
Las notas se agrupan por archivo y número de diapositiva. Las secciones
pertenecen a la tercera edición del texto guía aportado. Use la búsqueda
del navegador para localizar un tema.
\end{frame}

\section{semanas-01-05.qmd}\label{semanas-01-05.qmd}

\begin{frame}{1. Sistemas y Señales Biomédicos}
\phantomsection\label{sistemas-y-seuxf1ales-biomuxe9dicos}
\textbf{Libro guía:} John Semmlow, \emph{Circuits, Signals, and Systems
for Bioengineers}, 3.ª ed., 2018.

§1.1, \textbf{Why Biomedical Engineers Need To Analyze Signals And
Systems} (página 8 del visor PDF); §2.2, \textbf{Time Domain
Measurements} (página 57 del visor PDF); §2.4, \textbf{Time Domain
Analysis} (página 81 del visor PDF); §3.3, \textbf{Time--Frequency
Transformation Of Continuous Signals} (página 118 del visor PDF).

\textbf{Correspondencia:} Síntesis o encuadre que reúne varias
secciones; no corresponde a un único apartado del libro. Lectura
principal: Semmlow (2018), tercera edición. La correspondencia usa la
edición del archivo aportado, que difiere de la segunda edición de 2012
citada en el microcurrículo.

\textbf{Lectura:} use estas secciones como ruta de preparación y
consulta. La bibliografía aparece al final del bloque.
\end{frame}

\begin{frame}{2. Cómo estudiar estas cinco semanas}
\phantomsection\label{cuxf3mo-estudiar-estas-cinco-semanas}
\textbf{Libro guía:} John Semmlow, \emph{Circuits, Signals, and Systems
for Bioengineers}, 3.ª ed., 2018.

§1.1, \textbf{Why Biomedical Engineers Need To Analyze Signals And
Systems} (página 8 del visor PDF); §2.2, \textbf{Time Domain
Measurements} (página 57 del visor PDF); §2.4, \textbf{Time Domain
Analysis} (página 81 del visor PDF); §3.3, \textbf{Time--Frequency
Transformation Of Continuous Signals} (página 118 del visor PDF).

\textbf{Correspondencia:} Síntesis o encuadre que reúne varias
secciones; no corresponde a un único apartado del libro.

\textbf{Lectura:} use estas secciones como ruta de preparación y
consulta. La bibliografía aparece al final del bloque.
\end{frame}

\begin{frame}{3. Una ruta desde el cuerpo hasta el espectro}
\phantomsection\label{una-ruta-desde-el-cuerpo-hasta-el-espectro}
\textbf{Libro guía:} John Semmlow, \emph{Circuits, Signals, and Systems
for Bioengineers}, 3.ª ed., 2018.

§1.1, \textbf{Why Biomedical Engineers Need To Analyze Signals And
Systems} (página 8 del visor PDF); §2.2, \textbf{Time Domain
Measurements} (página 57 del visor PDF); §2.4, \textbf{Time Domain
Analysis} (página 81 del visor PDF); §3.3, \textbf{Time--Frequency
Transformation Of Continuous Signals} (página 118 del visor PDF).

\textbf{Correspondencia:} Síntesis o encuadre que reúne varias
secciones; no corresponde a un único apartado del libro.

\textbf{Lectura:} use estas secciones como ruta de preparación y
consulta. La bibliografía aparece al final del bloque.
\end{frame}

\begin{frame}{4. Mapa de las diez sesiones}
\phantomsection\label{mapa-de-las-diez-sesiones}
\textbf{Libro guía:} John Semmlow, \emph{Circuits, Signals, and Systems
for Bioengineers}, 3.ª ed., 2018.

§1.1, \textbf{Why Biomedical Engineers Need To Analyze Signals And
Systems} (página 8 del visor PDF); §2.2, \textbf{Time Domain
Measurements} (página 57 del visor PDF); §2.4, \textbf{Time Domain
Analysis} (página 81 del visor PDF); §3.3, \textbf{Time--Frequency
Transformation Of Continuous Signals} (página 118 del visor PDF).

\textbf{Correspondencia:} Síntesis o encuadre que reúne varias
secciones; no corresponde a un único apartado del libro.

\textbf{Lectura:} use estas secciones como ruta de preparación y
consulta. La bibliografía aparece al final del bloque.
\end{frame}

\begin{frame}{5. Organización de cada sesión y trabajo independiente}
\phantomsection\label{organizaciuxf3n-de-cada-sesiuxf3n-y-trabajo-independiente}
\textbf{Libro guía:} John Semmlow, \emph{Circuits, Signals, and Systems
for Bioengineers}, 3.ª ed., 2018.

§1.1.1, \textbf{Goals Of This Book} (página 10 del visor PDF).

\textbf{Correspondencia:} Organización docente basada en las actividades
y RAE del microcurrículo. No corresponde a una sección temática del
libro. Microcurrículo: actividades de aprendizaje y RAE, PDF 2--3.
Distribución docente propuesta para las sesiones de 90 minutos.

\textbf{Lectura:} use estas secciones como ruta de preparación y
consulta. La bibliografía aparece al final del bloque.
\end{frame}

\begin{frame}{6. Evidencias de aprendizaje y criterios de revisión}
\phantomsection\label{evidencias-de-aprendizaje-y-criterios-de-revisiuxf3n}
\textbf{Libro guía:} John Semmlow, \emph{Circuits, Signals, and Systems
for Bioengineers}, 3.ª ed., 2018.

§1.1.1, \textbf{Goals Of This Book} (página 10 del visor PDF).

\textbf{Correspondencia:} Organización docente basada en las actividades
y RAE del microcurrículo. No corresponde a una sección temática del
libro. Microcurrículo: RAE y actividades de evaluación, PDF 2--3. No se
resuelven aquí las inconsistencias administrativas del documento.

\textbf{Lectura:} use estas secciones como ruta de preparación y
consulta. La bibliografía aparece al final del bloque.
\end{frame}

\begin{frame}{7. Convenciones de los ejemplos en Python}
\phantomsection\label{convenciones-de-los-ejemplos-en-python}
\textbf{Libro guía:} John Semmlow, \emph{Circuits, Signals, and Systems
for Bioengineers}, 3.ª ed., 2018.

§1.1.1, \textbf{Goals Of This Book} (página 10 del visor PDF).

\textbf{Correspondencia:} Ejemplo, figura o implementación de
elaboración docente a partir de estos fundamentos. No es una
transcripción del código MATLAB del libro. El lenguaje del
microcurrículo suministrado es Python. El libro guía utiliza MATLAB; se
conservan conceptos y se explicitan diferencias de convenciones.

\textbf{Lectura:} use estas secciones como ruta de preparación y
consulta. La bibliografía aparece al final del bloque.
\end{frame}

\begin{frame}{8. Semana 1 · Señales, sistemas y bioseñales}
\phantomsection\label{semana-1-seuxf1ales-sistemas-y-bioseuxf1ales}
\textbf{Libro guía:} John Semmlow, \emph{Circuits, Signals, and Systems
for Bioengineers}, 3.ª ed., 2018.

§1.1, \textbf{Why Biomedical Engineers Need To Analyze Signals And
Systems} (página 8 del visor PDF); §2.2, \textbf{Time Domain
Measurements} (página 57 del visor PDF); §2.4, \textbf{Time Domain
Analysis} (página 81 del visor PDF); §3.3, \textbf{Time--Frequency
Transformation Of Continuous Signals} (página 118 del visor PDF).

\textbf{Correspondencia:} Síntesis o encuadre que reúne varias
secciones; no corresponde a un único apartado del libro.

\textbf{Lectura:} use estas secciones como ruta de preparación y
consulta. La bibliografía aparece al final del bloque.
\end{frame}

\begin{frame}{9. Sesión 1 · La pregunta clínica define qué información
importa}
\phantomsection\label{sesiuxf3n-1-la-pregunta-cluxednica-define-quuxe9-informaciuxf3n-importa}
\textbf{Libro guía:} John Semmlow, \emph{Circuits, Signals, and Systems
for Bioengineers}, 3.ª ed., 2018.

§1.1, \textbf{Why Biomedical Engineers Need To Analyze Signals And
Systems} (página 8 del visor PDF); §1.2.2, \textbf{Biotransducers And
Common Physiological Measurements} (página 12 del visor PDF).

\textbf{Correspondencia:} El tema se desarrolla en las secciones
indicadas. La redacción es una síntesis docente.

\textbf{Microcurrículo:} semana 1, sesión 1 (sesión global 1). Los
numerales del cronograma son identificadores curriculares; no son los
apartados del libro citados arriba.
\end{frame}

\begin{frame}{10. Señal: una variable que porta información}
\phantomsection\label{seuxf1al-una-variable-que-porta-informaciuxf3n}
\textbf{Libro guía:} John Semmlow, \emph{Circuits, Signals, and Systems
for Bioengineers}, 3.ª ed., 2018.

§1.2.3, \textbf{Signal Representation: Continuous (Analog) Signals
Versus Discrete (Digital) Signals} (página 14 del visor PDF).

\textbf{Correspondencia:} El tema se desarrolla en las secciones
indicadas. La redacción es una síntesis docente.

\textbf{Microcurrículo:} semana 1, sesión 1 (sesión global 1). Los
numerales del cronograma son identificadores curriculares; no son los
apartados del libro citados arriba.
\end{frame}

\begin{frame}{11. Sistema: una transformación con propósito}
\phantomsection\label{sistema-una-transformaciuxf3n-con-propuxf3sito}
\textbf{Libro guía:} John Semmlow, \emph{Circuits, Signals, and Systems
for Bioengineers}, 3.ª ed., 2018.

§1.4, \textbf{Biological Systems} (página 34 del visor PDF); §1.4.5,
\textbf{Biosystems Modeling} (página 46 del visor PDF).

\textbf{Correspondencia:} El tema se desarrolla en las secciones
indicadas. La redacción es una síntesis docente.

\textbf{Microcurrículo:} semana 1, sesión 1 (sesión global 1). Los
numerales del cronograma son identificadores curriculares; no son los
apartados del libro citados arriba.
\end{frame}

\begin{frame}{12. Modelo entrada--sistema--salida}
\phantomsection\label{modelo-entradasistemasalida}
\textbf{Libro guía:} John Semmlow, \emph{Circuits, Signals, and Systems
for Bioengineers}, 3.ª ed., 2018.

§1.4, \textbf{Biological Systems} (página 34 del visor PDF); §1.4.5,
\textbf{Biosystems Modeling} (página 46 del visor PDF).

\textbf{Correspondencia:} El tema se desarrolla en las secciones
indicadas. La redacción es una síntesis docente.

\textbf{Microcurrículo:} semana 1, sesión 1 (sesión global 1). Los
numerales del cronograma son identificadores curriculares; no son los
apartados del libro citados arriba.
\end{frame}

\begin{frame}{13. Fuente, estímulo y respuesta}
\phantomsection\label{fuente-estuxedmulo-y-respuesta}
\textbf{Libro guía:} John Semmlow, \emph{Circuits, Signals, and Systems
for Bioengineers}, 3.ª ed., 2018.

§1.4, \textbf{Biological Systems} (página 34 del visor PDF); §1.4.5,
\textbf{Biosystems Modeling} (página 46 del visor PDF).

\textbf{Correspondencia:} El tema se desarrolla en las secciones
indicadas. La redacción es una síntesis docente.

\textbf{Microcurrículo:} semana 1, sesión 1 (sesión global 1). Los
numerales del cronograma son identificadores curriculares; no son los
apartados del libro citados arriba.
\end{frame}

\begin{frame}{14. Jerarquía de una observación biomédica}
\phantomsection\label{jerarquuxeda-de-una-observaciuxf3n-biomuxe9dica}
\textbf{Libro guía:} John Semmlow, \emph{Circuits, Signals, and Systems
for Bioengineers}, 3.ª ed., 2018.

§1.2.2, \textbf{Biotransducers And Common Physiological Measurements}
(página 12 del visor PDF).

\textbf{Correspondencia:} Ejemplo, figura o implementación de
elaboración docente a partir de estos fundamentos. No es una
transcripción del código MATLAB del libro.

\textbf{Microcurrículo:} semana 1, sesión 1 (sesión global 1). Los
numerales del cronograma son identificadores curriculares; no son los
apartados del libro citados arriba.
\end{frame}

\begin{frame}{15. Caso integrador: frecuencia cardíaca desde ECG}
\phantomsection\label{caso-integrador-frecuencia-carduxedaca-desde-ecg}
\textbf{Libro guía:} John Semmlow, \emph{Circuits, Signals, and Systems
for Bioengineers}, 3.ª ed., 2018.

§1.2.2, \textbf{Biotransducers And Common Physiological Measurements}
(página 12 del visor PDF).

\textbf{Correspondencia:} Ejemplo, figura o implementación de
elaboración docente a partir de estos fundamentos. No es una
transcripción del código MATLAB del libro.

\textbf{Microcurrículo:} semana 1, sesión 1 (sesión global 1). Los
numerales del cronograma son identificadores curriculares; no son los
apartados del libro citados arriba.
\end{frame}

\begin{frame}{16. Verificación de la sesión 1}
\phantomsection\label{verificaciuxf3n-de-la-sesiuxf3n-1}
\textbf{Libro guía:} John Semmlow, \emph{Circuits, Signals, and Systems
for Bioengineers}, 3.ª ed., 2018.

§1.2.1, \textbf{Biological Signal Sources} (página 11 del visor PDF);
§1.2.2, \textbf{Biotransducers And Common Physiological Measurements}
(página 12 del visor PDF).

\textbf{Correspondencia:} Actividad o interpretación biomédica de
elaboración docente apoyada en las secciones indicadas.

\textbf{Microcurrículo:} semana 1, sesión 1 (sesión global 1). Los
numerales del cronograma son identificadores curriculares; no son los
apartados del libro citados arriba.
\end{frame}

\begin{frame}{17. Sesión 2 · Las bioseñales expresan distintas formas de
energía}
\phantomsection\label{sesiuxf3n-2-las-bioseuxf1ales-expresan-distintas-formas-de-energuxeda}
\textbf{Libro guía:} John Semmlow, \emph{Circuits, Signals, and Systems
for Bioengineers}, 3.ª ed., 2018.

§1.2.1, \textbf{Biological Signal Sources} (página 11 del visor PDF);
§1.2.2, \textbf{Biotransducers And Common Physiological Measurements}
(página 12 del visor PDF).

\textbf{Correspondencia:} El tema se desarrolla en las secciones
indicadas. La redacción es una síntesis docente.

\textbf{Microcurrículo:} semana 1, sesión 2 (sesión global 2). Los
numerales del cronograma son identificadores curriculares; no son los
apartados del libro citados arriba.
\end{frame}

\begin{frame}{18. Del fenómeno al dato: la cadena de observación}
\phantomsection\label{del-fenuxf3meno-al-dato-la-cadena-de-observaciuxf3n}
\textbf{Libro guía:} John Semmlow, \emph{Circuits, Signals, and Systems
for Bioengineers}, 3.ª ed., 2018.

§1.2.2, \textbf{Biotransducers And Common Physiological Measurements}
(página 12 del visor PDF); §4.2, \textbf{Data Acquisition And Storage}
(página 175 del visor PDF).

\textbf{Correspondencia:} Ejemplo, figura o implementación de
elaboración docente a partir de estos fundamentos. No es una
transcripción del código MATLAB del libro.

\textbf{Microcurrículo:} semana 1, sesión 2 (sesión global 2). Los
numerales del cronograma son identificadores curriculares; no son los
apartados del libro citados arriba.
\end{frame}

\begin{frame}{19. Un modelo aditivo útil, pero incompleto}
\phantomsection\label{un-modelo-aditivo-uxfatil-pero-incompleto}
\textbf{Libro guía:} John Semmlow, \emph{Circuits, Signals, and Systems
for Bioengineers}, 3.ª ed., 2018.

§1.3.1, \textbf{Biological Noise Sources} (página 26 del visor PDF);
§1.3.2, \textbf{Noise Properties: Additive Gaussian Noise} (página 27
del visor PDF).

\textbf{Correspondencia:} El tema se desarrolla en las secciones
indicadas. La redacción es una síntesis docente.

\textbf{Microcurrículo:} semana 1, sesión 2 (sesión global 2). Los
numerales del cronograma son identificadores curriculares; no son los
apartados del libro citados arriba.
\end{frame}

\begin{frame}{20. Ruido, interferencia y artefacto no son sinónimos}
\phantomsection\label{ruido-interferencia-y-artefacto-no-son-sinuxf3nimos}
\textbf{Libro guía:} John Semmlow, \emph{Circuits, Signals, and Systems
for Bioengineers}, 3.ª ed., 2018.

§1.3.1, \textbf{Biological Noise Sources} (página 26 del visor PDF);
§1.3.2, \textbf{Noise Properties: Additive Gaussian Noise} (página 27
del visor PDF).

\textbf{Correspondencia:} El tema se desarrolla en las secciones
indicadas. La redacción es una síntesis docente.

\textbf{Microcurrículo:} semana 1, sesión 2 (sesión global 2). Los
numerales del cronograma son identificadores curriculares; no son los
apartados del libro citados arriba.
\end{frame}

\begin{frame}{21. La relación señal--ruido necesita contexto}
\phantomsection\label{la-relaciuxf3n-seuxf1alruido-necesita-contexto}
\textbf{Libro guía:} John Semmlow, \emph{Circuits, Signals, and Systems
for Bioengineers}, 3.ª ed., 2018.

§1.3.1, \textbf{Biological Noise Sources} (página 26 del visor PDF);
§1.3.2, \textbf{Noise Properties: Additive Gaussian Noise} (página 27
del visor PDF).

\textbf{Correspondencia:} El tema se desarrolla en las secciones
indicadas. La redacción es una síntesis docente.

\textbf{Microcurrículo:} semana 1, sesión 2 (sesión global 2). Los
numerales del cronograma son identificadores curriculares; no son los
apartados del libro citados arriba.
\end{frame}

\begin{frame}{22. Artefactos frecuentes por modalidad}
\phantomsection\label{artefactos-frecuentes-por-modalidad}
\textbf{Libro guía:} John Semmlow, \emph{Circuits, Signals, and Systems
for Bioengineers}, 3.ª ed., 2018.

§1.3.1, \textbf{Biological Noise Sources} (página 26 del visor PDF).

\textbf{Correspondencia:} Actividad o interpretación biomédica de
elaboración docente apoyada en las secciones indicadas.

\textbf{Microcurrículo:} semana 1, sesión 2 (sesión global 2). Los
numerales del cronograma son identificadores curriculares; no son los
apartados del libro citados arriba.
\end{frame}

\begin{frame}{23. ¿Artefacto o fisiología?}
\phantomsection\label{artefacto-o-fisiologuxeda}
\textbf{Libro guía:} John Semmlow, \emph{Circuits, Signals, and Systems
for Bioengineers}, 3.ª ed., 2018.

§1.3.1, \textbf{Biological Noise Sources} (página 26 del visor PDF).

\textbf{Correspondencia:} Actividad o interpretación biomédica de
elaboración docente apoyada en las secciones indicadas.

\textbf{Microcurrículo:} semana 1, sesión 2 (sesión global 2). Los
numerales del cronograma son identificadores curriculares; no son los
apartados del libro citados arriba.
\end{frame}

\begin{frame}{24. Cierre de la semana 1}
\phantomsection\label{cierre-de-la-semana-1}
\textbf{Libro guía:} John Semmlow, \emph{Circuits, Signals, and Systems
for Bioengineers}, 3.ª ed., 2018.

§1.2, \textbf{Biosignal: Signals Produced By Living Systems} (página 11
del visor PDF); §1.3, \textbf{Noise} (página 26 del visor PDF).

\textbf{Correspondencia:} Síntesis o encuadre que reúne varias
secciones; no corresponde a un único apartado del libro.

\textbf{Microcurrículo:} semana 1, sesión 2 (sesión global 2). Los
numerales del cronograma son identificadores curriculares; no son los
apartados del libro citados arriba.
\end{frame}

\begin{frame}{25. Semana 2 · Representación y descripción temporal}
\phantomsection\label{semana-2-representaciuxf3n-y-descripciuxf3n-temporal}
\textbf{Libro guía:} John Semmlow, \emph{Circuits, Signals, and Systems
for Bioengineers}, 3.ª ed., 2018.

§1.1, \textbf{Why Biomedical Engineers Need To Analyze Signals And
Systems} (página 8 del visor PDF); §2.2, \textbf{Time Domain
Measurements} (página 57 del visor PDF); §2.4, \textbf{Time Domain
Analysis} (página 81 del visor PDF); §3.3, \textbf{Time--Frequency
Transformation Of Continuous Signals} (página 118 del visor PDF).

\textbf{Correspondencia:} Síntesis o encuadre que reúne varias
secciones; no corresponde a un único apartado del libro.

\textbf{Lectura:} use estas secciones como ruta de preparación y
consulta. La bibliografía aparece al final del bloque.
\end{frame}

\begin{frame}{26. Sesión 3 · Variable independiente y variable
dependiente}
\phantomsection\label{sesiuxf3n-3-variable-independiente-y-variable-dependiente}
\textbf{Libro guía:} John Semmlow, \emph{Circuits, Signals, and Systems
for Bioengineers}, 3.ª ed., 2018.

§1.2.3, \textbf{Signal Representation: Continuous (Analog) Signals
Versus Discrete (Digital) Signals} (página 14 del visor PDF).

\textbf{Correspondencia:} El tema se desarrolla en las secciones
indicadas. La redacción es una síntesis docente.

\textbf{Microcurrículo:} semana 2, sesión 1 (sesión global 3). Los
numerales del cronograma son identificadores curriculares; no son los
apartados del libro citados arriba.
\end{frame}

\begin{frame}{27. Continuo, discreto, analógico y digital}
\phantomsection\label{continuo-discreto-analuxf3gico-y-digital}
\textbf{Libro guía:} John Semmlow, \emph{Circuits, Signals, and Systems
for Bioengineers}, 3.ª ed., 2018.

§1.2.3, \textbf{Signal Representation: Continuous (Analog) Signals
Versus Discrete (Digital) Signals} (página 14 del visor PDF).

\textbf{Correspondencia:} El tema se desarrolla en las secciones
indicadas. La redacción es una síntesis docente.

\textbf{Microcurrículo:} semana 2, sesión 1 (sesión global 3). Los
numerales del cronograma son identificadores curriculares; no son los
apartados del libro citados arriba.
\end{frame}

\begin{frame}{28. Determinista y aleatoria: una distinción de modelo}
\phantomsection\label{determinista-y-aleatoria-una-distinciuxf3n-de-modelo}
\textbf{Libro guía:} John Semmlow, \emph{Circuits, Signals, and Systems
for Bioengineers}, 3.ª ed., 2018.

§1.4.1, \textbf{Deterministic Versus Stochastic Signals And Systems}
(página 34 del visor PDF).

\textbf{Correspondencia:} El tema se desarrolla en las secciones
indicadas. La redacción es una síntesis docente.

\textbf{Microcurrículo:} semana 2, sesión 1 (sesión global 3). Los
numerales del cronograma son identificadores curriculares; no son los
apartados del libro citados arriba.
\end{frame}

\begin{frame}{29. Media y mediana responden preguntas distintas}
\phantomsection\label{media-y-mediana-responden-preguntas-distintas}
\textbf{Libro guía:} John Semmlow, \emph{Circuits, Signals, and Systems
for Bioengineers}, 3.ª ed., 2018.

§2.2.1, \textbf{Mean And Root Mean Squared Signal Measurements} (página
58 del visor PDF).

\textbf{Correspondencia:} Las secciones aportan el fundamento. La
diapositiva amplía la formulación o precisa condiciones que no se
presentan con idéntica notación en el libro. La media se desarrolla en
§2.2.1; la mediana es un complemento del cronograma.

\textbf{Microcurrículo:} semana 2, sesión 1 (sesión global 3). Los
numerales del cronograma son identificadores curriculares; no son los
apartados del libro citados arriba.
\end{frame}

\begin{frame}{30. Varianza y desviación estándar cuantifican dispersión}
\phantomsection\label{varianza-y-desviaciuxf3n-estuxe1ndar-cuantifican-dispersiuxf3n}
\textbf{Libro guía:} John Semmlow, \emph{Circuits, Signals, and Systems
for Bioengineers}, 3.ª ed., 2018.

§2.2.1, \textbf{Mean And Root Mean Squared Signal Measurements} (página
58 del visor PDF); §2.2.2, \textbf{Variance And Standard Deviation}
(página 64 del visor PDF).

\textbf{Correspondencia:} El tema se desarrolla en las secciones
indicadas. La redacción es una síntesis docente.

\textbf{Microcurrículo:} semana 2, sesión 1 (sesión global 3). Los
numerales del cronograma son identificadores curriculares; no son los
apartados del libro citados arriba.
\end{frame}

\begin{frame}{31. Valor RMS: amplitud efectiva}
\phantomsection\label{valor-rms-amplitud-efectiva}
\textbf{Libro guía:} John Semmlow, \emph{Circuits, Signals, and Systems
for Bioengineers}, 3.ª ed., 2018.

§2.2.1, \textbf{Mean And Root Mean Squared Signal Measurements} (página
58 del visor PDF); §2.2.2, \textbf{Variance And Standard Deviation}
(página 64 del visor PDF).

\textbf{Correspondencia:} El tema se desarrolla en las secciones
indicadas. La redacción es una síntesis docente.

\textbf{Microcurrículo:} semana 2, sesión 1 (sesión global 3). Los
numerales del cronograma son identificadores curriculares; no son los
apartados del libro citados arriba.
\end{frame}

\begin{frame}{32. Energía y potencia media}
\phantomsection\label{energuxeda-y-potencia-media}
\textbf{Libro guía:} John Semmlow, \emph{Circuits, Signals, and Systems
for Bioengineers}, 3.ª ed., 2018.

§2.2.1, \textbf{Mean And Root Mean Squared Signal Measurements} (página
58 del visor PDF); §4.3, \textbf{Power Spectrum} (página 191 del visor
PDF).

\textbf{Correspondencia:} Las secciones aportan el fundamento. La
diapositiva amplía la formulación o precisa condiciones que no se
presentan con idéntica notación en el libro.

\textbf{Microcurrículo:} semana 2, sesión 1 (sesión global 3). Los
numerales del cronograma son identificadores curriculares; no son los
apartados del libro citados arriba.
\end{frame}

\begin{frame}{33. Ejemplo computacional: descriptores explícitos}
\phantomsection\label{ejemplo-computacional-descriptores-expluxedcitos}
\textbf{Libro guía:} John Semmlow, \emph{Circuits, Signals, and Systems
for Bioengineers}, 3.ª ed., 2018.

§2.2.1, \textbf{Mean And Root Mean Squared Signal Measurements} (página
58 del visor PDF); §2.2.2, \textbf{Variance And Standard Deviation}
(página 64 del visor PDF).

\textbf{Correspondencia:} Ejemplo, figura o implementación de
elaboración docente a partir de estos fundamentos. No es una
transcripción del código MATLAB del libro.

\textbf{Microcurrículo:} semana 2, sesión 1 (sesión global 3). Los
numerales del cronograma son identificadores curriculares; no son los
apartados del libro citados arriba.
\end{frame}

\begin{frame}{34. Promediado de realizaciones y reducción de ruido}
\phantomsection\label{promediado-de-realizaciones-y-reducciuxf3n-de-ruido}
\textbf{Libro guía:} John Semmlow, \emph{Circuits, Signals, and Systems
for Bioengineers}, 3.ª ed., 2018.

§2.2.3, \textbf{Averaging For Noise Reduction} (página 68 del visor
PDF); §2.2.4, \textbf{Ensemble Averaging} (página 70 del visor PDF).

\textbf{Correspondencia:} El tema se desarrolla en las secciones
indicadas. La redacción es una síntesis docente.

\textbf{Microcurrículo:} semana 2, sesión 1 (sesión global 3). Los
numerales del cronograma son identificadores curriculares; no son los
apartados del libro citados arriba.
\end{frame}

\begin{frame}{35. Sesión 4 · La sinusoide como coordenada de oscilación}
\phantomsection\label{sesiuxf3n-4-la-sinusoide-como-coordenada-de-oscilaciuxf3n}
\textbf{Libro guía:} John Semmlow, \emph{Circuits, Signals, and Systems
for Bioengineers}, 3.ª ed., 2018.

§2.3.1, \textbf{Sinusoids As Real-Valued Signals} (página 73 del visor
PDF); §6.3, \textbf{The Response Of System Elements To Sinusoidal
Inputs: Phasor Analysis} (página 251 del visor PDF).

\textbf{Correspondencia:} El tema se desarrolla en las secciones
indicadas. La redacción es una síntesis docente.

\textbf{Microcurrículo:} semana 2, sesión 2 (sesión global 4). Los
numerales del cronograma son identificadores curriculares; no son los
apartados del libro citados arriba.
\end{frame}

\begin{frame}{36. Frecuencia no significa rapidez de muestreo}
\phantomsection\label{frecuencia-no-significa-rapidez-de-muestreo}
\textbf{Libro guía:} John Semmlow, \emph{Circuits, Signals, and Systems
for Bioengineers}, 3.ª ed., 2018.

§1.2.3, \textbf{Signal Representation: Continuous (Analog) Signals
Versus Discrete (Digital) Signals} (página 14 del visor PDF); §3.4.3,
\textbf{Details Of The Dft Spectrum} (página 153 del visor PDF).

\textbf{Correspondencia:} El tema se desarrolla en las secciones
indicadas. La redacción es una síntesis docente.

\textbf{Microcurrículo:} semana 2, sesión 2 (sesión global 4). Los
numerales del cronograma son identificadores curriculares; no son los
apartados del libro citados arriba.
\end{frame}

\begin{frame}{37. La fase ubica el ciclo respecto de una referencia}
\phantomsection\label{la-fase-ubica-el-ciclo-respecto-de-una-referencia}
\textbf{Libro guía:} John Semmlow, \emph{Circuits, Signals, and Systems
for Bioengineers}, 3.ª ed., 2018.

§2.3.1, \textbf{Sinusoids As Real-Valued Signals} (página 73 del visor
PDF).

\textbf{Correspondencia:} El tema se desarrolla en las secciones
indicadas. La redacción es una síntesis docente.

\textbf{Microcurrículo:} semana 2, sesión 2 (sesión global 4). Los
numerales del cronograma son identificadores curriculares; no son los
apartados del libro citados arriba.
\end{frame}

\begin{frame}{38. Senoides complejas: representación compacta}
\phantomsection\label{senoides-complejas-representaciuxf3n-compacta}
\textbf{Libro guía:} John Semmlow, \emph{Circuits, Signals, and Systems
for Bioengineers}, 3.ª ed., 2018.

§2.3.2, \textbf{Complex Representation Of Sinusoids} (página 78 del
visor PDF); §3.3.6, \textbf{Complex Representation} (página 132 del
visor PDF).

\textbf{Correspondencia:} El tema se desarrolla en las secciones
indicadas. La redacción es una síntesis docente.

\textbf{Microcurrículo:} semana 2, sesión 2 (sesión global 4). Los
numerales del cronograma son identificadores curriculares; no son los
apartados del libro citados arriba.
\end{frame}

\begin{frame}{39. Una señal real requiere frecuencias conjugadas}
\phantomsection\label{una-seuxf1al-real-requiere-frecuencias-conjugadas}
\textbf{Libro guía:} John Semmlow, \emph{Circuits, Signals, and Systems
for Bioengineers}, 3.ª ed., 2018.

§2.3.2, \textbf{Complex Representation Of Sinusoids} (página 78 del
visor PDF); §3.3.6, \textbf{Complex Representation} (página 132 del
visor PDF).

\textbf{Correspondencia:} El tema se desarrolla en las secciones
indicadas. La redacción es una síntesis docente.

\textbf{Microcurrículo:} semana 2, sesión 2 (sesión global 4). Los
numerales del cronograma son identificadores curriculares; no son los
apartados del libro citados arriba.
\end{frame}

\begin{frame}{40. Verificación de la semana 2}
\phantomsection\label{verificaciuxf3n-de-la-semana-2}
\textbf{Libro guía:} John Semmlow, \emph{Circuits, Signals, and Systems
for Bioengineers}, 3.ª ed., 2018.

§2.2.1, \textbf{Mean And Root Mean Squared Signal Measurements} (página
58 del visor PDF); §2.2.2, \textbf{Variance And Standard Deviation}
(página 64 del visor PDF); §2.3.1, \textbf{Sinusoids As Real-Valued
Signals} (página 73 del visor PDF).

\textbf{Correspondencia:} Síntesis o encuadre que reúne varias
secciones; no corresponde a un único apartado del libro.

\textbf{Microcurrículo:} semana 2, sesión 2 (sesión global 4). Los
numerales del cronograma son identificadores curriculares; no son los
apartados del libro citados arriba.
\end{frame}

\begin{frame}{41. Semana 3 · Operaciones y rasgos temporales}
\phantomsection\label{semana-3-operaciones-y-rasgos-temporales}
\textbf{Libro guía:} John Semmlow, \emph{Circuits, Signals, and Systems
for Bioengineers}, 3.ª ed., 2018.

§1.1, \textbf{Why Biomedical Engineers Need To Analyze Signals And
Systems} (página 8 del visor PDF); §2.2, \textbf{Time Domain
Measurements} (página 57 del visor PDF); §2.4, \textbf{Time Domain
Analysis} (página 81 del visor PDF); §3.3, \textbf{Time--Frequency
Transformation Of Continuous Signals} (página 118 del visor PDF).

\textbf{Correspondencia:} Síntesis o encuadre que reúne varias
secciones; no corresponde a un único apartado del libro.

\textbf{Lectura:} use estas secciones como ruta de preparación y
consulta. La bibliografía aparece al final del bloque.
\end{frame}

\begin{frame}{42. Sesión 5 · Desplazamiento temporal}
\phantomsection\label{sesiuxf3n-5-desplazamiento-temporal}
\textbf{Libro guía:} John Semmlow, \emph{Circuits, Signals, and Systems
for Bioengineers}, 3.ª ed., 2018.

§1.4.2, \textbf{Deterministic Signal Properties: Periodic, Aperiodic,
And Transient} (página 37 del visor PDF); §2.4.4, \textbf{Shifted
Correlations: Cross-Correlation} (página 89 del visor PDF).

\textbf{Correspondencia:} Tema previsto en el cronograma, desarrollado
como ampliación docente. Las referencias del libro son antecedentes
conceptuales. El cronograma prescribe este tema. Las secciones citadas
aportan antecedentes; el libro no dedica una sección específica a cada
técnica de preprocesamiento o detección presentada aquí.

\textbf{Microcurrículo:} semana 3, sesión 1 (sesión global 5). Los
numerales del cronograma son identificadores curriculares; no son los
apartados del libro citados arriba.
\end{frame}

\begin{frame}{43. Escalamiento temporal cambia la velocidad}
\phantomsection\label{escalamiento-temporal-cambia-la-velocidad}
\textbf{Libro guía:} John Semmlow, \emph{Circuits, Signals, and Systems
for Bioengineers}, 3.ª ed., 2018.

§1.4.2, \textbf{Deterministic Signal Properties: Periodic, Aperiodic,
And Transient} (página 37 del visor PDF); §2.4.4, \textbf{Shifted
Correlations: Cross-Correlation} (página 89 del visor PDF).

\textbf{Correspondencia:} Tema previsto en el cronograma, desarrollado
como ampliación docente. Las referencias del libro son antecedentes
conceptuales. El cronograma prescribe este tema. Las secciones citadas
aportan antecedentes; el libro no dedica una sección específica a cada
técnica de preprocesamiento o detección presentada aquí.

\textbf{Microcurrículo:} semana 3, sesión 1 (sesión global 5). Los
numerales del cronograma son identificadores curriculares; no son los
apartados del libro citados arriba.
\end{frame}

\begin{frame}{44. Inversión temporal: una operación matemática}
\phantomsection\label{inversiuxf3n-temporal-una-operaciuxf3n-matemuxe1tica}
\textbf{Libro guía:} John Semmlow, \emph{Circuits, Signals, and Systems
for Bioengineers}, 3.ª ed., 2018.

§1.4.2, \textbf{Deterministic Signal Properties: Periodic, Aperiodic,
And Transient} (página 37 del visor PDF); §2.4.4, \textbf{Shifted
Correlations: Cross-Correlation} (página 89 del visor PDF).

\textbf{Correspondencia:} Tema previsto en el cronograma, desarrollado
como ampliación docente. Las referencias del libro son antecedentes
conceptuales. El cronograma prescribe este tema. Las secciones citadas
aportan antecedentes; el libro no dedica una sección específica a cada
técnica de preprocesamiento o detección presentada aquí.

\textbf{Microcurrículo:} semana 3, sesión 1 (sesión global 5). Los
numerales del cronograma son identificadores curriculares; no son los
apartados del libro citados arriba.
\end{frame}

\begin{frame}{45. El orden de las transformaciones importa}
\phantomsection\label{el-orden-de-las-transformaciones-importa}
\textbf{Libro guía:} John Semmlow, \emph{Circuits, Signals, and Systems
for Bioengineers}, 3.ª ed., 2018.

§1.4.2, \textbf{Deterministic Signal Properties: Periodic, Aperiodic,
And Transient} (página 37 del visor PDF); §2.4.4, \textbf{Shifted
Correlations: Cross-Correlation} (página 89 del visor PDF).

\textbf{Correspondencia:} Tema previsto en el cronograma, desarrollado
como ampliación docente. Las referencias del libro son antecedentes
conceptuales. El cronograma prescribe este tema. Las secciones citadas
aportan antecedentes; el libro no dedica una sección específica a cada
técnica de preprocesamiento o detección presentada aquí.

\textbf{Microcurrículo:} semana 3, sesión 1 (sesión global 5). Los
numerales del cronograma son identificadores curriculares; no son los
apartados del libro citados arriba.
\end{frame}

\begin{frame}{46. Operaciones de amplitud}
\phantomsection\label{operaciones-de-amplitud}
\textbf{Libro guía:} John Semmlow, \emph{Circuits, Signals, and Systems
for Bioengineers}, 3.ª ed., 2018.

§1.4.2, \textbf{Deterministic Signal Properties: Periodic, Aperiodic,
And Transient} (página 37 del visor PDF); §2.4.4, \textbf{Shifted
Correlations: Cross-Correlation} (página 89 del visor PDF).

\textbf{Correspondencia:} Tema previsto en el cronograma, desarrollado
como ampliación docente. Las referencias del libro son antecedentes
conceptuales. El cronograma prescribe este tema. Las secciones citadas
aportan antecedentes; el libro no dedica una sección específica a cada
técnica de preprocesamiento o detección presentada aquí.

\textbf{Microcurrículo:} semana 3, sesión 1 (sesión global 5). Los
numerales del cronograma son identificadores curriculares; no son los
apartados del libro citados arriba.
\end{frame}

\begin{frame}{47. Ventaneo como multiplicación}
\phantomsection\label{ventaneo-como-multiplicaciuxf3n}
\textbf{Libro guía:} John Semmlow, \emph{Circuits, Signals, and Systems
for Bioengineers}, 3.ª ed., 2018.

§4.2.3, \textbf{Data Truncation} (página 184 del visor PDF); §4.2.4,
\textbf{Data Truncation---Window Functions} (página 188 del visor PDF).

\textbf{Correspondencia:} Ejemplo, figura o implementación de
elaboración docente a partir de estos fundamentos. No es una
transcripción del código MATLAB del libro. El cronograma prescribe este
tema. Las secciones citadas aportan antecedentes; el libro no dedica una
sección específica a cada técnica de preprocesamiento o detección
presentada aquí.

\textbf{Microcurrículo:} semana 3, sesión 1 (sesión global 5). Los
numerales del cronograma son identificadores curriculares; no son los
apartados del libro citados arriba.
\end{frame}

\begin{frame}{48. Código: tiempo y amplitud no se transforman igual}
\phantomsection\label{cuxf3digo-tiempo-y-amplitud-no-se-transforman-igual}
\textbf{Libro guía:} John Semmlow, \emph{Circuits, Signals, and Systems
for Bioengineers}, 3.ª ed., 2018.

§1.4.2, \textbf{Deterministic Signal Properties: Periodic, Aperiodic,
And Transient} (página 37 del visor PDF); §2.4.4, \textbf{Shifted
Correlations: Cross-Correlation} (página 89 del visor PDF).

\textbf{Correspondencia:} Tema previsto en el cronograma, desarrollado
como ampliación docente. Las referencias del libro son antecedentes
conceptuales. El cronograma prescribe este tema. Las secciones citadas
aportan antecedentes; el libro no dedica una sección específica a cada
técnica de preprocesamiento o detección presentada aquí.

\textbf{Microcurrículo:} semana 3, sesión 1 (sesión global 5). Los
numerales del cronograma son identificadores curriculares; no son los
apartados del libro citados arriba.
\end{frame}

\begin{frame}{49. Sesión 6 · Un rasgo temporal comprime información}
\phantomsection\label{sesiuxf3n-6-un-rasgo-temporal-comprime-informaciuxf3n}
\textbf{Libro guía:} John Semmlow, \emph{Circuits, Signals, and Systems
for Bioengineers}, 3.ª ed., 2018.

§2.2, \textbf{Time Domain Measurements} (página 57 del visor PDF); §2.4,
\textbf{Time Domain Analysis} (página 81 del visor PDF).

\textbf{Correspondencia:} Tema previsto en el cronograma, desarrollado
como ampliación docente. Las referencias del libro son antecedentes
conceptuales. El cronograma prescribe este tema. Las secciones citadas
aportan antecedentes; el libro no dedica una sección específica a cada
técnica de preprocesamiento o detección presentada aquí.

\textbf{Microcurrículo:} semana 3, sesión 2 (sesión global 6). Los
numerales del cronograma son identificadores curriculares; no son los
apartados del libro citados arriba.
\end{frame}

\begin{frame}{50. Envolvente: describir variaciones lentas de amplitud}
\phantomsection\label{envolvente-describir-variaciones-lentas-de-amplitud}
\textbf{Libro guía:} John Semmlow, \emph{Circuits, Signals, and Systems
for Bioengineers}, 3.ª ed., 2018.

§2.4, \textbf{Time Domain Analysis} (página 81 del visor PDF); §4.6,
\textbf{Time--Frequency Analysis} (página 205 del visor PDF).

\textbf{Correspondencia:} No existe una sección específica del libro
guía para el tema completo de esta diapositiva. Las secciones anteriores
son antecedentes, no referencias directas. La envolvente de Hilbert no
tiene sección propia en el libro guía. Complemento: Semmlow y Griffel
(2014), §6.3.3, The Analytic Signal. La interpretación de envolvente
exige supuestos adicionales.

\textbf{Microcurrículo:} semana 3, sesión 2 (sesión global 6). Los
numerales del cronograma son identificadores curriculares; no son los
apartados del libro citados arriba.
\end{frame}

\begin{frame}{51. Cruces por cero}
\phantomsection\label{cruces-por-cero}
\textbf{Libro guía:} John Semmlow, \emph{Circuits, Signals, and Systems
for Bioengineers}, 3.ª ed., 2018.

§2.2, \textbf{Time Domain Measurements} (página 57 del visor PDF); §2.4,
\textbf{Time Domain Analysis} (página 81 del visor PDF).

\textbf{Correspondencia:} Tema previsto en el cronograma, desarrollado
como ampliación docente. Las referencias del libro son antecedentes
conceptuales. El cronograma prescribe este tema. Las secciones citadas
aportan antecedentes; el libro no dedica una sección específica a cada
técnica de preprocesamiento o detección presentada aquí.

\textbf{Microcurrículo:} semana 3, sesión 2 (sesión global 6). Los
numerales del cronograma son identificadores curriculares; no son los
apartados del libro citados arriba.
\end{frame}

\begin{frame}{52. Máximos y mínimos no equivalen a eventos}
\phantomsection\label{muxe1ximos-y-muxednimos-no-equivalen-a-eventos}
\textbf{Libro guía:} John Semmlow, \emph{Circuits, Signals, and Systems
for Bioengineers}, 3.ª ed., 2018.

§2.2, \textbf{Time Domain Measurements} (página 57 del visor PDF); §2.4,
\textbf{Time Domain Analysis} (página 81 del visor PDF).

\textbf{Correspondencia:} Tema previsto en el cronograma, desarrollado
como ampliación docente. Las referencias del libro son antecedentes
conceptuales. El cronograma prescribe este tema. Las secciones citadas
aportan antecedentes; el libro no dedica una sección específica a cada
técnica de preprocesamiento o detección presentada aquí.

\textbf{Microcurrículo:} semana 3, sesión 2 (sesión global 6). Los
numerales del cronograma son identificadores curriculares; no son los
apartados del libro citados arriba.
\end{frame}

\begin{frame}{53. Intervalos y tasas}
\phantomsection\label{intervalos-y-tasas}
\textbf{Libro guía:} John Semmlow, \emph{Circuits, Signals, and Systems
for Bioengineers}, 3.ª ed., 2018.

§2.2, \textbf{Time Domain Measurements} (página 57 del visor PDF); §2.4,
\textbf{Time Domain Analysis} (página 81 del visor PDF).

\textbf{Correspondencia:} Tema previsto en el cronograma, desarrollado
como ampliación docente. Las referencias del libro son antecedentes
conceptuales. El cronograma prescribe este tema. Las secciones citadas
aportan antecedentes; el libro no dedica una sección específica a cada
técnica de preprocesamiento o detección presentada aquí.

\textbf{Microcurrículo:} semana 3, sesión 2 (sesión global 6). Los
numerales del cronograma son identificadores curriculares; no son los
apartados del libro citados arriba.
\end{frame}

\begin{frame}{54. Código: detección con restricciones}
\phantomsection\label{cuxf3digo-detecciuxf3n-con-restricciones}
\textbf{Libro guía:} John Semmlow, \emph{Circuits, Signals, and Systems
for Bioengineers}, 3.ª ed., 2018.

§2.2, \textbf{Time Domain Measurements} (página 57 del visor PDF); §2.4,
\textbf{Time Domain Analysis} (página 81 del visor PDF).

\textbf{Correspondencia:} Tema previsto en el cronograma, desarrollado
como ampliación docente. Las referencias del libro son antecedentes
conceptuales. El cronograma prescribe este tema. Las secciones citadas
aportan antecedentes; el libro no dedica una sección específica a cada
técnica de preprocesamiento o detección presentada aquí.

\textbf{Microcurrículo:} semana 3, sesión 2 (sesión global 6). Los
numerales del cronograma son identificadores curriculares; no son los
apartados del libro citados arriba.
\end{frame}

\begin{frame}{55. Validación de un detector de eventos}
\phantomsection\label{validaciuxf3n-de-un-detector-de-eventos}
\textbf{Libro guía:} John Semmlow, \emph{Circuits, Signals, and Systems
for Bioengineers}, 3.ª ed., 2018.

§2.4.4, \textbf{Shifted Correlations: Cross-Correlation} (página 89 del
visor PDF).

\textbf{Correspondencia:} No existe una sección específica del libro
guía para el tema completo de esta diapositiva. Las secciones anteriores
son antecedentes, no referencias directas. Antecedente: detección por
semejanza. Complemento para métricas: Semmlow y Griffel (2014), §16.3.
El emparejamiento temporal uno a uno es la adaptación docente para
eventos.

\textbf{Microcurrículo:} semana 3, sesión 2 (sesión global 6). Los
numerales del cronograma son identificadores curriculares; no son los
apartados del libro citados arriba.
\end{frame}

\begin{frame}{56. Cierre de la semana 3}
\phantomsection\label{cierre-de-la-semana-3}
\textbf{Libro guía:} John Semmlow, \emph{Circuits, Signals, and Systems
for Bioengineers}, 3.ª ed., 2018.

§2.2, \textbf{Time Domain Measurements} (página 57 del visor PDF); §2.4,
\textbf{Time Domain Analysis} (página 81 del visor PDF).

\textbf{Correspondencia:} Tema previsto en el cronograma, desarrollado
como ampliación docente. Las referencias del libro son antecedentes
conceptuales. El cronograma prescribe este tema. Las secciones citadas
aportan antecedentes; el libro no dedica una sección específica a cada
técnica de preprocesamiento o detección presentada aquí.

\textbf{Microcurrículo:} semana 3, sesión 2 (sesión global 6). Los
numerales del cronograma son identificadores curriculares; no son los
apartados del libro citados arriba.
\end{frame}

\begin{frame}{57. Correlación como comparación de morfologías}
\phantomsection\label{correlaciuxf3n-como-comparaciuxf3n-de-morfologuxedas}
\textbf{Libro guía:} John Semmlow, \emph{Circuits, Signals, and Systems
for Bioengineers}, 3.ª ed., 2018.

§2.4.1, \textbf{Comparison Through Correlation} (página 81 del visor
PDF); §2.4.3, \textbf{Correlations Between Signals} (página 86 del visor
PDF).

\textbf{Correspondencia:} El tema se desarrolla en las secciones
indicadas. La redacción es una síntesis docente.

\textbf{Microcurrículo:} semana 3, sesión 2 (sesión global 6). Los
numerales del cronograma son identificadores curriculares; no son los
apartados del libro citados arriba.
\end{frame}

\begin{frame}{58. Correlación cruzada y retardo entre canales}
\phantomsection\label{correlaciuxf3n-cruzada-y-retardo-entre-canales}
\textbf{Libro guía:} John Semmlow, \emph{Circuits, Signals, and Systems
for Bioengineers}, 3.ª ed., 2018.

§2.4.4, \textbf{Shifted Correlations: Cross-Correlation} (página 89 del
visor PDF).

\textbf{Correspondencia:} El tema se desarrolla en las secciones
indicadas. La redacción es una síntesis docente.

\textbf{Microcurrículo:} semana 3, sesión 2 (sesión global 6). Los
numerales del cronograma son identificadores curriculares; no son los
apartados del libro citados arriba.
\end{frame}

\begin{frame}{59. Comprobación de un retardo conocido}
\phantomsection\label{comprobaciuxf3n-de-un-retardo-conocido}
\textbf{Libro guía:} John Semmlow, \emph{Circuits, Signals, and Systems
for Bioengineers}, 3.ª ed., 2018.

§2.4.4, \textbf{Shifted Correlations: Cross-Correlation} (página 89 del
visor PDF).

\textbf{Correspondencia:} Ejemplo, figura o implementación de
elaboración docente a partir de estos fundamentos. No es una
transcripción del código MATLAB del libro. Código original en Python del
principio de correlación cruzada.

\textbf{Microcurrículo:} semana 3, sesión 2 (sesión global 6). Los
numerales del cronograma son identificadores curriculares; no son los
apartados del libro citados arriba.
\end{frame}

\begin{frame}{60. Retardo y máximo de correlación}
\phantomsection\label{retardo-y-muxe1ximo-de-correlaciuxf3n}
\textbf{Libro guía:} John Semmlow, \emph{Circuits, Signals, and Systems
for Bioengineers}, 3.ª ed., 2018.

§2.4.4, \textbf{Shifted Correlations: Cross-Correlation} (página 89 del
visor PDF).

\textbf{Correspondencia:} Ejemplo, figura o implementación de
elaboración docente a partir de estos fundamentos. No es una
transcripción del código MATLAB del libro. Figura de cálculo propio con
señales sintéticas. Código reproducible:
codigo/ejemplos\_verificados.py.

\textbf{Microcurrículo:} semana 3, sesión 2 (sesión global 6). Los
numerales del cronograma son identificadores curriculares; no son los
apartados del libro citados arriba.
\end{frame}

\begin{frame}{61. Semana 4 · Segmentación, normalización y límites del
tiempo}
\phantomsection\label{semana-4-segmentaciuxf3n-normalizaciuxf3n-y-luxedmites-del-tiempo}
\textbf{Libro guía:} John Semmlow, \emph{Circuits, Signals, and Systems
for Bioengineers}, 3.ª ed., 2018.

§1.1, \textbf{Why Biomedical Engineers Need To Analyze Signals And
Systems} (página 8 del visor PDF); §2.2, \textbf{Time Domain
Measurements} (página 57 del visor PDF); §2.4, \textbf{Time Domain
Analysis} (página 81 del visor PDF); §3.3, \textbf{Time--Frequency
Transformation Of Continuous Signals} (página 118 del visor PDF).

\textbf{Correspondencia:} Síntesis o encuadre que reúne varias
secciones; no corresponde a un único apartado del libro.

\textbf{Lectura:} use estas secciones como ruta de preparación y
consulta. La bibliografía aparece al final del bloque.
\end{frame}

\begin{frame}{62. Sesión 7 · Segmentar es definir la unidad de análisis}
\phantomsection\label{sesiuxf3n-7-segmentar-es-definir-la-unidad-de-anuxe1lisis}
\textbf{Libro guía:} John Semmlow, \emph{Circuits, Signals, and Systems
for Bioengineers}, 3.ª ed., 2018.

§2.2.4, \textbf{Ensemble Averaging} (página 70 del visor PDF); §2.4.3,
\textbf{Correlations Between Signals} (página 86 del visor PDF).

\textbf{Correspondencia:} Tema previsto en el cronograma, desarrollado
como ampliación docente. Las referencias del libro son antecedentes
conceptuales. El cronograma prescribe este tema. Las secciones citadas
aportan antecedentes; el libro no dedica una sección específica a cada
técnica de preprocesamiento o detección presentada aquí.

\textbf{Microcurrículo:} semana 4, sesión 1 (sesión global 7). Los
numerales del cronograma son identificadores curriculares; no son los
apartados del libro citados arriba.
\end{frame}

\begin{frame}{63. Longitud de ventana: compromiso fundamental}
\phantomsection\label{longitud-de-ventana-compromiso-fundamental}
\textbf{Libro guía:} John Semmlow, \emph{Circuits, Signals, and Systems
for Bioengineers}, 3.ª ed., 2018.

§4.2.3, \textbf{Data Truncation} (página 184 del visor PDF); §4.6,
\textbf{Time--Frequency Analysis} (página 205 del visor PDF).

\textbf{Correspondencia:} Tema previsto en el cronograma, desarrollado
como ampliación docente. Las referencias del libro son antecedentes
conceptuales. El cronograma prescribe este tema. Las secciones citadas
aportan antecedentes; el libro no dedica una sección específica a cada
técnica de preprocesamiento o detección presentada aquí.

\textbf{Microcurrículo:} semana 4, sesión 1 (sesión global 7). Los
numerales del cronograma son identificadores curriculares; no son los
apartados del libro citados arriba.
\end{frame}

\begin{frame}{64. Solapamiento aumenta ejemplos, no información
independiente}
\phantomsection\label{solapamiento-aumenta-ejemplos-no-informaciuxf3n-independiente}
\textbf{Libro guía:} John Semmlow, \emph{Circuits, Signals, and Systems
for Bioengineers}, 3.ª ed., 2018.

§4.4, \textbf{Spectral Averaging} (página 195 del visor PDF).

\textbf{Correspondencia:} Las secciones aportan el fundamento. La
diapositiva amplía la formulación o precisa condiciones que no se
presentan con idéntica notación en el libro. El solapamiento se trata en
promediado espectral. La fuga entre entrenamiento y prueba es una
ampliación para diseños predictivos.

\textbf{Microcurrículo:} semana 4, sesión 1 (sesión global 7). Los
numerales del cronograma son identificadores curriculares; no son los
apartados del libro citados arriba.
\end{frame}

\begin{frame}{65. Segmentar por evento conserva contexto fisiológico}
\phantomsection\label{segmentar-por-evento-conserva-contexto-fisioluxf3gico}
\textbf{Libro guía:} John Semmlow, \emph{Circuits, Signals, and Systems
for Bioengineers}, 3.ª ed., 2018.

§2.2.4, \textbf{Ensemble Averaging} (página 70 del visor PDF).

\textbf{Correspondencia:} Actividad o interpretación biomédica de
elaboración docente apoyada en las secciones indicadas. El cronograma
prescribe este tema. Las secciones citadas aportan antecedentes; el
libro no dedica una sección específica a cada técnica de
preprocesamiento o detección presentada aquí.

\textbf{Microcurrículo:} semana 4, sesión 1 (sesión global 7). Los
numerales del cronograma son identificadores curriculares; no son los
apartados del libro citados arriba.
\end{frame}

\begin{frame}{66. Autocorrelación y repetición temporal}
\phantomsection\label{autocorrelaciuxf3n-y-repeticiuxf3n-temporal}
\textbf{Libro guía:} John Semmlow, \emph{Circuits, Signals, and Systems
for Bioengineers}, 3.ª ed., 2018.

§2.4.5, \textbf{Autocorrelation} (página 104 del visor PDF); §2.4.6,
\textbf{Autocovariance And Cross-Covariance} (página 108 del visor PDF).

\textbf{Correspondencia:} El tema se desarrolla en las secciones
indicadas. La redacción es una síntesis docente.

\textbf{Microcurrículo:} semana 4, sesión 1 (sesión global 7). Los
numerales del cronograma son identificadores curriculares; no son los
apartados del libro citados arriba.
\end{frame}

\begin{frame}{67. Normalización de amplitud}
\phantomsection\label{normalizaciuxf3n-de-amplitud}
\textbf{Libro guía:} John Semmlow, \emph{Circuits, Signals, and Systems
for Bioengineers}, 3.ª ed., 2018.

§2.2.4, \textbf{Ensemble Averaging} (página 70 del visor PDF); §2.4.3,
\textbf{Correlations Between Signals} (página 86 del visor PDF).

\textbf{Correspondencia:} Tema previsto en el cronograma, desarrollado
como ampliación docente. Las referencias del libro son antecedentes
conceptuales. El cronograma prescribe este tema. Las secciones citadas
aportan antecedentes; el libro no dedica una sección específica a cada
técnica de preprocesamiento o detección presentada aquí.

\textbf{Microcurrículo:} semana 4, sesión 1 (sesión global 7). Los
numerales del cronograma son identificadores curriculares; no son los
apartados del libro citados arriba.
\end{frame}

\begin{frame}{68. Normalización fisiológica: ventajas y riesgos}
\phantomsection\label{normalizaciuxf3n-fisioluxf3gica-ventajas-y-riesgos}
\textbf{Libro guía:} John Semmlow, \emph{Circuits, Signals, and Systems
for Bioengineers}, 3.ª ed., 2018.

§2.2.4, \textbf{Ensemble Averaging} (página 70 del visor PDF); §2.4.3,
\textbf{Correlations Between Signals} (página 86 del visor PDF).

\textbf{Correspondencia:} Tema previsto en el cronograma, desarrollado
como ampliación docente. Las referencias del libro son antecedentes
conceptuales. El cronograma prescribe este tema. Las secciones citadas
aportan antecedentes; el libro no dedica una sección específica a cada
técnica de preprocesamiento o detección presentada aquí.

\textbf{Microcurrículo:} semana 4, sesión 1 (sesión global 7). Los
numerales del cronograma son identificadores curriculares; no son los
apartados del libro citados arriba.
\end{frame}

\begin{frame}{69. Evitar fuga de información en normalización}
\phantomsection\label{evitar-fuga-de-informaciuxf3n-en-normalizaciuxf3n}
\textbf{Libro guía:} John Semmlow, \emph{Circuits, Signals, and Systems
for Bioengineers}, 3.ª ed., 2018.

§2.4.3, \textbf{Correlations Between Signals} (página 86 del visor PDF).

\textbf{Correspondencia:} No existe una sección específica del libro
guía para el tema completo de esta diapositiva. Las secciones anteriores
son antecedentes, no referencias directas. No hay sección específica de
partición de conjuntos en el libro guía. Complemento conceptual: Semmlow
y Griffel (2014), §16.1.1 y §16.6. Desarrollo docente sobre separación
por sujeto.

\textbf{Microcurrículo:} semana 4, sesión 1 (sesión global 7). Los
numerales del cronograma son identificadores curriculares; no son los
apartados del libro citados arriba.
\end{frame}

\begin{frame}{70. Lista de control de un segmento analizable}
\phantomsection\label{lista-de-control-de-un-segmento-analizable}
\textbf{Libro guía:} John Semmlow, \emph{Circuits, Signals, and Systems
for Bioengineers}, 3.ª ed., 2018.

§2.2.4, \textbf{Ensemble Averaging} (página 70 del visor PDF); §2.4.3,
\textbf{Correlations Between Signals} (página 86 del visor PDF).

\textbf{Correspondencia:} Tema previsto en el cronograma, desarrollado
como ampliación docente. Las referencias del libro son antecedentes
conceptuales. El cronograma prescribe este tema. Las secciones citadas
aportan antecedentes; el libro no dedica una sección específica a cada
técnica de preprocesamiento o detección presentada aquí.

\textbf{Microcurrículo:} semana 4, sesión 1 (sesión global 7). Los
numerales del cronograma son identificadores curriculares; no son los
apartados del libro citados arriba.
\end{frame}

\begin{frame}{71. Sesión 8 · El dominio temporal puede ocultar
estructura}
\phantomsection\label{sesiuxf3n-8-el-dominio-temporal-puede-ocultar-estructura}
\textbf{Libro guía:} John Semmlow, \emph{Circuits, Signals, and Systems
for Bioengineers}, 3.ª ed., 2018.

§2.4, \textbf{Time Domain Analysis} (página 81 del visor PDF); §3.2,
\textbf{Time--Frequency Domains: General Concepts} (página 117 del visor
PDF).

\textbf{Correspondencia:} El tema se desarrolla en las secciones
indicadas. La redacción es una síntesis docente.

\textbf{Microcurrículo:} semana 4, sesión 2 (sesión global 8). Los
numerales del cronograma son identificadores curriculares; no son los
apartados del libro citados arriba.
\end{frame}

\begin{frame}{72. Frecuencia como tasa de repetición}
\phantomsection\label{frecuencia-como-tasa-de-repeticiuxf3n}
\textbf{Libro guía:} John Semmlow, \emph{Circuits, Signals, and Systems
for Bioengineers}, 3.ª ed., 2018.

§2.3.1, \textbf{Sinusoids As Real-Valued Signals} (página 73 del visor
PDF); §3.3.1, \textbf{Sinusoidal Properties In The Time And Frequency
Domains} (página 118 del visor PDF).

\textbf{Correspondencia:} El tema se desarrolla en las secciones
indicadas. La redacción es una síntesis docente.

\textbf{Microcurrículo:} semana 4, sesión 2 (sesión global 8). Los
numerales del cronograma son identificadores curriculares; no son los
apartados del libro citados arriba.
\end{frame}

\begin{frame}{73. Ortogonalidad: fundamento de la descomposición}
\phantomsection\label{ortogonalidad-fundamento-de-la-descomposiciuxf3n}
\textbf{Libro guía:} John Semmlow, \emph{Circuits, Signals, and Systems
for Bioengineers}, 3.ª ed., 2018.

§2.4.2, \textbf{Orthogonal Signals And Orthogonality} (página 86 del
visor PDF); §3.3.4, \textbf{Finding The Fourier Coefficients} (página
123 del visor PDF).

\textbf{Correspondencia:} El tema se desarrolla en las secciones
indicadas. La redacción es una síntesis docente.

\textbf{Microcurrículo:} semana 4, sesión 2 (sesión global 8). Los
numerales del cronograma son identificadores curriculares; no son los
apartados del libro citados arriba.
\end{frame}

\begin{frame}{74. Una señal compleja puede verse como suma de
oscilaciones}
\phantomsection\label{una-seuxf1al-compleja-puede-verse-como-suma-de-oscilaciones}
\textbf{Libro guía:} John Semmlow, \emph{Circuits, Signals, and Systems
for Bioengineers}, 3.ª ed., 2018.

§3.3.2, \textbf{Sinusoidal Decomposition: The Fourier Series} (página
119 del visor PDF).

\textbf{Correspondencia:} El tema se desarrolla en las secciones
indicadas. La redacción es una síntesis docente.

\textbf{Microcurrículo:} semana 4, sesión 2 (sesión global 8). Los
numerales del cronograma son identificadores curriculares; no son los
apartados del libro citados arriba.
\end{frame}

\begin{frame}{75. Tiempo y frecuencia son representaciones
complementarias}
\phantomsection\label{tiempo-y-frecuencia-son-representaciones-complementarias}
\textbf{Libro guía:} John Semmlow, \emph{Circuits, Signals, and Systems
for Bioengineers}, 3.ª ed., 2018.

§3.2, \textbf{Time--Frequency Domains: General Concepts} (página 117 del
visor PDF); §3.3.7, \textbf{The Continuous Fourier Transform} (página
138 del visor PDF).

\textbf{Correspondencia:} El tema se desarrolla en las secciones
indicadas. La redacción es una síntesis docente.

\textbf{Microcurrículo:} semana 4, sesión 2 (sesión global 8). Los
numerales del cronograma son identificadores curriculares; no son los
apartados del libro citados arriba.
\end{frame}

\begin{frame}{76. Estacionariedad: el supuesto que suele quedar oculto}
\phantomsection\label{estacionariedad-el-supuesto-que-suele-quedar-oculto}
\textbf{Libro guía:} John Semmlow, \emph{Circuits, Signals, and Systems
for Bioengineers}, 3.ª ed., 2018.

§10.2, \textbf{Stochastic Processes: Stationarity And Ergodicity}
(página 450 del visor PDF).

\textbf{Correspondencia:} Las secciones aportan el fundamento. La
diapositiva amplía la formulación o precisa condiciones que no se
presentan con idéntica notación en el libro. La estacionariedad se
desarrolla en el capítulo 10 del libro guía, fuera del núcleo de
capítulos 1--8. Se anticipa para justificar ventanas del cronograma.

\textbf{Microcurrículo:} semana 4, sesión 2 (sesión global 8). Los
numerales del cronograma son identificadores curriculares; no son los
apartados del libro citados arriba.
\end{frame}

\begin{frame}{77. Preparación responsable antes del espectro}
\phantomsection\label{preparaciuxf3n-responsable-antes-del-espectro}
\textbf{Libro guía:} John Semmlow, \emph{Circuits, Signals, and Systems
for Bioengineers}, 3.ª ed., 2018.

§4.2, \textbf{Data Acquisition And Storage} (página 175 del visor PDF);
§4.6, \textbf{Time--Frequency Analysis} (página 205 del visor PDF).

\textbf{Correspondencia:} Actividad o interpretación biomédica de
elaboración docente apoyada en las secciones indicadas.

\textbf{Microcurrículo:} semana 4, sesión 2 (sesión global 8). Los
numerales del cronograma son identificadores curriculares; no son los
apartados del libro citados arriba.
\end{frame}

\begin{frame}{78. Verificación de la semana 4}
\phantomsection\label{verificaciuxf3n-de-la-semana-4}
\textbf{Libro guía:} John Semmlow, \emph{Circuits, Signals, and Systems
for Bioengineers}, 3.ª ed., 2018.

§2.2.4, \textbf{Ensemble Averaging} (página 70 del visor PDF); §4.6,
\textbf{Time--Frequency Analysis} (página 205 del visor PDF); §10.2,
\textbf{Stochastic Processes: Stationarity And Ergodicity} (página 450
del visor PDF).

\textbf{Correspondencia:} Actividad o interpretación biomédica de
elaboración docente apoyada en las secciones indicadas.

\textbf{Microcurrículo:} semana 4, sesión 2 (sesión global 8). Los
numerales del cronograma son identificadores curriculares; no son los
apartados del libro citados arriba.
\end{frame}

\begin{frame}{79. Semana 5 · De la periodicidad al espectro}
\phantomsection\label{semana-5-de-la-periodicidad-al-espectro}
\textbf{Libro guía:} John Semmlow, \emph{Circuits, Signals, and Systems
for Bioengineers}, 3.ª ed., 2018.

§1.1, \textbf{Why Biomedical Engineers Need To Analyze Signals And
Systems} (página 8 del visor PDF); §2.2, \textbf{Time Domain
Measurements} (página 57 del visor PDF); §2.4, \textbf{Time Domain
Analysis} (página 81 del visor PDF); §3.3, \textbf{Time--Frequency
Transformation Of Continuous Signals} (página 118 del visor PDF).

\textbf{Correspondencia:} Síntesis o encuadre que reúne varias
secciones; no corresponde a un único apartado del libro.

\textbf{Lectura:} use estas secciones como ruta de preparación y
consulta. La bibliografía aparece al final del bloque.
\end{frame}

\begin{frame}{80. Sesión 9 · Una señal periódica se repite}
\phantomsection\label{sesiuxf3n-9-una-seuxf1al-periuxf3dica-se-repite}
\textbf{Libro guía:} John Semmlow, \emph{Circuits, Signals, and Systems
for Bioengineers}, 3.ª ed., 2018.

§1.4.2, \textbf{Deterministic Signal Properties: Periodic, Aperiodic,
And Transient} (página 37 del visor PDF); §3.3.2, \textbf{Sinusoidal
Decomposition: The Fourier Series} (página 119 del visor PDF).

\textbf{Correspondencia:} El tema se desarrolla en las secciones
indicadas. La redacción es una síntesis docente.

\textbf{Microcurrículo:} semana 5, sesión 1 (sesión global 9). Los
numerales del cronograma son identificadores curriculares; no son los
apartados del libro citados arriba.
\end{frame}

\begin{frame}{81. Armónicos: múltiplos de la frecuencia fundamental}
\phantomsection\label{armuxf3nicos-muxfaltiplos-de-la-frecuencia-fundamental}
\textbf{Libro guía:} John Semmlow, \emph{Circuits, Signals, and Systems
for Bioengineers}, 3.ª ed., 2018.

§3.3.2, \textbf{Sinusoidal Decomposition: The Fourier Series} (página
119 del visor PDF).

\textbf{Correspondencia:} El tema se desarrolla en las secciones
indicadas. La redacción es una síntesis docente.

\textbf{Microcurrículo:} semana 5, sesión 1 (sesión global 9). Los
numerales del cronograma son identificadores curriculares; no son los
apartados del libro citados arriba.
\end{frame}

\begin{frame}{82. Serie trigonométrica de Fourier}
\phantomsection\label{serie-trigonomuxe9trica-de-fourier}
\textbf{Libro guía:} John Semmlow, \emph{Circuits, Signals, and Systems
for Bioengineers}, 3.ª ed., 2018.

§3.3.4, \textbf{Finding The Fourier Coefficients} (página 123 del visor
PDF).

\textbf{Correspondencia:} El tema se desarrolla en las secciones
indicadas. La redacción es una síntesis docente.

\textbf{Microcurrículo:} semana 5, sesión 1 (sesión global 9). Los
numerales del cronograma son identificadores curriculares; no son los
apartados del libro citados arriba.
\end{frame}

\begin{frame}{83. Los coeficientes cuantifican cada armónico}
\phantomsection\label{los-coeficientes-cuantifican-cada-armuxf3nico}
\textbf{Libro guía:} John Semmlow, \emph{Circuits, Signals, and Systems
for Bioengineers}, 3.ª ed., 2018.

§3.3.4, \textbf{Finding The Fourier Coefficients} (página 123 del visor
PDF).

\textbf{Correspondencia:} El tema se desarrolla en las secciones
indicadas. La redacción es una síntesis docente.

\textbf{Microcurrículo:} semana 5, sesión 1 (sesión global 9). Los
numerales del cronograma son identificadores curriculares; no son los
apartados del libro citados arriba.
\end{frame}

\begin{frame}{84. Forma compleja de la serie}
\phantomsection\label{forma-compleja-de-la-serie}
\textbf{Libro guía:} John Semmlow, \emph{Circuits, Signals, and Systems
for Bioengineers}, 3.ª ed., 2018.

§3.3.6, \textbf{Complex Representation} (página 132 del visor PDF).

\textbf{Correspondencia:} El tema se desarrolla en las secciones
indicadas. La redacción es una síntesis docente.

\textbf{Microcurrículo:} semana 5, sesión 1 (sesión global 9). Los
numerales del cronograma son identificadores curriculares; no son los
apartados del libro citados arriba.
\end{frame}

\begin{frame}{85. Simetría reduce el trabajo}
\phantomsection\label{simetruxeda-reduce-el-trabajo}
\textbf{Libro guía:} John Semmlow, \emph{Circuits, Signals, and Systems
for Bioengineers}, 3.ª ed., 2018.

§3.3.5, \textbf{Symmetry} (página 127 del visor PDF).

\textbf{Correspondencia:} El tema se desarrolla en las secciones
indicadas. La redacción es una síntesis docente.

\textbf{Microcurrículo:} semana 5, sesión 1 (sesión global 9). Los
numerales del cronograma son identificadores curriculares; no son los
apartados del libro citados arriba.
\end{frame}

\begin{frame}{86. Ejemplo conceptual: pulso periódico}
\phantomsection\label{ejemplo-conceptual-pulso-periuxf3dico}
\textbf{Libro guía:} John Semmlow, \emph{Circuits, Signals, and Systems
for Bioengineers}, 3.ª ed., 2018.

§3.3.4, \textbf{Finding The Fourier Coefficients} (página 123 del visor
PDF); §3.4.4, \textbf{The Inverse Discrete Fourier Transform} (página
163 del visor PDF).

\textbf{Correspondencia:} Actividad o interpretación biomédica de
elaboración docente apoyada en las secciones indicadas.

\textbf{Microcurrículo:} semana 5, sesión 1 (sesión global 9). Los
numerales del cronograma son identificadores curriculares; no son los
apartados del libro citados arriba.
\end{frame}

\begin{frame}{87. Parseval conecta tiempo y armónicos}
\phantomsection\label{parseval-conecta-tiempo-y-armuxf3nicos}
\textbf{Libro guía:} John Semmlow, \emph{Circuits, Signals, and Systems
for Bioengineers}, 3.ª ed., 2018.

§3.3.4, \textbf{Finding The Fourier Coefficients} (página 123 del visor
PDF); §4.3, \textbf{Power Spectrum} (página 191 del visor PDF).

\textbf{Correspondencia:} Las secciones aportan el fundamento. La
diapositiva amplía la formulación o precisa condiciones que no se
presentan con idéntica notación en el libro.

\textbf{Microcurrículo:} semana 5, sesión 1 (sesión global 9). Los
numerales del cronograma son identificadores curriculares; no son los
apartados del libro citados arriba.
\end{frame}

\begin{frame}{88. Reconstrucción truncada}
\phantomsection\label{reconstrucciuxf3n-truncada}
\textbf{Libro guía:} John Semmlow, \emph{Circuits, Signals, and Systems
for Bioengineers}, 3.ª ed., 2018.

§3.3.4, \textbf{Finding The Fourier Coefficients} (página 123 del visor
PDF); §3.4.4, \textbf{The Inverse Discrete Fourier Transform} (página
163 del visor PDF).

\textbf{Correspondencia:} Actividad o interpretación biomédica de
elaboración docente apoyada en las secciones indicadas.

\textbf{Microcurrículo:} semana 5, sesión 1 (sesión global 9). Los
numerales del cronograma son identificadores curriculares; no son los
apartados del libro citados arriba.
\end{frame}

\begin{frame}{89. Pulso periódico: coeficientes calculados}
\phantomsection\label{pulso-periuxf3dico-coeficientes-calculados}
\textbf{Libro guía:} John Semmlow, \emph{Circuits, Signals, and Systems
for Bioengineers}, 3.ª ed., 2018.

§3.3.4, \textbf{Finding The Fourier Coefficients} (página 123 del visor
PDF); §3.3.5, \textbf{Symmetry} (página 127 del visor PDF).

\textbf{Correspondencia:} Ejemplo, figura o implementación de
elaboración docente a partir de estos fundamentos. No es una
transcripción del código MATLAB del libro. Ejemplo didáctico original,
no reproducción de un ejercicio del libro.

\textbf{Microcurrículo:} semana 5, sesión 1 (sesión global 9). Los
numerales del cronograma son identificadores curriculares; no son los
apartados del libro citados arriba.
\end{frame}

\begin{frame}{90. Reconstrucción y discontinuidades}
\phantomsection\label{reconstrucciuxf3n-y-discontinuidades}
\textbf{Libro guía:} John Semmlow, \emph{Circuits, Signals, and Systems
for Bioengineers}, 3.ª ed., 2018.

§3.3.4, \textbf{Finding The Fourier Coefficients} (página 123 del visor
PDF); §3.3.5, \textbf{Symmetry} (página 127 del visor PDF).

\textbf{Correspondencia:} Ejemplo, figura o implementación de
elaboración docente a partir de estos fundamentos. No es una
transcripción del código MATLAB del libro. Figura de cálculo propio con
señales sintéticas. Código reproducible:
codigo/ejemplos\_verificados.py.

\textbf{Microcurrículo:} semana 5, sesión 1 (sesión global 9). Los
numerales del cronograma son identificadores curriculares; no son los
apartados del libro citados arriba.
\end{frame}

\begin{frame}{91. Sesión 10 · De la serie a la transformada de Fourier}
\phantomsection\label{sesiuxf3n-10-de-la-serie-a-la-transformada-de-fourier}
\textbf{Libro guía:} John Semmlow, \emph{Circuits, Signals, and Systems
for Bioengineers}, 3.ª ed., 2018.

§3.3.7, \textbf{The Continuous Fourier Transform} (página 138 del visor
PDF).

\textbf{Correspondencia:} El tema se desarrolla en las secciones
indicadas. La redacción es una síntesis docente.

\textbf{Microcurrículo:} semana 5, sesión 2 (sesión global 10). Los
numerales del cronograma son identificadores curriculares; no son los
apartados del libro citados arriba.
\end{frame}

\begin{frame}{92. Magnitud y fase cumplen funciones distintas}
\phantomsection\label{magnitud-y-fase-cumplen-funciones-distintas}
\textbf{Libro guía:} John Semmlow, \emph{Circuits, Signals, and Systems
for Bioengineers}, 3.ª ed., 2018.

§3.3.6, \textbf{Complex Representation} (página 132 del visor PDF);
§3.4.3, \textbf{Details Of The Dft Spectrum} (página 153 del visor PDF).

\textbf{Correspondencia:} El tema se desarrolla en las secciones
indicadas. La redacción es una síntesis docente.

\textbf{Microcurrículo:} semana 5, sesión 2 (sesión global 10). Los
numerales del cronograma son identificadores curriculares; no son los
apartados del libro citados arriba.
\end{frame}

\begin{frame}{93. Simetría para señales reales}
\phantomsection\label{simetruxeda-para-seuxf1ales-reales}
\textbf{Libro guía:} John Semmlow, \emph{Circuits, Signals, and Systems
for Bioengineers}, 3.ª ed., 2018.

§3.3.5, \textbf{Symmetry} (página 127 del visor PDF); §3.3.6,
\textbf{Complex Representation} (página 132 del visor PDF).

\textbf{Correspondencia:} El tema se desarrolla en las secciones
indicadas. La redacción es una síntesis docente.

\textbf{Microcurrículo:} semana 5, sesión 2 (sesión global 10). Los
numerales del cronograma son identificadores curriculares; no son los
apartados del libro citados arriba.
\end{frame}

\begin{frame}{94. Pares esenciales de transformada}
\phantomsection\label{pares-esenciales-de-transformada}
\textbf{Libro guía:} John Semmlow, \emph{Circuits, Signals, and Systems
for Bioengineers}, 3.ª ed., 2018.

§3.3.7, \textbf{The Continuous Fourier Transform} (página 138 del visor
PDF).

\textbf{Correspondencia:} Las secciones aportan el fundamento. La
diapositiva amplía la formulación o precisa condiciones que no se
presentan con idéntica notación en el libro. La diapositiva formaliza
propiedades de la transformada continua. No supone que la tabla o la
notación coincidan literalmente con el libro.

\textbf{Microcurrículo:} semana 5, sesión 2 (sesión global 10). Los
numerales del cronograma son identificadores curriculares; no son los
apartados del libro citados arriba.
\end{frame}

\begin{frame}{95. Duración y ancho de banda se compensan}
\phantomsection\label{duraciuxf3n-y-ancho-de-banda-se-compensan}
\textbf{Libro guía:} John Semmlow, \emph{Circuits, Signals, and Systems
for Bioengineers}, 3.ª ed., 2018.

§3.3.7, \textbf{The Continuous Fourier Transform} (página 138 del visor
PDF).

\textbf{Correspondencia:} Las secciones aportan el fundamento. La
diapositiva amplía la formulación o precisa condiciones que no se
presentan con idéntica notación en el libro. La diapositiva formaliza
propiedades de la transformada continua. No supone que la tabla o la
notación coincidan literalmente con el libro.

\textbf{Microcurrículo:} semana 5, sesión 2 (sesión global 10). Los
numerales del cronograma son identificadores curriculares; no son los
apartados del libro citados arriba.
\end{frame}

\begin{frame}{96. Desplazamiento temporal modifica fase, no magnitud}
\phantomsection\label{desplazamiento-temporal-modifica-fase-no-magnitud}
\textbf{Libro guía:} John Semmlow, \emph{Circuits, Signals, and Systems
for Bioengineers}, 3.ª ed., 2018.

§3.3.7, \textbf{The Continuous Fourier Transform} (página 138 del visor
PDF).

\textbf{Correspondencia:} Las secciones aportan el fundamento. La
diapositiva amplía la formulación o precisa condiciones que no se
presentan con idéntica notación en el libro. La diapositiva formaliza
propiedades de la transformada continua. No supone que la tabla o la
notación coincidan literalmente con el libro.

\textbf{Microcurrículo:} semana 5, sesión 2 (sesión global 10). Los
numerales del cronograma son identificadores curriculares; no son los
apartados del libro citados arriba.
\end{frame}

\begin{frame}{97. Escalamiento temporal modifica el espectro}
\phantomsection\label{escalamiento-temporal-modifica-el-espectro}
\textbf{Libro guía:} John Semmlow, \emph{Circuits, Signals, and Systems
for Bioengineers}, 3.ª ed., 2018.

§3.3.7, \textbf{The Continuous Fourier Transform} (página 138 del visor
PDF).

\textbf{Correspondencia:} Las secciones aportan el fundamento. La
diapositiva amplía la formulación o precisa condiciones que no se
presentan con idéntica notación en el libro. La diapositiva formaliza
propiedades de la transformada continua. No supone que la tabla o la
notación coincidan literalmente con el libro.

\textbf{Microcurrículo:} semana 5, sesión 2 (sesión global 10). Los
numerales del cronograma son identificadores curriculares; no son los
apartados del libro citados arriba.
\end{frame}

\begin{frame}{98. Linealidad y superposición espectral}
\phantomsection\label{linealidad-y-superposiciuxf3n-espectral}
\textbf{Libro guía:} John Semmlow, \emph{Circuits, Signals, and Systems
for Bioengineers}, 3.ª ed., 2018.

§3.3.7, \textbf{The Continuous Fourier Transform} (página 138 del visor
PDF).

\textbf{Correspondencia:} Las secciones aportan el fundamento. La
diapositiva amplía la formulación o precisa condiciones que no se
presentan con idéntica notación en el libro. La diapositiva formaliza
propiedades de la transformada continua. No supone que la tabla o la
notación coincidan literalmente con el libro.

\textbf{Microcurrículo:} semana 5, sesión 2 (sesión global 10). Los
numerales del cronograma son identificadores curriculares; no son los
apartados del libro citados arriba.
\end{frame}

\begin{frame}{99. Código: aproximación discreta para explorar}
\phantomsection\label{cuxf3digo-aproximaciuxf3n-discreta-para-explorar}
\textbf{Libro guía:} John Semmlow, \emph{Circuits, Signals, and Systems
for Bioengineers}, 3.ª ed., 2018.

§3.4.1, \textbf{The Discrete Fourier Transform} (página 142 del visor
PDF); §3.4.2, \textbf{Matlab Implementation Of The Discrete Fourier
Transform} (página 148 del visor PDF).

\textbf{Correspondencia:} Ejemplo, figura o implementación de
elaboración docente a partir de estos fundamentos. No es una
transcripción del código MATLAB del libro.

\textbf{Microcurrículo:} semana 5, sesión 2 (sesión global 10). Los
numerales del cronograma son identificadores curriculares; no son los
apartados del libro citados arriba.
\end{frame}

\begin{frame}{100. Lectura biomédica de un espectro}
\phantomsection\label{lectura-biomuxe9dica-de-un-espectro}
\textbf{Libro guía:} John Semmlow, \emph{Circuits, Signals, and Systems
for Bioengineers}, 3.ª ed., 2018.

§3.4.3, \textbf{Details Of The Dft Spectrum} (página 153 del visor PDF);
§4.2.3, \textbf{Data Truncation} (página 184 del visor PDF).

\textbf{Correspondencia:} Actividad o interpretación biomédica de
elaboración docente apoyada en las secciones indicadas.

\textbf{Microcurrículo:} semana 5, sesión 2 (sesión global 10). Los
numerales del cronograma son identificadores curriculares; no son los
apartados del libro citados arriba.
\end{frame}

\begin{frame}{101. Ejercicio integrador de la semana 5}
\phantomsection\label{ejercicio-integrador-de-la-semana-5}
\textbf{Libro guía:} John Semmlow, \emph{Circuits, Signals, and Systems
for Bioengineers}, 3.ª ed., 2018.

§2.3.1, \textbf{Sinusoids As Real-Valued Signals} (página 73 del visor
PDF); §3.3.6, \textbf{Complex Representation} (página 132 del visor
PDF); §3.3.7, \textbf{The Continuous Fourier Transform} (página 138 del
visor PDF).

\textbf{Correspondencia:} Actividad o interpretación biomédica de
elaboración docente apoyada en las secciones indicadas.

\textbf{Microcurrículo:} semana 5, sesión 2 (sesión global 10). Los
numerales del cronograma son identificadores curriculares; no son los
apartados del libro citados arriba.
\end{frame}

\begin{frame}{102. Síntesis y estudio autónomo}
\phantomsection\label{suxedntesis-y-estudio-autuxf3nomo}
\textbf{Libro guía:} John Semmlow, \emph{Circuits, Signals, and Systems
for Bioengineers}, 3.ª ed., 2018.

§1.1, \textbf{Why Biomedical Engineers Need To Analyze Signals And
Systems} (página 8 del visor PDF); §2.2, \textbf{Time Domain
Measurements} (página 57 del visor PDF); §2.4, \textbf{Time Domain
Analysis} (página 81 del visor PDF); §3.3, \textbf{Time--Frequency
Transformation Of Continuous Signals} (página 118 del visor PDF).

\textbf{Correspondencia:} Síntesis o encuadre que reúne varias
secciones; no corresponde a un único apartado del libro.

\textbf{Lectura:} use estas secciones como ruta de preparación y
consulta. La bibliografía aparece al final del bloque.
\end{frame}

\begin{frame}{103. La cadena completa de razonamiento}
\phantomsection\label{la-cadena-completa-de-razonamiento}
\textbf{Libro guía:} John Semmlow, \emph{Circuits, Signals, and Systems
for Bioengineers}, 3.ª ed., 2018.

§1.1, \textbf{Why Biomedical Engineers Need To Analyze Signals And
Systems} (página 8 del visor PDF); §2.2, \textbf{Time Domain
Measurements} (página 57 del visor PDF); §2.4, \textbf{Time Domain
Analysis} (página 81 del visor PDF); §3.3, \textbf{Time--Frequency
Transformation Of Continuous Signals} (página 118 del visor PDF).

\textbf{Correspondencia:} Síntesis o encuadre que reúne varias
secciones; no corresponde a un único apartado del libro.

\textbf{Lectura:} use estas secciones como ruta de preparación y
consulta. La bibliografía aparece al final del bloque.
\end{frame}

\begin{frame}{104. Errores conceptuales que deben evitarse}
\phantomsection\label{errores-conceptuales-que-deben-evitarse}
\textbf{Libro guía:} John Semmlow, \emph{Circuits, Signals, and Systems
for Bioengineers}, 3.ª ed., 2018.

§1.1, \textbf{Why Biomedical Engineers Need To Analyze Signals And
Systems} (página 8 del visor PDF); §2.2, \textbf{Time Domain
Measurements} (página 57 del visor PDF); §2.4, \textbf{Time Domain
Analysis} (página 81 del visor PDF); §3.3, \textbf{Time--Frequency
Transformation Of Continuous Signals} (página 118 del visor PDF).

\textbf{Correspondencia:} Síntesis o encuadre que reúne varias
secciones; no corresponde a un único apartado del libro.

\textbf{Lectura:} use estas secciones como ruta de preparación y
consulta. La bibliografía aparece al final del bloque.
\end{frame}

\begin{frame}{105. Autoevaluación acumulativa}
\phantomsection\label{autoevaluaciuxf3n-acumulativa}
\textbf{Libro guía:} John Semmlow, \emph{Circuits, Signals, and Systems
for Bioengineers}, 3.ª ed., 2018.

§1.1, \textbf{Why Biomedical Engineers Need To Analyze Signals And
Systems} (página 8 del visor PDF); §2.2, \textbf{Time Domain
Measurements} (página 57 del visor PDF); §2.4, \textbf{Time Domain
Analysis} (página 81 del visor PDF); §3.3, \textbf{Time--Frequency
Transformation Of Continuous Signals} (página 118 del visor PDF).

\textbf{Correspondencia:} Síntesis o encuadre que reúne varias
secciones; no corresponde a un único apartado del libro.

\textbf{Lectura:} use estas secciones como ruta de preparación y
consulta. La bibliografía aparece al final del bloque.
\end{frame}

\begin{frame}{106. Actividad integradora sugerida}
\phantomsection\label{actividad-integradora-sugerida}
\textbf{Libro guía:} John Semmlow, \emph{Circuits, Signals, and Systems
for Bioengineers}, 3.ª ed., 2018.

§1.1, \textbf{Why Biomedical Engineers Need To Analyze Signals And
Systems} (página 8 del visor PDF); §2.2, \textbf{Time Domain
Measurements} (página 57 del visor PDF); §2.4, \textbf{Time Domain
Analysis} (página 81 del visor PDF); §3.3, \textbf{Time--Frequency
Transformation Of Continuous Signals} (página 118 del visor PDF).

\textbf{Correspondencia:} Síntesis o encuadre que reúne varias
secciones; no corresponde a un único apartado del libro.

\textbf{Lectura:} use estas secciones como ruta de preparación y
consulta. La bibliografía aparece al final del bloque.
\end{frame}

\begin{frame}{107. Referencias y alcance}
\phantomsection\label{referencias-y-alcance}
\textbf{Libro guía:} John Semmlow, \emph{Circuits, Signals, and Systems
for Bioengineers}, 3.ª ed., 2018.

§1.1, \textbf{Why Biomedical Engineers Need To Analyze Signals And
Systems} (página 8 del visor PDF); §2.2, \textbf{Time Domain
Measurements} (página 57 del visor PDF); §2.4, \textbf{Time Domain
Analysis} (página 81 del visor PDF); §3.3, \textbf{Time--Frequency
Transformation Of Continuous Signals} (página 118 del visor PDF).

\textbf{Correspondencia:} Síntesis o encuadre que reúne varias
secciones; no corresponde a un único apartado del libro.

\textbf{Lectura:} use estas secciones como ruta de preparación y
consulta. La bibliografía aparece al final del bloque.
\end{frame}

\begin{frame}{108. Cierre: criterio de dominio}
\phantomsection\label{cierre-criterio-de-dominio}
\textbf{Libro guía:} John Semmlow, \emph{Circuits, Signals, and Systems
for Bioengineers}, 3.ª ed., 2018.

§1.1, \textbf{Why Biomedical Engineers Need To Analyze Signals And
Systems} (página 8 del visor PDF); §2.2, \textbf{Time Domain
Measurements} (página 57 del visor PDF); §2.4, \textbf{Time Domain
Analysis} (página 81 del visor PDF); §3.3, \textbf{Time--Frequency
Transformation Of Continuous Signals} (página 118 del visor PDF).

\textbf{Correspondencia:} Síntesis o encuadre que reúne varias
secciones; no corresponde a un único apartado del libro.

\textbf{Lectura:} use estas secciones como ruta de preparación y
consulta. La bibliografía aparece al final del bloque.
\end{frame}

\begin{frame}{109. Referencias}
\phantomsection\label{referencias}
\textbf{Libro guía:} John Semmlow, \emph{Circuits, Signals, and Systems
for Bioengineers}, 3.ª ed., 2018.

§1.1, \textbf{Why Biomedical Engineers Need To Analyze Signals And
Systems} (página 8 del visor PDF); §2.2, \textbf{Time Domain
Measurements} (página 57 del visor PDF); §2.4, \textbf{Time Domain
Analysis} (página 81 del visor PDF); §3.3, \textbf{Time--Frequency
Transformation Of Continuous Signals} (página 118 del visor PDF).

\textbf{Correspondencia:} Síntesis o encuadre que reúne varias
secciones; no corresponde a un único apartado del libro.

\textbf{Lectura:} use estas secciones como ruta de preparación y
consulta. La bibliografía aparece al final del bloque.
\end{frame}

\section{semanas-06-10.qmd}\label{semanas-06-10.qmd}

\begin{frame}{1. Sistemas y Señales Biomédicos}
\phantomsection\label{sistemas-y-seuxf1ales-biomuxe9dicos-1}
\textbf{Libro guía:} John Semmlow, \emph{Circuits, Signals, and Systems
for Bioengineers}, 3.ª ed., 2018.

§3.4, \textbf{Time--Frequency Transformation In The Discrete Domain}
(página 142 del visor PDF); §4.2, \textbf{Data Acquisition And Storage}
(página 175 del visor PDF); §4.6, \textbf{Time--Frequency Analysis}
(página 205 del visor PDF); §5.4, \textbf{Using The Impulse Response To
Find A Systems Output To Any Input---Convolution} (página 219 del visor
PDF).

\textbf{Correspondencia:} Síntesis o encuadre que reúne varias
secciones; no corresponde a un único apartado del libro. Lectura
principal: Semmlow (2018), tercera edición. La correspondencia usa la
edición del archivo aportado, que difiere de la segunda edición de 2012
citada en el microcurrículo.

\textbf{Lectura:} use estas secciones como ruta de preparación y
consulta. La bibliografía aparece al final del bloque.
\end{frame}

\begin{frame}{2. Cómo estudiar las semanas 6--10}
\phantomsection\label{cuxf3mo-estudiar-las-semanas-610}
\textbf{Libro guía:} John Semmlow, \emph{Circuits, Signals, and Systems
for Bioengineers}, 3.ª ed., 2018.

§3.4, \textbf{Time--Frequency Transformation In The Discrete Domain}
(página 142 del visor PDF); §4.2, \textbf{Data Acquisition And Storage}
(página 175 del visor PDF); §4.6, \textbf{Time--Frequency Analysis}
(página 205 del visor PDF); §5.4, \textbf{Using The Impulse Response To
Find A Systems Output To Any Input---Convolution} (página 219 del visor
PDF).

\textbf{Correspondencia:} Síntesis o encuadre que reúne varias
secciones; no corresponde a un único apartado del libro.

\textbf{Lectura:} use estas secciones como ruta de preparación y
consulta. La bibliografía aparece al final del bloque.
\end{frame}

\begin{frame}{3. Del espectro ideal a decisiones reproducibles}
\phantomsection\label{del-espectro-ideal-a-decisiones-reproducibles}
\textbf{Libro guía:} John Semmlow, \emph{Circuits, Signals, and Systems
for Bioengineers}, 3.ª ed., 2018.

§3.4, \textbf{Time--Frequency Transformation In The Discrete Domain}
(página 142 del visor PDF); §4.2, \textbf{Data Acquisition And Storage}
(página 175 del visor PDF); §4.6, \textbf{Time--Frequency Analysis}
(página 205 del visor PDF); §5.4, \textbf{Using The Impulse Response To
Find A Systems Output To Any Input---Convolution} (página 219 del visor
PDF).

\textbf{Correspondencia:} Síntesis o encuadre que reúne varias
secciones; no corresponde a un único apartado del libro.

\textbf{Lectura:} use estas secciones como ruta de preparación y
consulta. La bibliografía aparece al final del bloque.
\end{frame}

\begin{frame}{4. Mapa de las diez sesiones}
\phantomsection\label{mapa-de-las-diez-sesiones-1}
\textbf{Libro guía:} John Semmlow, \emph{Circuits, Signals, and Systems
for Bioengineers}, 3.ª ed., 2018.

§3.4, \textbf{Time--Frequency Transformation In The Discrete Domain}
(página 142 del visor PDF); §4.2, \textbf{Data Acquisition And Storage}
(página 175 del visor PDF); §4.6, \textbf{Time--Frequency Analysis}
(página 205 del visor PDF); §5.4, \textbf{Using The Impulse Response To
Find A Systems Output To Any Input---Convolution} (página 219 del visor
PDF).

\textbf{Correspondencia:} Síntesis o encuadre que reúne varias
secciones; no corresponde a un único apartado del libro.

\textbf{Lectura:} use estas secciones como ruta de preparación y
consulta. La bibliografía aparece al final del bloque.
\end{frame}

\begin{frame}{5. Organización de cada sesión y trabajo independiente}
\phantomsection\label{organizaciuxf3n-de-cada-sesiuxf3n-y-trabajo-independiente-1}
\textbf{Libro guía:} John Semmlow, \emph{Circuits, Signals, and Systems
for Bioengineers}, 3.ª ed., 2018.

§1.1.1, \textbf{Goals Of This Book} (página 10 del visor PDF).

\textbf{Correspondencia:} Organización docente basada en las actividades
y RAE del microcurrículo. No corresponde a una sección temática del
libro. Microcurrículo: actividades de aprendizaje y RAE, PDF 2--3.
Distribución docente propuesta para las sesiones de 90 minutos.

\textbf{Lectura:} use estas secciones como ruta de preparación y
consulta. La bibliografía aparece al final del bloque.
\end{frame}

\begin{frame}{6. Evidencias de aprendizaje y criterios de revisión}
\phantomsection\label{evidencias-de-aprendizaje-y-criterios-de-revisiuxf3n-1}
\textbf{Libro guía:} John Semmlow, \emph{Circuits, Signals, and Systems
for Bioengineers}, 3.ª ed., 2018.

§1.1.1, \textbf{Goals Of This Book} (página 10 del visor PDF).

\textbf{Correspondencia:} Organización docente basada en las actividades
y RAE del microcurrículo. No corresponde a una sección temática del
libro. Microcurrículo: RAE y actividades de evaluación, PDF 2--3. No se
resuelven aquí las inconsistencias administrativas del documento.

\textbf{Lectura:} use estas secciones como ruta de preparación y
consulta. La bibliografía aparece al final del bloque.
\end{frame}

\begin{frame}{7. Convenciones de los ejemplos en Python}
\phantomsection\label{convenciones-de-los-ejemplos-en-python-1}
\textbf{Libro guía:} John Semmlow, \emph{Circuits, Signals, and Systems
for Bioengineers}, 3.ª ed., 2018.

§1.1.1, \textbf{Goals Of This Book} (página 10 del visor PDF).

\textbf{Correspondencia:} Ejemplo, figura o implementación de
elaboración docente a partir de estos fundamentos. No es una
transcripción del código MATLAB del libro. El lenguaje del
microcurrículo suministrado es Python. El libro guía utiliza MATLAB; se
conservan conceptos y se explicitan diferencias de convenciones.

\textbf{Lectura:} use estas secciones como ruta de preparación y
consulta. La bibliografía aparece al final del bloque.
\end{frame}

\begin{frame}{8. Convenciones del bloque}
\phantomsection\label{convenciones-del-bloque}
\textbf{Libro guía:} John Semmlow, \emph{Circuits, Signals, and Systems
for Bioengineers}, 3.ª ed., 2018.

§3.4, \textbf{Time--Frequency Transformation In The Discrete Domain}
(página 142 del visor PDF); §4.2, \textbf{Data Acquisition And Storage}
(página 175 del visor PDF); §4.6, \textbf{Time--Frequency Analysis}
(página 205 del visor PDF); §5.4, \textbf{Using The Impulse Response To
Find A Systems Output To Any Input---Convolution} (página 219 del visor
PDF).

\textbf{Correspondencia:} Síntesis o encuadre que reúne varias
secciones; no corresponde a un único apartado del libro.

\textbf{Lectura:} use estas secciones como ruta de preparación y
consulta. La bibliografía aparece al final del bloque.
\end{frame}

\begin{frame}{9. Semana 6 · Propiedades de Fourier y DFT}
\phantomsection\label{semana-6-propiedades-de-fourier-y-dft}
\textbf{Libro guía:} John Semmlow, \emph{Circuits, Signals, and Systems
for Bioengineers}, 3.ª ed., 2018.

§3.4, \textbf{Time--Frequency Transformation In The Discrete Domain}
(página 142 del visor PDF); §4.2, \textbf{Data Acquisition And Storage}
(página 175 del visor PDF); §4.6, \textbf{Time--Frequency Analysis}
(página 205 del visor PDF); §5.4, \textbf{Using The Impulse Response To
Find A Systems Output To Any Input---Convolution} (página 219 del visor
PDF).

\textbf{Correspondencia:} Síntesis o encuadre que reúne varias
secciones; no corresponde a un único apartado del libro.

\textbf{Lectura:} use estas secciones como ruta de preparación y
consulta. La bibliografía aparece al final del bloque.
\end{frame}

\begin{frame}{10. Sesión 11 · Las propiedades permiten razonar sin
reintegrar}
\phantomsection\label{sesiuxf3n-11-las-propiedades-permiten-razonar-sin-reintegrar}
\textbf{Libro guía:} John Semmlow, \emph{Circuits, Signals, and Systems
for Bioengineers}, 3.ª ed., 2018.

§3.3.7, \textbf{The Continuous Fourier Transform} (página 138 del visor
PDF); §5.6, \textbf{Convolution In The Frequency Domain} (página 237 del
visor PDF).

\textbf{Correspondencia:} Las secciones aportan el fundamento. La
diapositiva amplía la formulación o precisa condiciones que no se
presentan con idéntica notación en el libro. Desarrollo de propiedades
con convención en Hz; las secciones son el fundamento y no una tabla
idéntica a la diapositiva.

\textbf{Microcurrículo:} semana 6, sesión 1 (sesión global 11). Los
numerales del cronograma son identificadores curriculares; no son los
apartados del libro citados arriba.
\end{frame}

\begin{frame}{11. Linealidad conserva la combinación}
\phantomsection\label{linealidad-conserva-la-combinaciuxf3n}
\textbf{Libro guía:} John Semmlow, \emph{Circuits, Signals, and Systems
for Bioengineers}, 3.ª ed., 2018.

§3.3.7, \textbf{The Continuous Fourier Transform} (página 138 del visor
PDF); §5.6, \textbf{Convolution In The Frequency Domain} (página 237 del
visor PDF).

\textbf{Correspondencia:} Las secciones aportan el fundamento. La
diapositiva amplía la formulación o precisa condiciones que no se
presentan con idéntica notación en el libro. Desarrollo de propiedades
con convención en Hz; las secciones son el fundamento y no una tabla
idéntica a la diapositiva.

\textbf{Microcurrículo:} semana 6, sesión 1 (sesión global 11). Los
numerales del cronograma son identificadores curriculares; no son los
apartados del libro citados arriba.
\end{frame}

\begin{frame}{12. Linealidad no separa fuentes automáticamente}
\phantomsection\label{linealidad-no-separa-fuentes-automuxe1ticamente}
\textbf{Libro guía:} John Semmlow, \emph{Circuits, Signals, and Systems
for Bioengineers}, 3.ª ed., 2018.

§3.3.7, \textbf{The Continuous Fourier Transform} (página 138 del visor
PDF); §5.6, \textbf{Convolution In The Frequency Domain} (página 237 del
visor PDF).

\textbf{Correspondencia:} Las secciones aportan el fundamento. La
diapositiva amplía la formulación o precisa condiciones que no se
presentan con idéntica notación en el libro. Desarrollo de propiedades
con convención en Hz; las secciones son el fundamento y no una tabla
idéntica a la diapositiva.

\textbf{Microcurrículo:} semana 6, sesión 1 (sesión global 11). Los
numerales del cronograma son identificadores curriculares; no son los
apartados del libro citados arriba.
\end{frame}

\begin{frame}{13. Un retraso añade fase lineal}
\phantomsection\label{un-retraso-auxf1ade-fase-lineal}
\textbf{Libro guía:} John Semmlow, \emph{Circuits, Signals, and Systems
for Bioengineers}, 3.ª ed., 2018.

§3.3.7, \textbf{The Continuous Fourier Transform} (página 138 del visor
PDF); §5.6, \textbf{Convolution In The Frequency Domain} (página 237 del
visor PDF).

\textbf{Correspondencia:} Las secciones aportan el fundamento. La
diapositiva amplía la formulación o precisa condiciones que no se
presentan con idéntica notación en el libro. Desarrollo de propiedades
con convención en Hz; las secciones son el fundamento y no una tabla
idéntica a la diapositiva.

\textbf{Microcurrículo:} semana 6, sesión 1 (sesión global 11). Los
numerales del cronograma son identificadores curriculares; no son los
apartados del libro citados arriba.
\end{frame}

\begin{frame}{14. La fase lineal representa un retardo constante}
\phantomsection\label{la-fase-lineal-representa-un-retardo-constante}
\textbf{Libro guía:} John Semmlow, \emph{Circuits, Signals, and Systems
for Bioengineers}, 3.ª ed., 2018.

§6.4.1, \textbf{The Spectrum Of A Transfer Function} (página 259 del
visor PDF); §8.4.1, \textbf{Filter Properties} (página 360 del visor
PDF).

\textbf{Correspondencia:} Las secciones aportan el fundamento. La
diapositiva amplía la formulación o precisa condiciones que no se
presentan con idéntica notación en el libro. Formalización del retardo
de grupo a partir de la fase. El criterio de no distorsión también exige
ganancia constante.

\textbf{Microcurrículo:} semana 6, sesión 1 (sesión global 11). Los
numerales del cronograma son identificadores curriculares; no son los
apartados del libro citados arriba.
\end{frame}

\begin{frame}{15. Desplazar en frecuencia equivale a modular}
\phantomsection\label{desplazar-en-frecuencia-equivale-a-modular}
\textbf{Libro guía:} John Semmlow, \emph{Circuits, Signals, and Systems
for Bioengineers}, 3.ª ed., 2018.

§3.3.7, \textbf{The Continuous Fourier Transform} (página 138 del visor
PDF); §5.6, \textbf{Convolution In The Frequency Domain} (página 237 del
visor PDF).

\textbf{Correspondencia:} Las secciones aportan el fundamento. La
diapositiva amplía la formulación o precisa condiciones que no se
presentan con idéntica notación en el libro. Desarrollo de propiedades
con convención en Hz; las secciones son el fundamento y no una tabla
idéntica a la diapositiva.

\textbf{Microcurrículo:} semana 6, sesión 1 (sesión global 11). Los
numerales del cronograma son identificadores curriculares; no son los
apartados del libro citados arriba.
\end{frame}

\begin{frame}{16. El escalamiento temporal modifica ancho y amplitud}
\phantomsection\label{el-escalamiento-temporal-modifica-ancho-y-amplitud}
\textbf{Libro guía:} John Semmlow, \emph{Circuits, Signals, and Systems
for Bioengineers}, 3.ª ed., 2018.

§3.3.7, \textbf{The Continuous Fourier Transform} (página 138 del visor
PDF); §5.6, \textbf{Convolution In The Frequency Domain} (página 237 del
visor PDF).

\textbf{Correspondencia:} Las secciones aportan el fundamento. La
diapositiva amplía la formulación o precisa condiciones que no se
presentan con idéntica notación en el libro. Desarrollo de propiedades
con convención en Hz; las secciones son el fundamento y no una tabla
idéntica a la diapositiva.

\textbf{Microcurrículo:} semana 6, sesión 1 (sesión global 11). Los
numerales del cronograma son identificadores curriculares; no son los
apartados del libro citados arriba.
\end{frame}

\begin{frame}{17. Derivar resalta componentes rápidas}
\phantomsection\label{derivar-resalta-componentes-ruxe1pidas}
\textbf{Libro guía:} John Semmlow, \emph{Circuits, Signals, and Systems
for Bioengineers}, 3.ª ed., 2018.

§6.5.2, \textbf{Derivative Element} (página 264 del visor PDF).

\textbf{Correspondencia:} Las secciones aportan el fundamento. La
diapositiva amplía la formulación o precisa condiciones que no se
presentan con idéntica notación en el libro. Desarrollo de propiedades
con convención en Hz; las secciones son el fundamento y no una tabla
idéntica a la diapositiva.

\textbf{Microcurrículo:} semana 6, sesión 1 (sesión global 11). Los
numerales del cronograma son identificadores curriculares; no son los
apartados del libro citados arriba.
\end{frame}

\begin{frame}{18. Integrar favorece variaciones lentas}
\phantomsection\label{integrar-favorece-variaciones-lentas}
\textbf{Libro guía:} John Semmlow, \emph{Circuits, Signals, and Systems
for Bioengineers}, 3.ª ed., 2018.

§6.5.3, \textbf{Integrator Element} (página 266 del visor PDF); §7.2.2,
\textbf{Calculus Operations In The Laplace Domain} (página 301 del visor
PDF).

\textbf{Correspondencia:} Las secciones aportan el fundamento. La
diapositiva amplía la formulación o precisa condiciones que no se
presentan con idéntica notación en el libro. Desarrollo de propiedades
con convención en Hz; las secciones son el fundamento y no una tabla
idéntica a la diapositiva.

\textbf{Microcurrículo:} semana 6, sesión 1 (sesión global 11). Los
numerales del cronograma son identificadores curriculares; no son los
apartados del libro citados arriba.
\end{frame}

\begin{frame}{19. La convolución se convierte en multiplicación}
\phantomsection\label{la-convoluciuxf3n-se-convierte-en-multiplicaciuxf3n}
\textbf{Libro guía:} John Semmlow, \emph{Circuits, Signals, and Systems
for Bioengineers}, 3.ª ed., 2018.

§5.6, \textbf{Convolution In The Frequency Domain} (página 237 del visor
PDF).

\textbf{Correspondencia:} Las secciones aportan el fundamento. La
diapositiva amplía la formulación o precisa condiciones que no se
presentan con idéntica notación en el libro. Desarrollo de propiedades
con convención en Hz; las secciones son el fundamento y no una tabla
idéntica a la diapositiva.

\textbf{Microcurrículo:} semana 6, sesión 1 (sesión global 11). Los
numerales del cronograma son identificadores curriculares; no son los
apartados del libro citados arriba.
\end{frame}

\begin{frame}{20. Multiplicar en tiempo produce convolución espectral}
\phantomsection\label{multiplicar-en-tiempo-produce-convoluciuxf3n-espectral}
\textbf{Libro guía:} John Semmlow, \emph{Circuits, Signals, and Systems
for Bioengineers}, 3.ª ed., 2018.

§4.2.3, \textbf{Data Truncation} (página 184 del visor PDF); §4.2.4,
\textbf{Data Truncation---Window Functions} (página 188 del visor PDF).

\textbf{Correspondencia:} Las secciones aportan el fundamento. La
diapositiva amplía la formulación o precisa condiciones que no se
presentan con idéntica notación en el libro. Desarrollo de propiedades
con convención en Hz; las secciones son el fundamento y no una tabla
idéntica a la diapositiva.

\textbf{Microcurrículo:} semana 6, sesión 1 (sesión global 11). Los
numerales del cronograma son identificadores curriculares; no son los
apartados del libro citados arriba.
\end{frame}

\begin{frame}{21. Dualidad: tiempo y frecuencia intercambian
estructuras}
\phantomsection\label{dualidad-tiempo-y-frecuencia-intercambian-estructuras}
\textbf{Libro guía:} John Semmlow, \emph{Circuits, Signals, and Systems
for Bioengineers}, 3.ª ed., 2018.

§3.3.7, \textbf{The Continuous Fourier Transform} (página 138 del visor
PDF); §5.6, \textbf{Convolution In The Frequency Domain} (página 237 del
visor PDF).

\textbf{Correspondencia:} Las secciones aportan el fundamento. La
diapositiva amplía la formulación o precisa condiciones que no se
presentan con idéntica notación en el libro. Desarrollo de propiedades
con convención en Hz; las secciones son el fundamento y no una tabla
idéntica a la diapositiva.

\textbf{Microcurrículo:} semana 6, sesión 1 (sesión global 11). Los
numerales del cronograma son identificadores curriculares; no son los
apartados del libro citados arriba.
\end{frame}

\begin{frame}{22. Tabla de predicción rápida}
\phantomsection\label{tabla-de-predicciuxf3n-ruxe1pida}
\textbf{Libro guía:} John Semmlow, \emph{Circuits, Signals, and Systems
for Bioengineers}, 3.ª ed., 2018.

§3.3.7, \textbf{The Continuous Fourier Transform} (página 138 del visor
PDF); §5.6, \textbf{Convolution In The Frequency Domain} (página 237 del
visor PDF).

\textbf{Correspondencia:} Las secciones aportan el fundamento. La
diapositiva amplía la formulación o precisa condiciones que no se
presentan con idéntica notación en el libro. Desarrollo de propiedades
con convención en Hz; las secciones son el fundamento y no una tabla
idéntica a la diapositiva.

\textbf{Microcurrículo:} semana 6, sesión 1 (sesión global 11). Los
numerales del cronograma son identificadores curriculares; no son los
apartados del libro citados arriba.
\end{frame}

\begin{frame}{23. Caso biomédico: tres modificaciones del ECG}
\phantomsection\label{caso-biomuxe9dico-tres-modificaciones-del-ecg}
\textbf{Libro guía:} John Semmlow, \emph{Circuits, Signals, and Systems
for Bioengineers}, 3.ª ed., 2018.

§3.3.7, \textbf{The Continuous Fourier Transform} (página 138 del visor
PDF); §5.6, \textbf{Convolution In The Frequency Domain} (página 237 del
visor PDF).

\textbf{Correspondencia:} Las secciones aportan el fundamento. La
diapositiva amplía la formulación o precisa condiciones que no se
presentan con idéntica notación en el libro. Desarrollo de propiedades
con convención en Hz; las secciones son el fundamento y no una tabla
idéntica a la diapositiva.

\textbf{Microcurrículo:} semana 6, sesión 1 (sesión global 11). Los
numerales del cronograma son identificadores curriculares; no son los
apartados del libro citados arriba.
\end{frame}

\begin{frame}{24. Verificación de la sesión 11}
\phantomsection\label{verificaciuxf3n-de-la-sesiuxf3n-11}
\textbf{Libro guía:} John Semmlow, \emph{Circuits, Signals, and Systems
for Bioengineers}, 3.ª ed., 2018.

§3.3.7, \textbf{The Continuous Fourier Transform} (página 138 del visor
PDF); §5.6, \textbf{Convolution In The Frequency Domain} (página 237 del
visor PDF).

\textbf{Correspondencia:} Las secciones aportan el fundamento. La
diapositiva amplía la formulación o precisa condiciones que no se
presentan con idéntica notación en el libro. Desarrollo de propiedades
con convención en Hz; las secciones son el fundamento y no una tabla
idéntica a la diapositiva.

\textbf{Microcurrículo:} semana 6, sesión 1 (sesión global 11). Los
numerales del cronograma son identificadores curriculares; no son los
apartados del libro citados arriba.
\end{frame}

\begin{frame}{25. Sesión 12 · La DFT analiza un bloque finito}
\phantomsection\label{sesiuxf3n-12-la-dft-analiza-un-bloque-finito}
\textbf{Libro guía:} John Semmlow, \emph{Circuits, Signals, and Systems
for Bioengineers}, 3.ª ed., 2018.

§3.4.1, \textbf{The Discrete Fourier Transform} (página 142 del visor
PDF).

\textbf{Correspondencia:} El tema se desarrolla en las secciones
indicadas. La redacción es una síntesis docente.

\textbf{Microcurrículo:} semana 6, sesión 2 (sesión global 12). Los
numerales del cronograma son identificadores curriculares; no son los
apartados del libro citados arriba.
\end{frame}

\begin{frame}{26. La transformada inversa reconstruye el bloque}
\phantomsection\label{la-transformada-inversa-reconstruye-el-bloque}
\textbf{Libro guía:} John Semmlow, \emph{Circuits, Signals, and Systems
for Bioengineers}, 3.ª ed., 2018.

§3.4.4, \textbf{The Inverse Discrete Fourier Transform} (página 163 del
visor PDF).

\textbf{Correspondencia:} El tema se desarrolla en las secciones
indicadas. La redacción es una síntesis docente.

\textbf{Microcurrículo:} semana 6, sesión 2 (sesión global 12). Los
numerales del cronograma son identificadores curriculares; no son los
apartados del libro citados arriba.
\end{frame}

\begin{frame}{27. Cada índice corresponde a una frecuencia}
\phantomsection\label{cada-uxedndice-corresponde-a-una-frecuencia}
\textbf{Libro guía:} John Semmlow, \emph{Circuits, Signals, and Systems
for Bioengineers}, 3.ª ed., 2018.

§3.4.1, \textbf{The Discrete Fourier Transform} (página 142 del visor
PDF); §3.4.3, \textbf{Details Of The Dft Spectrum} (página 153 del visor
PDF).

\textbf{Correspondencia:} El tema se desarrolla en las secciones
indicadas. La redacción es una síntesis docente.

\textbf{Microcurrículo:} semana 6, sesión 2 (sesión global 12). Los
numerales del cronograma son identificadores curriculares; no son los
apartados del libro citados arriba.
\end{frame}

\begin{frame}{28. Los índices superiores representan frecuencias
negativas}
\phantomsection\label{los-uxedndices-superiores-representan-frecuencias-negativas}
\textbf{Libro guía:} John Semmlow, \emph{Circuits, Signals, and Systems
for Bioengineers}, 3.ª ed., 2018.

§3.4.1, \textbf{The Discrete Fourier Transform} (página 142 del visor
PDF); §3.4.3, \textbf{Details Of The Dft Spectrum} (página 153 del visor
PDF).

\textbf{Correspondencia:} El tema se desarrolla en las secciones
indicadas. La redacción es una síntesis docente.

\textbf{Microcurrículo:} semana 6, sesión 2 (sesión global 12). Los
numerales del cronograma son identificadores curriculares; no son los
apartados del libro citados arriba.
\end{frame}

\begin{frame}{29. Señales reales poseen simetría conjugada}
\phantomsection\label{seuxf1ales-reales-poseen-simetruxeda-conjugada}
\textbf{Libro guía:} John Semmlow, \emph{Circuits, Signals, and Systems
for Bioengineers}, 3.ª ed., 2018.

§3.4.1, \textbf{The Discrete Fourier Transform} (página 142 del visor
PDF); §3.4.3, \textbf{Details Of The Dft Spectrum} (página 153 del visor
PDF).

\textbf{Correspondencia:} El tema se desarrolla en las secciones
indicadas. La redacción es una síntesis docente.

\textbf{Microcurrículo:} semana 6, sesión 2 (sesión global 12). Los
numerales del cronograma son identificadores curriculares; no son los
apartados del libro citados arriba.
\end{frame}

\begin{frame}{30. La FFT es un algoritmo}
\phantomsection\label{la-fft-es-un-algoritmo}
\textbf{Libro guía:} John Semmlow, \emph{Circuits, Signals, and Systems
for Bioengineers}, 3.ª ed., 2018.

§3.4.2, \textbf{Matlab Implementation Of The Discrete Fourier Transform}
(página 148 del visor PDF).

\textbf{Correspondencia:} El tema se desarrolla en las secciones
indicadas. La redacción es una síntesis docente.

\textbf{Microcurrículo:} semana 6, sesión 2 (sesión global 12). Los
numerales del cronograma son identificadores curriculares; no son los
apartados del libro citados arriba.
\end{frame}

\begin{frame}{31. La DFT supone extensión periódica del bloque}
\phantomsection\label{la-dft-supone-extensiuxf3n-periuxf3dica-del-bloque}
\textbf{Libro guía:} John Semmlow, \emph{Circuits, Signals, and Systems
for Bioengineers}, 3.ª ed., 2018.

§4.2.3, \textbf{Data Truncation} (página 184 del visor PDF).

\textbf{Correspondencia:} El tema se desarrolla en las secciones
indicadas. La redacción es una síntesis docente.

\textbf{Microcurrículo:} semana 6, sesión 2 (sesión global 12). Los
numerales del cronograma son identificadores curriculares; no son los
apartados del libro citados arriba.
\end{frame}

\begin{frame}{32. Magnitud bilateral y unilateral}
\phantomsection\label{magnitud-bilateral-y-unilateral}
\textbf{Libro guía:} John Semmlow, \emph{Circuits, Signals, and Systems
for Bioengineers}, 3.ª ed., 2018.

§3.4.3, \textbf{Details Of The Dft Spectrum} (página 153 del visor PDF).

\textbf{Correspondencia:} El tema se desarrolla en las secciones
indicadas. La redacción es una síntesis docente.

\textbf{Microcurrículo:} semana 6, sesión 2 (sesión global 12). Los
numerales del cronograma son identificadores curriculares; no son los
apartados del libro citados arriba.
\end{frame}

\begin{frame}{33. La fase necesita una referencia de amplitud}
\phantomsection\label{la-fase-necesita-una-referencia-de-amplitud}
\textbf{Libro guía:} John Semmlow, \emph{Circuits, Signals, and Systems
for Bioengineers}, 3.ª ed., 2018.

§3.4.3, \textbf{Details Of The Dft Spectrum} (página 153 del visor PDF).

\textbf{Correspondencia:} El tema se desarrolla en las secciones
indicadas. La redacción es una síntesis docente.

\textbf{Microcurrículo:} semana 6, sesión 2 (sesión global 12). Los
numerales del cronograma son identificadores curriculares; no son los
apartados del libro citados arriba.
\end{frame}

\begin{frame}{34. Convolución circular y periodicidad}
\phantomsection\label{convoluciuxf3n-circular-y-periodicidad}
\textbf{Libro guía:} John Semmlow, \emph{Circuits, Signals, and Systems
for Bioengineers}, 3.ª ed., 2018.

§5.6, \textbf{Convolution In The Frequency Domain} (página 237 del visor
PDF).

\textbf{Correspondencia:} Las secciones aportan el fundamento. La
diapositiva amplía la formulación o precisa condiciones que no se
presentan con idéntica notación en el libro. Convolución circular
mediante DFT: extensión del teorema para bloques finitos. El relleno
especificado evita aliasing circular.

\textbf{Microcurrículo:} semana 6, sesión 2 (sesión global 12). Los
numerales del cronograma son identificadores curriculares; no son los
apartados del libro citados arriba.
\end{frame}

\begin{frame}{35. Código: eje y espectro unilateral}
\phantomsection\label{cuxf3digo-eje-y-espectro-unilateral}
\textbf{Libro guía:} John Semmlow, \emph{Circuits, Signals, and Systems
for Bioengineers}, 3.ª ed., 2018.

§3.4.2, \textbf{Matlab Implementation Of The Discrete Fourier Transform}
(página 148 del visor PDF); §3.4.3, \textbf{Details Of The Dft Spectrum}
(página 153 del visor PDF).

\textbf{Correspondencia:} Ejemplo, figura o implementación de
elaboración docente a partir de estos fundamentos. No es una
transcripción del código MATLAB del libro.

\textbf{Microcurrículo:} semana 6, sesión 2 (sesión global 12). Los
numerales del cronograma son identificadores curriculares; no son los
apartados del libro citados arriba.
\end{frame}

\begin{frame}{36. Prueba de amplitud para longitudes par e impar}
\phantomsection\label{prueba-de-amplitud-para-longitudes-par-e-impar}
\textbf{Libro guía:} John Semmlow, \emph{Circuits, Signals, and Systems
for Bioengineers}, 3.ª ed., 2018.

§3.4.2, \textbf{Matlab Implementation Of The Discrete Fourier Transform}
(página 148 del visor PDF); §3.4.3, \textbf{Details Of The Dft Spectrum}
(página 153 del visor PDF).

\textbf{Correspondencia:} Ejemplo, figura o implementación de
elaboración docente a partir de estos fundamentos. No es una
transcripción del código MATLAB del libro.

\textbf{Microcurrículo:} semana 6, sesión 2 (sesión global 12). Los
numerales del cronograma son identificadores curriculares; no son los
apartados del libro citados arriba.
\end{frame}

\begin{frame}{37. Verificación de la sesión 12}
\phantomsection\label{verificaciuxf3n-de-la-sesiuxf3n-12}
\textbf{Libro guía:} John Semmlow, \emph{Circuits, Signals, and Systems
for Bioengineers}, 3.ª ed., 2018.

§3.4.1, \textbf{The Discrete Fourier Transform} (página 142 del visor
PDF); §3.4.2, \textbf{Matlab Implementation Of The Discrete Fourier
Transform} (página 148 del visor PDF); §3.4.3, \textbf{Details Of The
Dft Spectrum} (página 153 del visor PDF).

\textbf{Correspondencia:} El tema se desarrolla en las secciones
indicadas. La redacción es una síntesis docente.

\textbf{Microcurrículo:} semana 6, sesión 2 (sesión global 12). Los
numerales del cronograma son identificadores curriculares; no son los
apartados del libro citados arriba.
\end{frame}

\begin{frame}{38. Cierre de la semana 6}
\phantomsection\label{cierre-de-la-semana-6}
\textbf{Libro guía:} John Semmlow, \emph{Circuits, Signals, and Systems
for Bioengineers}, 3.ª ed., 2018.

§3.3.7, \textbf{The Continuous Fourier Transform} (página 138 del visor
PDF); §3.4, \textbf{Time--Frequency Transformation In The Discrete
Domain} (página 142 del visor PDF).

\textbf{Correspondencia:} Síntesis o encuadre que reúne varias
secciones; no corresponde a un único apartado del libro.

\textbf{Microcurrículo:} semana 6, sesión 2 (sesión global 12). Los
numerales del cronograma son identificadores curriculares; no son los
apartados del libro citados arriba.
\end{frame}

\begin{frame}{39. Semana 7 · Resolución, fuga y adquisición}
\phantomsection\label{semana-7-resoluciuxf3n-fuga-y-adquisiciuxf3n}
\textbf{Libro guía:} John Semmlow, \emph{Circuits, Signals, and Systems
for Bioengineers}, 3.ª ed., 2018.

§3.4, \textbf{Time--Frequency Transformation In The Discrete Domain}
(página 142 del visor PDF); §4.2, \textbf{Data Acquisition And Storage}
(página 175 del visor PDF); §4.6, \textbf{Time--Frequency Analysis}
(página 205 del visor PDF); §5.4, \textbf{Using The Impulse Response To
Find A Systems Output To Any Input---Convolution} (página 219 del visor
PDF).

\textbf{Correspondencia:} Síntesis o encuadre que reúne varias
secciones; no corresponde a un único apartado del libro.

\textbf{Lectura:} use estas secciones como ruta de preparación y
consulta. La bibliografía aparece al final del bloque.
\end{frame}

\begin{frame}{40. Sesión 13 · La resolución proviene de la duración}
\phantomsection\label{sesiuxf3n-13-la-resoluciuxf3n-proviene-de-la-duraciuxf3n}
\textbf{Libro guía:} John Semmlow, \emph{Circuits, Signals, and Systems
for Bioengineers}, 3.ª ed., 2018.

§3.4.3, \textbf{Details Of The Dft Spectrum} (página 153 del visor PDF);
§4.2.3, \textbf{Data Truncation} (página 184 del visor PDF).

\textbf{Correspondencia:} El tema se desarrolla en las secciones
indicadas. La redacción es una síntesis docente.

\textbf{Microcurrículo:} semana 7, sesión 1 (sesión global 13). Los
numerales del cronograma son identificadores curriculares; no son los
apartados del libro citados arriba.
\end{frame}

\begin{frame}{41. Resolución nominal no garantiza separación efectiva}
\phantomsection\label{resoluciuxf3n-nominal-no-garantiza-separaciuxf3n-efectiva}
\textbf{Libro guía:} John Semmlow, \emph{Circuits, Signals, and Systems
for Bioengineers}, 3.ª ed., 2018.

§4.2.4, \textbf{Data Truncation---Window Functions} (página 188 del
visor PDF).

\textbf{Correspondencia:} El tema se desarrolla en las secciones
indicadas. La redacción es una síntesis docente.

\textbf{Microcurrículo:} semana 7, sesión 1 (sesión global 13). Los
numerales del cronograma son identificadores curriculares; no son los
apartados del libro citados arriba.
\end{frame}

\begin{frame}{42. Una frecuencia alineada con un bin produce un caso
especial}
\phantomsection\label{una-frecuencia-alineada-con-un-bin-produce-un-caso-especial}
\textbf{Libro guía:} John Semmlow, \emph{Circuits, Signals, and Systems
for Bioengineers}, 3.ª ed., 2018.

§3.4.1, \textbf{The Discrete Fourier Transform} (página 142 del visor
PDF); §4.2.3, \textbf{Data Truncation} (página 184 del visor PDF).

\textbf{Correspondencia:} El tema se desarrolla en las secciones
indicadas. La redacción es una síntesis docente.

\textbf{Microcurrículo:} semana 7, sesión 1 (sesión global 13). Los
numerales del cronograma son identificadores curriculares; no son los
apartados del libro citados arriba.
\end{frame}

\begin{frame}{43. La fuga distribuye energía entre bins}
\phantomsection\label{la-fuga-distribuye-energuxeda-entre-bins}
\textbf{Libro guía:} John Semmlow, \emph{Circuits, Signals, and Systems
for Bioengineers}, 3.ª ed., 2018.

§4.2.3, \textbf{Data Truncation} (página 184 del visor PDF).

\textbf{Correspondencia:} El tema se desarrolla en las secciones
indicadas. La redacción es una síntesis docente.

\textbf{Microcurrículo:} semana 7, sesión 1 (sesión global 13). Los
numerales del cronograma son identificadores curriculares; no son los
apartados del libro citados arriba.
\end{frame}

\begin{frame}{44. Ventana rectangular: resolución y lóbulos laterales}
\phantomsection\label{ventana-rectangular-resoluciuxf3n-y-luxf3bulos-laterales}
\textbf{Libro guía:} John Semmlow, \emph{Circuits, Signals, and Systems
for Bioengineers}, 3.ª ed., 2018.

§4.2.4, \textbf{Data Truncation---Window Functions} (página 188 del
visor PDF).

\textbf{Correspondencia:} El tema se desarrolla en las secciones
indicadas. La redacción es una síntesis docente.

\textbf{Microcurrículo:} semana 7, sesión 1 (sesión global 13). Los
numerales del cronograma son identificadores curriculares; no son los
apartados del libro citados arriba.
\end{frame}

\begin{frame}{45. Hann reduce lóbulos laterales y ensancha el principal}
\phantomsection\label{hann-reduce-luxf3bulos-laterales-y-ensancha-el-principal}
\textbf{Libro guía:} John Semmlow, \emph{Circuits, Signals, and Systems
for Bioengineers}, 3.ª ed., 2018.

§4.2.4, \textbf{Data Truncation---Window Functions} (página 188 del
visor PDF).

\textbf{Correspondencia:} El tema se desarrolla en las secciones
indicadas. La redacción es una síntesis docente.

\textbf{Microcurrículo:} semana 7, sesión 1 (sesión global 13). Los
numerales del cronograma son identificadores curriculares; no son los
apartados del libro citados arriba.
\end{frame}

\begin{frame}{46. Hamming y Blackman ofrecen otros compromisos}
\phantomsection\label{hamming-y-blackman-ofrecen-otros-compromisos}
\textbf{Libro guía:} John Semmlow, \emph{Circuits, Signals, and Systems
for Bioengineers}, 3.ª ed., 2018.

§4.2.4, \textbf{Data Truncation---Window Functions} (página 188 del
visor PDF).

\textbf{Correspondencia:} El tema se desarrolla en las secciones
indicadas. La redacción es una síntesis docente.

\textbf{Microcurrículo:} semana 7, sesión 1 (sesión global 13). Los
numerales del cronograma son identificadores curriculares; no son los
apartados del libro citados arriba.
\end{frame}

\begin{frame}{47. La ganancia coherente corrige amplitud}
\phantomsection\label{la-ganancia-coherente-corrige-amplitud}
\textbf{Libro guía:} John Semmlow, \emph{Circuits, Signals, and Systems
for Bioengineers}, 3.ª ed., 2018.

§4.2.4, \textbf{Data Truncation---Window Functions} (página 188 del
visor PDF); §4.3, \textbf{Power Spectrum} (página 191 del visor PDF).

\textbf{Correspondencia:} Las secciones aportan el fundamento. La
diapositiva amplía la formulación o precisa condiciones que no se
presentan con idéntica notación en el libro. Ganancia coherente y ENBW
formalizan el escalamiento de ventanas. Consulte también la
documentación oficial de SciPy sobre escalamiento espectral.

\textbf{Microcurrículo:} semana 7, sesión 1 (sesión global 13). Los
numerales del cronograma son identificadores curriculares; no son los
apartados del libro citados arriba.
\end{frame}

\begin{frame}{48. El ancho de banda equivalente afecta el ruido}
\phantomsection\label{el-ancho-de-banda-equivalente-afecta-el-ruido}
\textbf{Libro guía:} John Semmlow, \emph{Circuits, Signals, and Systems
for Bioengineers}, 3.ª ed., 2018.

§4.2.4, \textbf{Data Truncation---Window Functions} (página 188 del
visor PDF); §4.3, \textbf{Power Spectrum} (página 191 del visor PDF).

\textbf{Correspondencia:} Las secciones aportan el fundamento. La
diapositiva amplía la formulación o precisa condiciones que no se
presentan con idéntica notación en el libro. Ganancia coherente y ENBW
formalizan el escalamiento de ventanas. Consulte también la
documentación oficial de SciPy sobre escalamiento espectral.

\textbf{Microcurrículo:} semana 7, sesión 1 (sesión global 13). Los
numerales del cronograma son identificadores curriculares; no son los
apartados del libro citados arriba.
\end{frame}

\begin{frame}{49. Zero padding densifica el eje}
\phantomsection\label{zero-padding-densifica-el-eje}
\textbf{Libro guía:} John Semmlow, \emph{Circuits, Signals, and Systems
for Bioengineers}, 3.ª ed., 2018.

§3.4.3, \textbf{Details Of The Dft Spectrum} (página 153 del visor PDF);
§4.2.3, \textbf{Data Truncation} (página 184 del visor PDF).

\textbf{Correspondencia:} El tema se desarrolla en las secciones
indicadas. La redacción es una síntesis docente.

\textbf{Microcurrículo:} semana 7, sesión 1 (sesión global 13). Los
numerales del cronograma son identificadores curriculares; no son los
apartados del libro citados arriba.
\end{frame}

\begin{frame}{50. Duraciones largas también tienen costos}
\phantomsection\label{duraciones-largas-tambiuxe9n-tienen-costos}
\textbf{Libro guía:} John Semmlow, \emph{Circuits, Signals, and Systems
for Bioengineers}, 3.ª ed., 2018.

§4.2.3, \textbf{Data Truncation} (página 184 del visor PDF); §10.2,
\textbf{Stochastic Processes: Stationarity And Ergodicity} (página 450
del visor PDF).

\textbf{Correspondencia:} El tema se desarrolla en las secciones
indicadas. La redacción es una síntesis docente.

\textbf{Microcurrículo:} semana 7, sesión 1 (sesión global 13). Los
numerales del cronograma son identificadores curriculares; no son los
apartados del libro citados arriba.
\end{frame}

\begin{frame}{51. Ejemplo biomédico: ritmo respiratorio}
\phantomsection\label{ejemplo-biomuxe9dico-ritmo-respiratorio}
\textbf{Libro guía:} John Semmlow, \emph{Circuits, Signals, and Systems
for Bioengineers}, 3.ª ed., 2018.

§4.2.3, \textbf{Data Truncation} (página 184 del visor PDF).

\textbf{Correspondencia:} Actividad o interpretación biomédica de
elaboración docente apoyada en las secciones indicadas.

\textbf{Microcurrículo:} semana 7, sesión 1 (sesión global 13). Los
numerales del cronograma son identificadores curriculares; no son los
apartados del libro citados arriba.
\end{frame}

\begin{frame}{52. Código: comparar ventanas}
\phantomsection\label{cuxf3digo-comparar-ventanas}
\textbf{Libro guía:} John Semmlow, \emph{Circuits, Signals, and Systems
for Bioengineers}, 3.ª ed., 2018.

§4.2.4, \textbf{Data Truncation---Window Functions} (página 188 del
visor PDF).

\textbf{Correspondencia:} Ejemplo, figura o implementación de
elaboración docente a partir de estos fundamentos. No es una
transcripción del código MATLAB del libro.

\textbf{Microcurrículo:} semana 7, sesión 1 (sesión global 13). Los
numerales del cronograma son identificadores curriculares; no son los
apartados del libro citados arriba.
\end{frame}

\begin{frame}{53. Verificación de la sesión 13}
\phantomsection\label{verificaciuxf3n-de-la-sesiuxf3n-13}
\textbf{Libro guía:} John Semmlow, \emph{Circuits, Signals, and Systems
for Bioengineers}, 3.ª ed., 2018.

§4.2.3, \textbf{Data Truncation} (página 184 del visor PDF); §4.2.4,
\textbf{Data Truncation---Window Functions} (página 188 del visor PDF).

\textbf{Correspondencia:} El tema se desarrolla en las secciones
indicadas. La redacción es una síntesis docente.

\textbf{Microcurrículo:} semana 7, sesión 1 (sesión global 13). Los
numerales del cronograma son identificadores curriculares; no son los
apartados del libro citados arriba.
\end{frame}

\begin{frame}{54. Sesión 14 · La cadena de adquisición condiciona el
dato}
\phantomsection\label{sesiuxf3n-14-la-cadena-de-adquisiciuxf3n-condiciona-el-dato}
\textbf{Libro guía:} John Semmlow, \emph{Circuits, Signals, and Systems
for Bioengineers}, 3.ª ed., 2018.

§1.2.2, \textbf{Biotransducers And Common Physiological Measurements}
(página 12 del visor PDF); §4.2.1, \textbf{Data Sampling: The Sampling
Theorem} (página 175 del visor PDF); §4.2.2, \textbf{Amplitude
Slicing---Quantization} (página 182 del visor PDF).

\textbf{Correspondencia:} El tema se desarrolla en las secciones
indicadas. La redacción es una síntesis docente.

\textbf{Microcurrículo:} semana 7, sesión 2 (sesión global 14). Los
numerales del cronograma son identificadores curriculares; no son los
apartados del libro citados arriba.
\end{frame}

\begin{frame}{55. El sensor define sensibilidad y ancho de banda}
\phantomsection\label{el-sensor-define-sensibilidad-y-ancho-de-banda}
\textbf{Libro guía:} John Semmlow, \emph{Circuits, Signals, and Systems
for Bioengineers}, 3.ª ed., 2018.

§1.2.2, \textbf{Biotransducers And Common Physiological Measurements}
(página 12 del visor PDF).

\textbf{Correspondencia:} Actividad o interpretación biomédica de
elaboración docente apoyada en las secciones indicadas.

\textbf{Microcurrículo:} semana 7, sesión 2 (sesión global 14). Los
numerales del cronograma son identificadores curriculares; no son los
apartados del libro citados arriba.
\end{frame}

\begin{frame}{56. El acoplamiento forma parte de la medición}
\phantomsection\label{el-acoplamiento-forma-parte-de-la-mediciuxf3n}
\textbf{Libro guía:} John Semmlow, \emph{Circuits, Signals, and Systems
for Bioengineers}, 3.ª ed., 2018.

§1.2.2, \textbf{Biotransducers And Common Physiological Measurements}
(página 12 del visor PDF).

\textbf{Correspondencia:} Actividad o interpretación biomédica de
elaboración docente apoyada en las secciones indicadas.

\textbf{Microcurrículo:} semana 7, sesión 2 (sesión global 14). Los
numerales del cronograma son identificadores curriculares; no son los
apartados del libro citados arriba.
\end{frame}

\begin{frame}{57. El acondicionamiento protege y adapta la señal}
\phantomsection\label{el-acondicionamiento-protege-y-adapta-la-seuxf1al}
\textbf{Libro guía:} John Semmlow, \emph{Circuits, Signals, and Systems
for Bioengineers}, 3.ª ed., 2018.

§1.2.2, \textbf{Biotransducers And Common Physiological Measurements}
(página 12 del visor PDF).

\textbf{Correspondencia:} Actividad o interpretación biomédica de
elaboración docente apoyada en las secciones indicadas.

\textbf{Microcurrículo:} semana 7, sesión 2 (sesión global 14). Los
numerales del cronograma son identificadores curriculares; no son los
apartados del libro citados arriba.
\end{frame}

\begin{frame}{58. El filtro antialias actúa antes del ADC}
\phantomsection\label{el-filtro-antialias-actuxfaa-antes-del-adc}
\textbf{Libro guía:} John Semmlow, \emph{Circuits, Signals, and Systems
for Bioengineers}, 3.ª ed., 2018.

§4.2.1, \textbf{Data Sampling: The Sampling Theorem} (página 175 del
visor PDF).

\textbf{Correspondencia:} El tema se desarrolla en las secciones
indicadas. La redacción es una síntesis docente.

\textbf{Microcurrículo:} semana 7, sesión 2 (sesión global 14). Los
numerales del cronograma son identificadores curriculares; no son los
apartados del libro citados arriba.
\end{frame}

\begin{frame}{59. El ADC discretiza tiempo y amplitud}
\phantomsection\label{el-adc-discretiza-tiempo-y-amplitud}
\textbf{Libro guía:} John Semmlow, \emph{Circuits, Signals, and Systems
for Bioengineers}, 3.ª ed., 2018.

§4.2.2, \textbf{Amplitude Slicing---Quantization} (página 182 del visor
PDF).

\textbf{Correspondencia:} El tema se desarrolla en las secciones
indicadas. La redacción es una síntesis docente.

\textbf{Microcurrículo:} semana 7, sesión 2 (sesión global 14). Los
numerales del cronograma son identificadores curriculares; no son los
apartados del libro citados arriba.
\end{frame}

\begin{frame}{60. El rango dinámico exige margen}
\phantomsection\label{el-rango-dinuxe1mico-exige-margen}
\textbf{Libro guía:} John Semmlow, \emph{Circuits, Signals, and Systems
for Bioengineers}, 3.ª ed., 2018.

§4.2.2, \textbf{Amplitude Slicing---Quantization} (página 182 del visor
PDF).

\textbf{Correspondencia:} El tema se desarrolla en las secciones
indicadas. La redacción es una síntesis docente.

\textbf{Microcurrículo:} semana 7, sesión 2 (sesión global 14). Los
numerales del cronograma son identificadores curriculares; no son los
apartados del libro citados arriba.
\end{frame}

\begin{frame}{61. El reloj determina el eje temporal}
\phantomsection\label{el-reloj-determina-el-eje-temporal}
\textbf{Libro guía:} John Semmlow, \emph{Circuits, Signals, and Systems
for Bioengineers}, 3.ª ed., 2018.

§1.2.3, \textbf{Signal Representation: Continuous (Analog) Signals
Versus Discrete (Digital) Signals} (página 14 del visor PDF); §4.2,
\textbf{Data Acquisition And Storage} (página 175 del visor PDF).

\textbf{Correspondencia:} Tema previsto en el cronograma, desarrollado
como ampliación docente. Las referencias del libro son antecedentes
conceptuales. Metadatos, contrato de datos y sincronización: extensión
de la cadena de adquisición para el trabajo reproducible.

\textbf{Microcurrículo:} semana 7, sesión 2 (sesión global 14). Los
numerales del cronograma son identificadores curriculares; no son los
apartados del libro citados arriba.
\end{frame}

\begin{frame}{62. La sincronización sostiene análisis multicanal}
\phantomsection\label{la-sincronizaciuxf3n-sostiene-anuxe1lisis-multicanal}
\textbf{Libro guía:} John Semmlow, \emph{Circuits, Signals, and Systems
for Bioengineers}, 3.ª ed., 2018.

§1.2.3, \textbf{Signal Representation: Continuous (Analog) Signals
Versus Discrete (Digital) Signals} (página 14 del visor PDF); §4.2,
\textbf{Data Acquisition And Storage} (página 175 del visor PDF).

\textbf{Correspondencia:} Tema previsto en el cronograma, desarrollado
como ampliación docente. Las referencias del libro son antecedentes
conceptuales. Metadatos, contrato de datos y sincronización: extensión
de la cadena de adquisición para el trabajo reproducible.

\textbf{Microcurrículo:} semana 7, sesión 2 (sesión global 14). Los
numerales del cronograma son identificadores curriculares; no son los
apartados del libro citados arriba.
\end{frame}

\begin{frame}{63. Almacenar valores exige conservar contexto}
\phantomsection\label{almacenar-valores-exige-conservar-contexto}
\textbf{Libro guía:} John Semmlow, \emph{Circuits, Signals, and Systems
for Bioengineers}, 3.ª ed., 2018.

§1.2.3, \textbf{Signal Representation: Continuous (Analog) Signals
Versus Discrete (Digital) Signals} (página 14 del visor PDF); §4.2,
\textbf{Data Acquisition And Storage} (página 175 del visor PDF).

\textbf{Correspondencia:} Tema previsto en el cronograma, desarrollado
como ampliación docente. Las referencias del libro son antecedentes
conceptuales. Metadatos, contrato de datos y sincronización: extensión
de la cadena de adquisición para el trabajo reproducible.

\textbf{Microcurrículo:} semana 7, sesión 2 (sesión global 14). Los
numerales del cronograma son identificadores curriculares; no son los
apartados del libro citados arriba.
\end{frame}

\begin{frame}{64. Formato tabular no siempre basta}
\phantomsection\label{formato-tabular-no-siempre-basta}
\textbf{Libro guía:} John Semmlow, \emph{Circuits, Signals, and Systems
for Bioengineers}, 3.ª ed., 2018.

§1.2.3, \textbf{Signal Representation: Continuous (Analog) Signals
Versus Discrete (Digital) Signals} (página 14 del visor PDF); §4.2,
\textbf{Data Acquisition And Storage} (página 175 del visor PDF).

\textbf{Correspondencia:} Tema previsto en el cronograma, desarrollado
como ampliación docente. Las referencias del libro son antecedentes
conceptuales. Metadatos, contrato de datos y sincronización: extensión
de la cadena de adquisición para el trabajo reproducible.

\textbf{Microcurrículo:} semana 7, sesión 2 (sesión global 14). Los
numerales del cronograma son identificadores curriculares; no son los
apartados del libro citados arriba.
\end{frame}

\begin{frame}{65. Control de calidad antes del análisis}
\phantomsection\label{control-de-calidad-antes-del-anuxe1lisis}
\textbf{Libro guía:} John Semmlow, \emph{Circuits, Signals, and Systems
for Bioengineers}, 3.ª ed., 2018.

§1.2.3, \textbf{Signal Representation: Continuous (Analog) Signals
Versus Discrete (Digital) Signals} (página 14 del visor PDF); §4.2,
\textbf{Data Acquisition And Storage} (página 175 del visor PDF).

\textbf{Correspondencia:} Tema previsto en el cronograma, desarrollado
como ampliación docente. Las referencias del libro son antecedentes
conceptuales. Metadatos, contrato de datos y sincronización: extensión
de la cadena de adquisición para el trabajo reproducible.

\textbf{Microcurrículo:} semana 7, sesión 2 (sesión global 14). Los
numerales del cronograma son identificadores curriculares; no son los
apartados del libro citados arriba.
\end{frame}

\begin{frame}{66. Verificación de la sesión 14}
\phantomsection\label{verificaciuxf3n-de-la-sesiuxf3n-14}
\textbf{Libro guía:} John Semmlow, \emph{Circuits, Signals, and Systems
for Bioengineers}, 3.ª ed., 2018.

§1.2.2, \textbf{Biotransducers And Common Physiological Measurements}
(página 12 del visor PDF); §4.2.1, \textbf{Data Sampling: The Sampling
Theorem} (página 175 del visor PDF); §4.2.2, \textbf{Amplitude
Slicing---Quantization} (página 182 del visor PDF).

\textbf{Correspondencia:} Actividad o interpretación biomédica de
elaboración docente apoyada en las secciones indicadas.

\textbf{Microcurrículo:} semana 7, sesión 2 (sesión global 14). Los
numerales del cronograma son identificadores curriculares; no son los
apartados del libro citados arriba.
\end{frame}

\begin{frame}{67. Cierre de la semana 7}
\phantomsection\label{cierre-de-la-semana-7}
\textbf{Libro guía:} John Semmlow, \emph{Circuits, Signals, and Systems
for Bioengineers}, 3.ª ed., 2018.

§4.2, \textbf{Data Acquisition And Storage} (página 175 del visor PDF).

\textbf{Correspondencia:} Síntesis o encuadre que reúne varias
secciones; no corresponde a un único apartado del libro.

\textbf{Microcurrículo:} semana 7, sesión 2 (sesión global 14). Los
numerales del cronograma son identificadores curriculares; no son los
apartados del libro citados arriba.
\end{frame}

\begin{frame}{68. Semana 8 · Muestreo y estimación de potencia}
\phantomsection\label{semana-8-muestreo-y-estimaciuxf3n-de-potencia}
\textbf{Libro guía:} John Semmlow, \emph{Circuits, Signals, and Systems
for Bioengineers}, 3.ª ed., 2018.

§3.4, \textbf{Time--Frequency Transformation In The Discrete Domain}
(página 142 del visor PDF); §4.2, \textbf{Data Acquisition And Storage}
(página 175 del visor PDF); §4.6, \textbf{Time--Frequency Analysis}
(página 205 del visor PDF); §5.4, \textbf{Using The Impulse Response To
Find A Systems Output To Any Input---Convolution} (página 219 del visor
PDF).

\textbf{Correspondencia:} Síntesis o encuadre que reúne varias
secciones; no corresponde a un único apartado del libro.

\textbf{Lectura:} use estas secciones como ruta de preparación y
consulta. La bibliografía aparece al final del bloque.
\end{frame}

\begin{frame}{69. Sesión 15 · Muestrear observa la señal en instantes
discretos}
\phantomsection\label{sesiuxf3n-15-muestrear-observa-la-seuxf1al-en-instantes-discretos}
\textbf{Libro guía:} John Semmlow, \emph{Circuits, Signals, and Systems
for Bioengineers}, 3.ª ed., 2018.

§4.2.1, \textbf{Data Sampling: The Sampling Theorem} (página 175 del
visor PDF); §5.6.1, \textbf{Data Sampling Revisited} (página 238 del
visor PDF).

\textbf{Correspondencia:} El tema se desarrolla en las secciones
indicadas. La redacción es una síntesis docente.

\textbf{Microcurrículo:} semana 8, sesión 1 (sesión global 15). Los
numerales del cronograma son identificadores curriculares; no son los
apartados del libro citados arriba.
\end{frame}

\begin{frame}{70. Muestrear replica el espectro}
\phantomsection\label{muestrear-replica-el-espectro}
\textbf{Libro guía:} John Semmlow, \emph{Circuits, Signals, and Systems
for Bioengineers}, 3.ª ed., 2018.

§5.6.1, \textbf{Data Sampling Revisited} (página 238 del visor PDF).

\textbf{Correspondencia:} El tema se desarrolla en las secciones
indicadas. La redacción es una síntesis docente.

\textbf{Microcurrículo:} semana 8, sesión 1 (sesión global 15). Los
numerales del cronograma son identificadores curriculares; no son los
apartados del libro citados arriba.
\end{frame}

\begin{frame}{71. Nyquist establece una condición ideal}
\phantomsection\label{nyquist-establece-una-condiciuxf3n-ideal}
\textbf{Libro guía:} John Semmlow, \emph{Circuits, Signals, and Systems
for Bioengineers}, 3.ª ed., 2018.

§4.2.1, \textbf{Data Sampling: The Sampling Theorem} (página 175 del
visor PDF).

\textbf{Correspondencia:} El tema se desarrolla en las secciones
indicadas. La redacción es una síntesis docente.

\textbf{Microcurrículo:} semana 8, sesión 1 (sesión global 15). Los
numerales del cronograma son identificadores curriculares; no son los
apartados del libro citados arriba.
\end{frame}

\begin{frame}{72. La frecuencia de Nyquist no es la tasa de Nyquist}
\phantomsection\label{la-frecuencia-de-nyquist-no-es-la-tasa-de-nyquist}
\textbf{Libro guía:} John Semmlow, \emph{Circuits, Signals, and Systems
for Bioengineers}, 3.ª ed., 2018.

§4.2.1, \textbf{Data Sampling: The Sampling Theorem} (página 175 del
visor PDF).

\textbf{Correspondencia:} El tema se desarrolla en las secciones
indicadas. La redacción es una síntesis docente.

\textbf{Microcurrículo:} semana 8, sesión 1 (sesión global 15). Los
numerales del cronograma son identificadores curriculares; no son los
apartados del libro citados arriba.
\end{frame}

\begin{frame}{73. El alias de una frecuencia depende de \(f_s\)}
\phantomsection\label{el-alias-de-una-frecuencia-depende-de-f_s}
\textbf{Libro guía:} John Semmlow, \emph{Circuits, Signals, and Systems
for Bioengineers}, 3.ª ed., 2018.

§4.2.1, \textbf{Data Sampling: The Sampling Theorem} (página 175 del
visor PDF).

\textbf{Correspondencia:} El tema se desarrolla en las secciones
indicadas. La redacción es una síntesis docente.

\textbf{Microcurrículo:} semana 8, sesión 1 (sesión global 15). Los
numerales del cronograma son identificadores curriculares; no son los
apartados del libro citados arriba.
\end{frame}

\begin{frame}{74. El filtro antialias necesita una banda de transición}
\phantomsection\label{el-filtro-antialias-necesita-una-banda-de-transiciuxf3n}
\textbf{Libro guía:} John Semmlow, \emph{Circuits, Signals, and Systems
for Bioengineers}, 3.ª ed., 2018.

§4.2.1, \textbf{Data Sampling: The Sampling Theorem} (página 175 del
visor PDF).

\textbf{Correspondencia:} El tema se desarrolla en las secciones
indicadas. La redacción es una síntesis docente.

\textbf{Microcurrículo:} semana 8, sesión 1 (sesión global 15). Los
numerales del cronograma son identificadores curriculares; no son los
apartados del libro citados arriba.
\end{frame}

\begin{frame}{75. Sobremuestrear simplifica parte del diseño}
\phantomsection\label{sobremuestrear-simplifica-parte-del-diseuxf1o}
\textbf{Libro guía:} John Semmlow, \emph{Circuits, Signals, and Systems
for Bioengineers}, 3.ª ed., 2018.

§4.2.1, \textbf{Data Sampling: The Sampling Theorem} (página 175 del
visor PDF).

\textbf{Correspondencia:} El tema se desarrolla en las secciones
indicadas. La redacción es una síntesis docente.

\textbf{Microcurrículo:} semana 8, sesión 1 (sesión global 15). Los
numerales del cronograma son identificadores curriculares; no son los
apartados del libro citados arriba.
\end{frame}

\begin{frame}{76. La resolución temporal básica es \(T_s\)}
\phantomsection\label{la-resoluciuxf3n-temporal-buxe1sica-es-t_s}
\textbf{Libro guía:} John Semmlow, \emph{Circuits, Signals, and Systems
for Bioengineers}, 3.ª ed., 2018.

§4.2.1, \textbf{Data Sampling: The Sampling Theorem} (página 175 del
visor PDF).

\textbf{Correspondencia:} El tema se desarrolla en las secciones
indicadas. La redacción es una síntesis docente.

\textbf{Microcurrículo:} semana 8, sesión 1 (sesión global 15). Los
numerales del cronograma son identificadores curriculares; no son los
apartados del libro citados arriba.
\end{frame}

\begin{frame}{77. Duración y muestreo controlan aspectos distintos}
\phantomsection\label{duraciuxf3n-y-muestreo-controlan-aspectos-distintos}
\textbf{Libro guía:} John Semmlow, \emph{Circuits, Signals, and Systems
for Bioengineers}, 3.ª ed., 2018.

§4.2.1, \textbf{Data Sampling: The Sampling Theorem} (página 175 del
visor PDF).

\textbf{Correspondencia:} El tema se desarrolla en las secciones
indicadas. La redacción es una síntesis docente.

\textbf{Microcurrículo:} semana 8, sesión 1 (sesión global 15). Los
numerales del cronograma son identificadores curriculares; no son los
apartados del libro citados arriba.
\end{frame}

\begin{frame}{78. Muestreo irregular exige otro modelo}
\phantomsection\label{muestreo-irregular-exige-otro-modelo}
\textbf{Libro guía:} John Semmlow, \emph{Circuits, Signals, and Systems
for Bioengineers}, 3.ª ed., 2018.

§4.2.1, \textbf{Data Sampling: The Sampling Theorem} (página 175 del
visor PDF).

\textbf{Correspondencia:} Las secciones aportan el fundamento. La
diapositiva amplía la formulación o precisa condiciones que no se
presentan con idéntica notación en el libro. El muestreo uniforme es el
modelo del teorema; datos irregulares requieren otro modelo o un
remuestreo justificado.

\textbf{Microcurrículo:} semana 8, sesión 1 (sesión global 15). Los
numerales del cronograma son identificadores curriculares; no son los
apartados del libro citados arriba.
\end{frame}

\begin{frame}{79. Código: observar aliasing}
\phantomsection\label{cuxf3digo-observar-aliasing}
\textbf{Libro guía:} John Semmlow, \emph{Circuits, Signals, and Systems
for Bioengineers}, 3.ª ed., 2018.

§4.2.1, \textbf{Data Sampling: The Sampling Theorem} (página 175 del
visor PDF).

\textbf{Correspondencia:} Actividad o interpretación biomédica de
elaboración docente apoyada en las secciones indicadas.

\textbf{Microcurrículo:} semana 8, sesión 1 (sesión global 15). Los
numerales del cronograma son identificadores curriculares; no son los
apartados del libro citados arriba.
\end{frame}

\begin{frame}{80. Selección de \(f_s\) para una bioseñal}
\phantomsection\label{selecciuxf3n-de-f_s-para-una-bioseuxf1al}
\textbf{Libro guía:} John Semmlow, \emph{Circuits, Signals, and Systems
for Bioengineers}, 3.ª ed., 2018.

§4.2.1, \textbf{Data Sampling: The Sampling Theorem} (página 175 del
visor PDF).

\textbf{Correspondencia:} Actividad o interpretación biomédica de
elaboración docente apoyada en las secciones indicadas.

\textbf{Microcurrículo:} semana 8, sesión 1 (sesión global 15). Los
numerales del cronograma son identificadores curriculares; no son los
apartados del libro citados arriba.
\end{frame}

\begin{frame}{81. Verificación de la sesión 15}
\phantomsection\label{verificaciuxf3n-de-la-sesiuxf3n-15}
\textbf{Libro guía:} John Semmlow, \emph{Circuits, Signals, and Systems
for Bioengineers}, 3.ª ed., 2018.

§4.2.1, \textbf{Data Sampling: The Sampling Theorem} (página 175 del
visor PDF).

\textbf{Correspondencia:} Actividad o interpretación biomédica de
elaboración docente apoyada en las secciones indicadas.

\textbf{Microcurrículo:} semana 8, sesión 1 (sesión global 15). Los
numerales del cronograma son identificadores curriculares; no son los
apartados del libro citados arriba.
\end{frame}

\begin{frame}{82. Sesión 16 · Potencia espectral describe distribución
cuadrática}
\phantomsection\label{sesiuxf3n-16-potencia-espectral-describe-distribuciuxf3n-cuadruxe1tica}
\textbf{Libro guía:} John Semmlow, \emph{Circuits, Signals, and Systems
for Bioengineers}, 3.ª ed., 2018.

§4.3, \textbf{Power Spectrum} (página 191 del visor PDF).

\textbf{Correspondencia:} El tema se desarrolla en las secciones
indicadas. La redacción es una síntesis docente.

\textbf{Microcurrículo:} semana 8, sesión 2 (sesión global 16). Los
numerales del cronograma son identificadores curriculares; no son los
apartados del libro citados arriba.
\end{frame}

\begin{frame}{83. Energía y potencia no son intercambiables}
\phantomsection\label{energuxeda-y-potencia-no-son-intercambiables}
\textbf{Libro guía:} John Semmlow, \emph{Circuits, Signals, and Systems
for Bioengineers}, 3.ª ed., 2018.

§2.2.1, \textbf{Mean And Root Mean Squared Signal Measurements} (página
58 del visor PDF); §4.3, \textbf{Power Spectrum} (página 191 del visor
PDF).

\textbf{Correspondencia:} Las secciones aportan el fundamento. La
diapositiva amplía la formulación o precisa condiciones que no se
presentan con idéntica notación en el libro.

\textbf{Microcurrículo:} semana 8, sesión 2 (sesión global 16). Los
numerales del cronograma son identificadores curriculares; no son los
apartados del libro citados arriba.
\end{frame}

\begin{frame}{84. Espectro de energía y PSD responden a modelos
distintos}
\phantomsection\label{espectro-de-energuxeda-y-psd-responden-a-modelos-distintos}
\textbf{Libro guía:} John Semmlow, \emph{Circuits, Signals, and Systems
for Bioengineers}, 3.ª ed., 2018.

§2.2.1, \textbf{Mean And Root Mean Squared Signal Measurements} (página
58 del visor PDF); §4.3, \textbf{Power Spectrum} (página 191 del visor
PDF).

\textbf{Correspondencia:} Las secciones aportan el fundamento. La
diapositiva amplía la formulación o precisa condiciones que no se
presentan con idéntica notación en el libro.

\textbf{Microcurrículo:} semana 8, sesión 2 (sesión global 16). Los
numerales del cronograma son identificadores curriculares; no son los
apartados del libro citados arriba.
\end{frame}

\begin{frame}{85. El periodograma parte de la DFT}
\phantomsection\label{el-periodograma-parte-de-la-dft}
\textbf{Libro guía:} John Semmlow, \emph{Circuits, Signals, and Systems
for Bioengineers}, 3.ª ed., 2018.

§3.4.1, \textbf{The Discrete Fourier Transform} (página 142 del visor
PDF); §4.3, \textbf{Power Spectrum} (página 191 del visor PDF).

\textbf{Correspondencia:} El tema se desarrolla en las secciones
indicadas. La redacción es una síntesis docente.

\textbf{Microcurrículo:} semana 8, sesión 2 (sesión global 16). Los
numerales del cronograma son identificadores curriculares; no son los
apartados del libro citados arriba.
\end{frame}

\begin{frame}{86. Las unidades distinguen espectro y densidad}
\phantomsection\label{las-unidades-distinguen-espectro-y-densidad}
\textbf{Libro guía:} John Semmlow, \emph{Circuits, Signals, and Systems
for Bioengineers}, 3.ª ed., 2018.

§4.3, \textbf{Power Spectrum} (página 191 del visor PDF).

\textbf{Correspondencia:} Las secciones aportan el fundamento. La
diapositiva amplía la formulación o precisa condiciones que no se
presentan con idéntica notación en el libro.

\textbf{Microcurrículo:} semana 8, sesión 2 (sesión global 16). Los
numerales del cronograma son identificadores curriculares; no son los
apartados del libro citados arriba.
\end{frame}

\begin{frame}{87. La PSD unilateral conserva potencia}
\phantomsection\label{la-psd-unilateral-conserva-potencia}
\textbf{Libro guía:} John Semmlow, \emph{Circuits, Signals, and Systems
for Bioengineers}, 3.ª ed., 2018.

§4.3, \textbf{Power Spectrum} (página 191 del visor PDF).

\textbf{Correspondencia:} Las secciones aportan el fundamento. La
diapositiva amplía la formulación o precisa condiciones que no se
presentan con idéntica notación en el libro.

\textbf{Microcurrículo:} semana 8, sesión 2 (sesión global 16). Los
numerales del cronograma son identificadores curriculares; no son los
apartados del libro citados arriba.
\end{frame}

\begin{frame}{88. Parseval ofrece una prueba de consistencia}
\phantomsection\label{parseval-ofrece-una-prueba-de-consistencia}
\textbf{Libro guía:} John Semmlow, \emph{Circuits, Signals, and Systems
for Bioengineers}, 3.ª ed., 2018.

§4.3, \textbf{Power Spectrum} (página 191 del visor PDF).

\textbf{Correspondencia:} Las secciones aportan el fundamento. La
diapositiva amplía la formulación o precisa condiciones que no se
presentan con idéntica notación en el libro.

\textbf{Microcurrículo:} semana 8, sesión 2 (sesión global 16). Los
numerales del cronograma son identificadores curriculares; no son los
apartados del libro citados arriba.
\end{frame}

\begin{frame}{89. Parseval con ventana: comparación correcta}
\phantomsection\label{parseval-con-ventana-comparaciuxf3n-correcta}
\textbf{Libro guía:} John Semmlow, \emph{Circuits, Signals, and Systems
for Bioengineers}, 3.ª ed., 2018.

§4.3, \textbf{Power Spectrum} (página 191 del visor PDF); §4.4,
\textbf{Spectral Averaging} (página 195 del visor PDF).

\textbf{Correspondencia:} Ejemplo, figura o implementación de
elaboración docente a partir de estos fundamentos. No es una
transcripción del código MATLAB del libro. Escalamiento contrastado con
la documentación oficial de scipy.signal.periodogram.

\textbf{Microcurrículo:} semana 8, sesión 2 (sesión global 16). Los
numerales del cronograma son identificadores curriculares; no son los
apartados del libro citados arriba.

\textbf{Implementación:}
\href{https://docs.scipy.org/doc/scipy/reference/generated/scipy.signal.periodogram.html}{SciPy:
periodogram y escalamiento espectral}.
\end{frame}

\begin{frame}{90. Retirar la media cambia la PSD cerca de 0 Hz}
\phantomsection\label{retirar-la-media-cambia-la-psd-cerca-de-0-hz}
\textbf{Libro guía:} John Semmlow, \emph{Circuits, Signals, and Systems
for Bioengineers}, 3.ª ed., 2018.

§2.2.1, \textbf{Mean And Root Mean Squared Signal Measurements} (página
58 del visor PDF); §4.3, \textbf{Power Spectrum} (página 191 del visor
PDF).

\textbf{Correspondencia:} El tema se desarrolla en las secciones
indicadas. La redacción es una síntesis docente.

\textbf{Microcurrículo:} semana 8, sesión 2 (sesión global 16). Los
numerales del cronograma son identificadores curriculares; no son los
apartados del libro citados arriba.
\end{frame}

\begin{frame}{91. Autocorrelación y PSD: el vínculo estadístico}
\phantomsection\label{autocorrelaciuxf3n-y-psd-el-vuxednculo-estaduxedstico}
\textbf{Libro guía:} John Semmlow, \emph{Circuits, Signals, and Systems
for Bioengineers}, 3.ª ed., 2018.

§2.4.5, \textbf{Autocorrelation} (página 104 del visor PDF); §2.4.6,
\textbf{Autocovariance And Cross-Covariance} (página 108 del visor PDF);
§4.3, \textbf{Power Spectrum} (página 191 del visor PDF).

\textbf{Correspondencia:} Las secciones aportan el fundamento. La
diapositiva amplía la formulación o precisa condiciones que no se
presentan con idéntica notación en el libro. Formalización estadística
del vínculo entre correlación y potencia; no confundir una estimación
finita con el parámetro del proceso.

\textbf{Microcurrículo:} semana 8, sesión 2 (sesión global 16). Los
numerales del cronograma son identificadores curriculares; no son los
apartados del libro citados arriba.
\end{frame}

\begin{frame}{92. El periodograma tiene alta variabilidad}
\phantomsection\label{el-periodograma-tiene-alta-variabilidad}
\textbf{Libro guía:} John Semmlow, \emph{Circuits, Signals, and Systems
for Bioengineers}, 3.ª ed., 2018.

§4.4, \textbf{Spectral Averaging} (página 195 del visor PDF).

\textbf{Correspondencia:} El tema se desarrolla en las secciones
indicadas. La redacción es una síntesis docente.

\textbf{Microcurrículo:} semana 8, sesión 2 (sesión global 16). Los
numerales del cronograma son identificadores curriculares; no son los
apartados del libro citados arriba.
\end{frame}

\begin{frame}{93. Potencia en una banda}
\phantomsection\label{potencia-en-una-banda}
\textbf{Libro guía:} John Semmlow, \emph{Circuits, Signals, and Systems
for Bioengineers}, 3.ª ed., 2018.

§4.3, \textbf{Power Spectrum} (página 191 del visor PDF); §4.5,
\textbf{Signal Bandwidth} (página 200 del visor PDF).

\textbf{Correspondencia:} El tema se desarrolla en las secciones
indicadas. La redacción es una síntesis docente.

\textbf{Microcurrículo:} semana 8, sesión 2 (sesión global 16). Los
numerales del cronograma son identificadores curriculares; no son los
apartados del libro citados arriba.
\end{frame}

\begin{frame}{94. Potencia absoluta y relativa}
\phantomsection\label{potencia-absoluta-y-relativa}
\textbf{Libro guía:} John Semmlow, \emph{Circuits, Signals, and Systems
for Bioengineers}, 3.ª ed., 2018.

§4.3, \textbf{Power Spectrum} (página 191 del visor PDF).

\textbf{Correspondencia:} Las secciones aportan el fundamento. La
diapositiva amplía la formulación o precisa condiciones que no se
presentan con idéntica notación en el libro.

\textbf{Microcurrículo:} semana 8, sesión 2 (sesión global 16). Los
numerales del cronograma son identificadores curriculares; no son los
apartados del libro citados arriba.
\end{frame}

\begin{frame}{95. Código: periodograma con unidades explícitas}
\phantomsection\label{cuxf3digo-periodograma-con-unidades-expluxedcitas}
\textbf{Libro guía:} John Semmlow, \emph{Circuits, Signals, and Systems
for Bioengineers}, 3.ª ed., 2018.

§4.3, \textbf{Power Spectrum} (página 191 del visor PDF).

\textbf{Correspondencia:} Actividad o interpretación biomédica de
elaboración docente apoyada en las secciones indicadas.

\textbf{Microcurrículo:} semana 8, sesión 2 (sesión global 16). Los
numerales del cronograma son identificadores curriculares; no son los
apartados del libro citados arriba.

\textbf{Implementación:}
\href{https://docs.scipy.org/doc/scipy/reference/generated/scipy.signal.periodogram.html}{SciPy:
periodogram y escalamiento espectral}.
\end{frame}

\begin{frame}{96. Lectura biomédica prudente}
\phantomsection\label{lectura-biomuxe9dica-prudente}
\textbf{Libro guía:} John Semmlow, \emph{Circuits, Signals, and Systems
for Bioengineers}, 3.ª ed., 2018.

§4.3, \textbf{Power Spectrum} (página 191 del visor PDF).

\textbf{Correspondencia:} Actividad o interpretación biomédica de
elaboración docente apoyada en las secciones indicadas.

\textbf{Microcurrículo:} semana 8, sesión 2 (sesión global 16). Los
numerales del cronograma son identificadores curriculares; no son los
apartados del libro citados arriba.
\end{frame}

\begin{frame}{97. Verificación de la sesión 16}
\phantomsection\label{verificaciuxf3n-de-la-sesiuxf3n-16}
\textbf{Libro guía:} John Semmlow, \emph{Circuits, Signals, and Systems
for Bioengineers}, 3.ª ed., 2018.

§4.3, \textbf{Power Spectrum} (página 191 del visor PDF).

\textbf{Correspondencia:} Actividad o interpretación biomédica de
elaboración docente apoyada en las secciones indicadas.

\textbf{Microcurrículo:} semana 8, sesión 2 (sesión global 16). Los
numerales del cronograma son identificadores curriculares; no son los
apartados del libro citados arriba.
\end{frame}

\begin{frame}{98. Cierre de la semana 8}
\phantomsection\label{cierre-de-la-semana-8}
\textbf{Libro guía:} John Semmlow, \emph{Circuits, Signals, and Systems
for Bioengineers}, 3.ª ed., 2018.

§4.2.1, \textbf{Data Sampling: The Sampling Theorem} (página 175 del
visor PDF); §4.3, \textbf{Power Spectrum} (página 191 del visor PDF).

\textbf{Correspondencia:} Síntesis o encuadre que reúne varias
secciones; no corresponde a un único apartado del libro.

\textbf{Microcurrículo:} semana 8, sesión 2 (sesión global 16). Los
numerales del cronograma son identificadores curriculares; no son los
apartados del libro citados arriba.
\end{frame}

\begin{frame}{99. Semana 9 · Welch y análisis tiempo-frecuencia}
\phantomsection\label{semana-9-welch-y-anuxe1lisis-tiempo-frecuencia}
\textbf{Libro guía:} John Semmlow, \emph{Circuits, Signals, and Systems
for Bioengineers}, 3.ª ed., 2018.

§3.4, \textbf{Time--Frequency Transformation In The Discrete Domain}
(página 142 del visor PDF); §4.2, \textbf{Data Acquisition And Storage}
(página 175 del visor PDF); §4.6, \textbf{Time--Frequency Analysis}
(página 205 del visor PDF); §5.4, \textbf{Using The Impulse Response To
Find A Systems Output To Any Input---Convolution} (página 219 del visor
PDF).

\textbf{Correspondencia:} Síntesis o encuadre que reúne varias
secciones; no corresponde a un único apartado del libro.

\textbf{Lectura:} use estas secciones como ruta de preparación y
consulta. La bibliografía aparece al final del bloque.
\end{frame}

\begin{frame}{100. Sesión 17 · Welch reduce variabilidad mediante
promediado}
\phantomsection\label{sesiuxf3n-17-welch-reduce-variabilidad-mediante-promediado}
\textbf{Libro guía:} John Semmlow, \emph{Circuits, Signals, and Systems
for Bioengineers}, 3.ª ed., 2018.

§4.4, \textbf{Spectral Averaging} (página 195 del visor PDF).

\textbf{Correspondencia:} El tema se desarrolla en las secciones
indicadas. La redacción es una síntesis docente.

\textbf{Microcurrículo:} semana 9, sesión 1 (sesión global 17). Los
numerales del cronograma son identificadores curriculares; no son los
apartados del libro citados arriba.
\end{frame}

\begin{frame}{101. El procedimiento contiene cuatro decisiones}
\phantomsection\label{el-procedimiento-contiene-cuatro-decisiones}
\textbf{Libro guía:} John Semmlow, \emph{Circuits, Signals, and Systems
for Bioengineers}, 3.ª ed., 2018.

§4.4, \textbf{Spectral Averaging} (página 195 del visor PDF).

\textbf{Correspondencia:} El tema se desarrolla en las secciones
indicadas. La redacción es una síntesis docente.

\textbf{Microcurrículo:} semana 9, sesión 1 (sesión global 17). Los
numerales del cronograma son identificadores curriculares; no son los
apartados del libro citados arriba.
\end{frame}

\begin{frame}{102. Segmentos cortos producen más promedios}
\phantomsection\label{segmentos-cortos-producen-muxe1s-promedios}
\textbf{Libro guía:} John Semmlow, \emph{Circuits, Signals, and Systems
for Bioengineers}, 3.ª ed., 2018.

§4.4, \textbf{Spectral Averaging} (página 195 del visor PDF).

\textbf{Correspondencia:} El tema se desarrolla en las secciones
indicadas. La redacción es una síntesis docente.

\textbf{Microcurrículo:} semana 9, sesión 1 (sesión global 17). Los
numerales del cronograma son identificadores curriculares; no son los
apartados del libro citados arriba.
\end{frame}

\begin{frame}{103. El solapamiento reutiliza muestras}
\phantomsection\label{el-solapamiento-reutiliza-muestras}
\textbf{Libro guía:} John Semmlow, \emph{Circuits, Signals, and Systems
for Bioengineers}, 3.ª ed., 2018.

§4.4, \textbf{Spectral Averaging} (página 195 del visor PDF).

\textbf{Correspondencia:} El tema se desarrolla en las secciones
indicadas. La redacción es una síntesis docente.

\textbf{Microcurrículo:} semana 9, sesión 1 (sesión global 17). Los
numerales del cronograma son identificadores curriculares; no son los
apartados del libro citados arriba.
\end{frame}

\begin{frame}{104. Hann y 50 \% forman una elección frecuente}
\phantomsection\label{hann-y-50-forman-una-elecciuxf3n-frecuente}
\textbf{Libro guía:} John Semmlow, \emph{Circuits, Signals, and Systems
for Bioengineers}, 3.ª ed., 2018.

§4.4, \textbf{Spectral Averaging} (página 195 del visor PDF).

\textbf{Correspondencia:} El tema se desarrolla en las secciones
indicadas. La redacción es una síntesis docente.

\textbf{Microcurrículo:} semana 9, sesión 1 (sesión global 17). Los
numerales del cronograma son identificadores curriculares; no son los
apartados del libro citados arriba.
\end{frame}

\begin{frame}{105. El escalamiento incluye la energía de la ventana}
\phantomsection\label{el-escalamiento-incluye-la-energuxeda-de-la-ventana}
\textbf{Libro guía:} John Semmlow, \emph{Circuits, Signals, and Systems
for Bioengineers}, 3.ª ed., 2018.

§4.4, \textbf{Spectral Averaging} (página 195 del visor PDF).

\textbf{Correspondencia:} El tema se desarrolla en las secciones
indicadas. La redacción es una síntesis docente.

\textbf{Microcurrículo:} semana 9, sesión 1 (sesión global 17). Los
numerales del cronograma son identificadores curriculares; no son los
apartados del libro citados arriba.
\end{frame}

\begin{frame}{106. \(N_{FFT}\) cambia la malla, no la evidencia}
\phantomsection\label{n_fft-cambia-la-malla-no-la-evidencia}
\textbf{Libro guía:} John Semmlow, \emph{Circuits, Signals, and Systems
for Bioengineers}, 3.ª ed., 2018.

§4.4, \textbf{Spectral Averaging} (página 195 del visor PDF).

\textbf{Correspondencia:} El tema se desarrolla en las secciones
indicadas. La redacción es una síntesis docente.

\textbf{Microcurrículo:} semana 9, sesión 1 (sesión global 17). Los
numerales del cronograma son identificadores curriculares; no son los
apartados del libro citados arriba.
\end{frame}

\begin{frame}{107. Promediar requiere segmentos comparables}
\phantomsection\label{promediar-requiere-segmentos-comparables}
\textbf{Libro guía:} John Semmlow, \emph{Circuits, Signals, and Systems
for Bioengineers}, 3.ª ed., 2018.

§4.4, \textbf{Spectral Averaging} (página 195 del visor PDF).

\textbf{Correspondencia:} El tema se desarrolla en las secciones
indicadas. La redacción es una síntesis docente.

\textbf{Microcurrículo:} semana 9, sesión 1 (sesión global 17). Los
numerales del cronograma son identificadores curriculares; no son los
apartados del libro citados arriba.
\end{frame}

\begin{frame}{108. Ancho de banda puede definirse de varias formas}
\phantomsection\label{ancho-de-banda-puede-definirse-de-varias-formas}
\textbf{Libro guía:} John Semmlow, \emph{Circuits, Signals, and Systems
for Bioengineers}, 3.ª ed., 2018.

§4.5, \textbf{Signal Bandwidth} (página 200 del visor PDF).

\textbf{Correspondencia:} El tema se desarrolla en las secciones
indicadas. La redacción es una síntesis docente.

\textbf{Microcurrículo:} semana 9, sesión 1 (sesión global 17). Los
numerales del cronograma son identificadores curriculares; no son los
apartados del libro citados arriba.
\end{frame}

\begin{frame}{109. Frecuencia mediana espectral}
\phantomsection\label{frecuencia-mediana-espectral}
\textbf{Libro guía:} John Semmlow, \emph{Circuits, Signals, and Systems
for Bioengineers}, 3.ª ed., 2018.

§4.3, \textbf{Power Spectrum} (página 191 del visor PDF); §4.5,
\textbf{Signal Bandwidth} (página 200 del visor PDF).

\textbf{Correspondencia:} Las secciones aportan el fundamento. La
diapositiva amplía la formulación o precisa condiciones que no se
presentan con idéntica notación en el libro. Frecuencia media/mediana
son descriptores de localización y no sinónimos de ancho de banda.

\textbf{Microcurrículo:} semana 9, sesión 1 (sesión global 17). Los
numerales del cronograma son identificadores curriculares; no son los
apartados del libro citados arriba.
\end{frame}

\begin{frame}{110. Frecuencia media espectral}
\phantomsection\label{frecuencia-media-espectral}
\textbf{Libro guía:} John Semmlow, \emph{Circuits, Signals, and Systems
for Bioengineers}, 3.ª ed., 2018.

§4.3, \textbf{Power Spectrum} (página 191 del visor PDF); §4.5,
\textbf{Signal Bandwidth} (página 200 del visor PDF).

\textbf{Correspondencia:} Las secciones aportan el fundamento. La
diapositiva amplía la formulación o precisa condiciones que no se
presentan con idéntica notación en el libro. Frecuencia media/mediana
son descriptores de localización y no sinónimos de ancho de banda.

\textbf{Microcurrículo:} semana 9, sesión 1 (sesión global 17). Los
numerales del cronograma son identificadores curriculares; no son los
apartados del libro citados arriba.
\end{frame}

\begin{frame}{111. Potencia de banda mediante integración numérica}
\phantomsection\label{potencia-de-banda-mediante-integraciuxf3n-numuxe9rica}
\textbf{Libro guía:} John Semmlow, \emph{Circuits, Signals, and Systems
for Bioengineers}, 3.ª ed., 2018.

§4.3, \textbf{Power Spectrum} (página 191 del visor PDF); §4.4,
\textbf{Spectral Averaging} (página 195 del visor PDF).

\textbf{Correspondencia:} Ejemplo, figura o implementación de
elaboración docente a partir de estos fundamentos. No es una
transcripción del código MATLAB del libro.

\textbf{Microcurrículo:} semana 9, sesión 1 (sesión global 17). Los
numerales del cronograma son identificadores curriculares; no son los
apartados del libro citados arriba.

\textbf{Implementación:}
\href{https://docs.scipy.org/doc/scipy/reference/generated/scipy.signal.periodogram.html}{SciPy:
periodogram y escalamiento espectral}.
\end{frame}

\begin{frame}{112. Selección razonada de parámetros}
\phantomsection\label{selecciuxf3n-razonada-de-paruxe1metros}
\textbf{Libro guía:} John Semmlow, \emph{Circuits, Signals, and Systems
for Bioengineers}, 3.ª ed., 2018.

§4.4, \textbf{Spectral Averaging} (página 195 del visor PDF); §4.5,
\textbf{Signal Bandwidth} (página 200 del visor PDF).

\textbf{Correspondencia:} Actividad o interpretación biomédica de
elaboración docente apoyada en las secciones indicadas.

\textbf{Microcurrículo:} semana 9, sesión 1 (sesión global 17). Los
numerales del cronograma son identificadores curriculares; no son los
apartados del libro citados arriba.
\end{frame}

\begin{frame}{113. Comparación entre condiciones}
\phantomsection\label{comparaciuxf3n-entre-condiciones}
\textbf{Libro guía:} John Semmlow, \emph{Circuits, Signals, and Systems
for Bioengineers}, 3.ª ed., 2018.

§4.4, \textbf{Spectral Averaging} (página 195 del visor PDF); §4.5,
\textbf{Signal Bandwidth} (página 200 del visor PDF).

\textbf{Correspondencia:} Actividad o interpretación biomédica de
elaboración docente apoyada en las secciones indicadas.

\textbf{Microcurrículo:} semana 9, sesión 1 (sesión global 17). Los
numerales del cronograma son identificadores curriculares; no son los
apartados del libro citados arriba.
\end{frame}

\begin{frame}{114. Verificación de la sesión 17}
\phantomsection\label{verificaciuxf3n-de-la-sesiuxf3n-17}
\textbf{Libro guía:} John Semmlow, \emph{Circuits, Signals, and Systems
for Bioengineers}, 3.ª ed., 2018.

§4.4, \textbf{Spectral Averaging} (página 195 del visor PDF); §4.5,
\textbf{Signal Bandwidth} (página 200 del visor PDF).

\textbf{Correspondencia:} Actividad o interpretación biomédica de
elaboración docente apoyada en las secciones indicadas.

\textbf{Microcurrículo:} semana 9, sesión 1 (sesión global 17). Los
numerales del cronograma son identificadores curriculares; no son los
apartados del libro citados arriba.
\end{frame}

\begin{frame}{115. Sesión 18 · Señales no estacionarias}
\phantomsection\label{sesiuxf3n-18-seuxf1ales-no-estacionarias}
\textbf{Libro guía:} John Semmlow, \emph{Circuits, Signals, and Systems
for Bioengineers}, 3.ª ed., 2018.

§4.6, \textbf{Time--Frequency Analysis} (página 205 del visor PDF).

\textbf{Correspondencia:} El tema se desarrolla en las secciones
indicadas. La redacción es una síntesis docente.

\textbf{Microcurrículo:} semana 9, sesión 2 (sesión global 18). Los
numerales del cronograma son identificadores curriculares; no son los
apartados del libro citados arriba.
\end{frame}

\begin{frame}{116. La estacionariedad es un supuesto de modelamiento}
\phantomsection\label{la-estacionariedad-es-un-supuesto-de-modelamiento}
\textbf{Libro guía:} John Semmlow, \emph{Circuits, Signals, and Systems
for Bioengineers}, 3.ª ed., 2018.

§10.2, \textbf{Stochastic Processes: Stationarity And Ergodicity}
(página 450 del visor PDF).

\textbf{Correspondencia:} El tema se desarrolla en las secciones
indicadas. La redacción es una síntesis docente.

\textbf{Microcurrículo:} semana 9, sesión 2 (sesión global 18). Los
numerales del cronograma son identificadores curriculares; no son los
apartados del libro citados arriba.
\end{frame}

\begin{frame}{117. La STFT aplica Fourier localmente}
\phantomsection\label{la-stft-aplica-fourier-localmente}
\textbf{Libro guía:} John Semmlow, \emph{Circuits, Signals, and Systems
for Bioengineers}, 3.ª ed., 2018.

§4.6, \textbf{Time--Frequency Analysis} (página 205 del visor PDF).

\textbf{Correspondencia:} El tema se desarrolla en las secciones
indicadas. La redacción es una síntesis docente.

\textbf{Microcurrículo:} semana 9, sesión 2 (sesión global 18). Los
numerales del cronograma son identificadores curriculares; no son los
apartados del libro citados arriba.
\end{frame}

\begin{frame}{118. El espectrograma muestra magnitud o potencia}
\phantomsection\label{el-espectrograma-muestra-magnitud-o-potencia}
\textbf{Libro guía:} John Semmlow, \emph{Circuits, Signals, and Systems
for Bioengineers}, 3.ª ed., 2018.

§4.6, \textbf{Time--Frequency Analysis} (página 205 del visor PDF).

\textbf{Correspondencia:} El tema se desarrolla en las secciones
indicadas. La redacción es una síntesis docente.

\textbf{Microcurrículo:} semana 9, sesión 2 (sesión global 18). Los
numerales del cronograma son identificadores curriculares; no son los
apartados del libro citados arriba.
\end{frame}

\begin{frame}{119. La ventana fija impone un compromiso}
\phantomsection\label{la-ventana-fija-impone-un-compromiso}
\textbf{Libro guía:} John Semmlow, \emph{Circuits, Signals, and Systems
for Bioengineers}, 3.ª ed., 2018.

§4.6, \textbf{Time--Frequency Analysis} (página 205 del visor PDF).

\textbf{Correspondencia:} El tema se desarrolla en las secciones
indicadas. La redacción es una síntesis docente.

\textbf{Microcurrículo:} semana 9, sesión 2 (sesión global 18). Los
numerales del cronograma son identificadores curriculares; no son los
apartados del libro citados arriba.
\end{frame}

\begin{frame}{120. La incertidumbre expresa ese límite}
\phantomsection\label{la-incertidumbre-expresa-ese-luxedmite}
\textbf{Libro guía:} John Semmlow, \emph{Circuits, Signals, and Systems
for Bioengineers}, 3.ª ed., 2018.

§4.6, \textbf{Time--Frequency Analysis} (página 205 del visor PDF).

\textbf{Correspondencia:} Las secciones aportan el fundamento. La
diapositiva amplía la formulación o precisa condiciones que no se
presentan con idéntica notación en el libro. Límite de dispersión
tiempo--frecuencia formalizado para la convención en Hz y desviaciones
de energía.

\textbf{Microcurrículo:} semana 9, sesión 2 (sesión global 18). Los
numerales del cronograma son identificadores curriculares; no son los
apartados del libro citados arriba.
\end{frame}

\begin{frame}{121. El avance controla densidad temporal}
\phantomsection\label{el-avance-controla-densidad-temporal}
\textbf{Libro guía:} John Semmlow, \emph{Circuits, Signals, and Systems
for Bioengineers}, 3.ª ed., 2018.

§4.6, \textbf{Time--Frequency Analysis} (página 205 del visor PDF).

\textbf{Correspondencia:} Actividad o interpretación biomédica de
elaboración docente apoyada en las secciones indicadas.

\textbf{Microcurrículo:} semana 9, sesión 2 (sesión global 18). Los
numerales del cronograma son identificadores curriculares; no son los
apartados del libro citados arriba.
\end{frame}

\begin{frame}{122. Las condiciones de borde necesitan atención}
\phantomsection\label{las-condiciones-de-borde-necesitan-atenciuxf3n}
\textbf{Libro guía:} John Semmlow, \emph{Circuits, Signals, and Systems
for Bioengineers}, 3.ª ed., 2018.

§4.6, \textbf{Time--Frequency Analysis} (página 205 del visor PDF).

\textbf{Correspondencia:} Actividad o interpretación biomédica de
elaboración docente apoyada en las secciones indicadas.

\textbf{Microcurrículo:} semana 9, sesión 2 (sesión global 18). Los
numerales del cronograma son identificadores curriculares; no son los
apartados del libro citados arriba.
\end{frame}

\begin{frame}{123. La escala logarítmica revela componentes débiles}
\phantomsection\label{la-escala-logaruxedtmica-revela-componentes-duxe9biles}
\textbf{Libro guía:} John Semmlow, \emph{Circuits, Signals, and Systems
for Bioengineers}, 3.ª ed., 2018.

§4.6, \textbf{Time--Frequency Analysis} (página 205 del visor PDF).

\textbf{Correspondencia:} Actividad o interpretación biomédica de
elaboración docente apoyada en las secciones indicadas.

\textbf{Microcurrículo:} semana 9, sesión 2 (sesión global 18). Los
numerales del cronograma son identificadores curriculares; no son los
apartados del libro citados arriba.
\end{frame}

\begin{frame}{124. STFT de una activación EMG}
\phantomsection\label{stft-de-una-activaciuxf3n-emg}
\textbf{Libro guía:} John Semmlow, \emph{Circuits, Signals, and Systems
for Bioengineers}, 3.ª ed., 2018.

§4.6, \textbf{Time--Frequency Analysis} (página 205 del visor PDF).

\textbf{Correspondencia:} Actividad o interpretación biomédica de
elaboración docente apoyada en las secciones indicadas.

\textbf{Microcurrículo:} semana 9, sesión 2 (sesión global 18). Los
numerales del cronograma son identificadores curriculares; no son los
apartados del libro citados arriba.
\end{frame}

\begin{frame}{125. STFT de una ráfaga EEG}
\phantomsection\label{stft-de-una-ruxe1faga-eeg}
\textbf{Libro guía:} John Semmlow, \emph{Circuits, Signals, and Systems
for Bioengineers}, 3.ª ed., 2018.

§4.6, \textbf{Time--Frequency Analysis} (página 205 del visor PDF).

\textbf{Correspondencia:} Actividad o interpretación biomédica de
elaboración docente apoyada en las secciones indicadas.

\textbf{Microcurrículo:} semana 9, sesión 2 (sesión global 18). Los
numerales del cronograma son identificadores curriculares; no son los
apartados del libro citados arriba.
\end{frame}

\begin{frame}{126. Código: STFT con densidad espectral}
\phantomsection\label{cuxf3digo-stft-con-densidad-espectral}
\textbf{Libro guía:} John Semmlow, \emph{Circuits, Signals, and Systems
for Bioengineers}, 3.ª ed., 2018.

§4.6, \textbf{Time--Frequency Analysis} (página 205 del visor PDF).

\textbf{Correspondencia:} Ejemplo, figura o implementación de
elaboración docente a partir de estos fundamentos. No es una
transcripción del código MATLAB del libro. Con scaling=``psd'', stft
devuelve aquí el semiespectro complejo. Su módulo cuadrado requiere
duplicar los bins interiores para combinar la potencia de las
frecuencias positivas y negativas.

\textbf{Microcurrículo:} semana 9, sesión 2 (sesión global 18). Los
numerales del cronograma son identificadores curriculares; no son los
apartados del libro citados arriba.

\textbf{Implementación:}
\href{https://docs.scipy.org/doc/scipy/reference/generated/scipy.signal.stft.html}{SciPy:
stft}. Función mantenida como legacy; se conserva para continuidad del
curso. Se explicita scaling y conversión de potencia unilateral.
\end{frame}

\begin{frame}{127. Validar un espectrograma}
\phantomsection\label{validar-un-espectrograma}
\textbf{Libro guía:} John Semmlow, \emph{Circuits, Signals, and Systems
for Bioengineers}, 3.ª ed., 2018.

§4.6, \textbf{Time--Frequency Analysis} (página 205 del visor PDF).

\textbf{Correspondencia:} Actividad o interpretación biomédica de
elaboración docente apoyada en las secciones indicadas.

\textbf{Microcurrículo:} semana 9, sesión 2 (sesión global 18). Los
numerales del cronograma son identificadores curriculares; no son los
apartados del libro citados arriba.
\end{frame}

\begin{frame}{128. Verificación de la sesión 18}
\phantomsection\label{verificaciuxf3n-de-la-sesiuxf3n-18}
\textbf{Libro guía:} John Semmlow, \emph{Circuits, Signals, and Systems
for Bioengineers}, 3.ª ed., 2018.

§4.6, \textbf{Time--Frequency Analysis} (página 205 del visor PDF).

\textbf{Correspondencia:} Actividad o interpretación biomédica de
elaboración docente apoyada en las secciones indicadas.

\textbf{Microcurrículo:} semana 9, sesión 2 (sesión global 18). Los
numerales del cronograma son identificadores curriculares; no son los
apartados del libro citados arriba.
\end{frame}

\begin{frame}{129. Cierre de la semana 9}
\phantomsection\label{cierre-de-la-semana-9}
\textbf{Libro guía:} John Semmlow, \emph{Circuits, Signals, and Systems
for Bioengineers}, 3.ª ed., 2018.

§4.4, \textbf{Spectral Averaging} (página 195 del visor PDF); §4.5,
\textbf{Signal Bandwidth} (página 200 del visor PDF); §4.6,
\textbf{Time--Frequency Analysis} (página 205 del visor PDF).

\textbf{Correspondencia:} Síntesis o encuadre que reúne varias
secciones; no corresponde a un único apartado del libro.

\textbf{Microcurrículo:} semana 9, sesión 2 (sesión global 18). Los
numerales del cronograma son identificadores curriculares; no son los
apartados del libro citados arriba.
\end{frame}

\begin{frame}{130. Semana 10 · Ventanas tiempo-frecuencia y sistemas
LTI}
\phantomsection\label{semana-10-ventanas-tiempo-frecuencia-y-sistemas-lti}
\textbf{Libro guía:} John Semmlow, \emph{Circuits, Signals, and Systems
for Bioengineers}, 3.ª ed., 2018.

§3.4, \textbf{Time--Frequency Transformation In The Discrete Domain}
(página 142 del visor PDF); §4.2, \textbf{Data Acquisition And Storage}
(página 175 del visor PDF); §4.6, \textbf{Time--Frequency Analysis}
(página 205 del visor PDF); §5.4, \textbf{Using The Impulse Response To
Find A Systems Output To Any Input---Convolution} (página 219 del visor
PDF).

\textbf{Correspondencia:} Síntesis o encuadre que reúne varias
secciones; no corresponde a un único apartado del libro.

\textbf{Lectura:} use estas secciones como ruta de preparación y
consulta. La bibliografía aparece al final del bloque.
\end{frame}

\begin{frame}{131. Sesión 19 · La ventana debe responder a la pregunta}
\phantomsection\label{sesiuxf3n-19-la-ventana-debe-responder-a-la-pregunta}
\textbf{Libro guía:} John Semmlow, \emph{Circuits, Signals, and Systems
for Bioengineers}, 3.ª ed., 2018.

§4.2.4, \textbf{Data Truncation---Window Functions} (página 188 del
visor PDF); §4.6, \textbf{Time--Frequency Analysis} (página 205 del
visor PDF).

\textbf{Correspondencia:} El tema se desarrolla en las secciones
indicadas. La redacción es una síntesis docente.

\textbf{Microcurrículo:} semana 10, sesión 1 (sesión global 19). Los
numerales del cronograma son identificadores curriculares; no son los
apartados del libro citados arriba.
\end{frame}

\begin{frame}{132. La longitud fija una escala de observación}
\phantomsection\label{la-longitud-fija-una-escala-de-observaciuxf3n}
\textbf{Libro guía:} John Semmlow, \emph{Circuits, Signals, and Systems
for Bioengineers}, 3.ª ed., 2018.

§4.2.4, \textbf{Data Truncation---Window Functions} (página 188 del
visor PDF); §4.6, \textbf{Time--Frequency Analysis} (página 205 del
visor PDF).

\textbf{Correspondencia:} El tema se desarrolla en las secciones
indicadas. La redacción es una síntesis docente.

\textbf{Microcurrículo:} semana 10, sesión 1 (sesión global 19). Los
numerales del cronograma son identificadores curriculares; no son los
apartados del libro citados arriba.
\end{frame}

\begin{frame}{133. Rectangular prioriza un lóbulo principal estrecho}
\phantomsection\label{rectangular-prioriza-un-luxf3bulo-principal-estrecho}
\textbf{Libro guía:} John Semmlow, \emph{Circuits, Signals, and Systems
for Bioengineers}, 3.ª ed., 2018.

§4.2.4, \textbf{Data Truncation---Window Functions} (página 188 del
visor PDF); §4.6, \textbf{Time--Frequency Analysis} (página 205 del
visor PDF).

\textbf{Correspondencia:} El tema se desarrolla en las secciones
indicadas. La redacción es una síntesis docente.

\textbf{Microcurrículo:} semana 10, sesión 1 (sesión global 19). Los
numerales del cronograma son identificadores curriculares; no son los
apartados del libro citados arriba.
\end{frame}

\begin{frame}{134. Hann ofrece un compromiso general}
\phantomsection\label{hann-ofrece-un-compromiso-general}
\textbf{Libro guía:} John Semmlow, \emph{Circuits, Signals, and Systems
for Bioengineers}, 3.ª ed., 2018.

§4.2.4, \textbf{Data Truncation---Window Functions} (página 188 del
visor PDF); §4.6, \textbf{Time--Frequency Analysis} (página 205 del
visor PDF).

\textbf{Correspondencia:} El tema se desarrolla en las secciones
indicadas. La redacción es una síntesis docente.

\textbf{Microcurrículo:} semana 10, sesión 1 (sesión global 19). Los
numerales del cronograma son identificadores curriculares; no son los
apartados del libro citados arriba.
\end{frame}

\begin{frame}{135. Hamming mantiene extremos pequeños, no nulos}
\phantomsection\label{hamming-mantiene-extremos-pequeuxf1os-no-nulos}
\textbf{Libro guía:} John Semmlow, \emph{Circuits, Signals, and Systems
for Bioengineers}, 3.ª ed., 2018.

§4.2.4, \textbf{Data Truncation---Window Functions} (página 188 del
visor PDF); §4.6, \textbf{Time--Frequency Analysis} (página 205 del
visor PDF).

\textbf{Correspondencia:} El tema se desarrolla en las secciones
indicadas. La redacción es una síntesis docente.

\textbf{Microcurrículo:} semana 10, sesión 1 (sesión global 19). Los
numerales del cronograma son identificadores curriculares; no son los
apartados del libro citados arriba.
\end{frame}

\begin{frame}{136. Blackman prioriza rango dinámico}
\phantomsection\label{blackman-prioriza-rango-dinuxe1mico}
\textbf{Libro guía:} John Semmlow, \emph{Circuits, Signals, and Systems
for Bioengineers}, 3.ª ed., 2018.

§4.2.4, \textbf{Data Truncation---Window Functions} (página 188 del
visor PDF); §4.6, \textbf{Time--Frequency Analysis} (página 205 del
visor PDF).

\textbf{Correspondencia:} El tema se desarrolla en las secciones
indicadas. La redacción es una síntesis docente.

\textbf{Microcurrículo:} semana 10, sesión 1 (sesión global 19). Los
numerales del cronograma son identificadores curriculares; no son los
apartados del libro citados arriba.
\end{frame}

\begin{frame}{137. Una ventana gaussiana localiza de forma equilibrada}
\phantomsection\label{una-ventana-gaussiana-localiza-de-forma-equilibrada}
\textbf{Libro guía:} John Semmlow, \emph{Circuits, Signals, and Systems
for Bioengineers}, 3.ª ed., 2018.

§4.2.4, \textbf{Data Truncation---Window Functions} (página 188 del
visor PDF); §4.6, \textbf{Time--Frequency Analysis} (página 205 del
visor PDF).

\textbf{Correspondencia:} El tema se desarrolla en las secciones
indicadas. La redacción es una síntesis docente.

\textbf{Microcurrículo:} semana 10, sesión 1 (sesión global 19). Los
numerales del cronograma son identificadores curriculares; no son los
apartados del libro citados arriba.
\end{frame}

\begin{frame}{138. Ventana y evento deben tener escalas compatibles}
\phantomsection\label{ventana-y-evento-deben-tener-escalas-compatibles}
\textbf{Libro guía:} John Semmlow, \emph{Circuits, Signals, and Systems
for Bioengineers}, 3.ª ed., 2018.

§4.6, \textbf{Time--Frequency Analysis} (página 205 del visor PDF).

\textbf{Correspondencia:} Actividad o interpretación biomédica de
elaboración docente apoyada en las secciones indicadas.

\textbf{Microcurrículo:} semana 10, sesión 1 (sesión global 19). Los
numerales del cronograma son identificadores curriculares; no son los
apartados del libro citados arriba.
\end{frame}

\begin{frame}{139. La frecuencia de interés determina ciclos observados}
\phantomsection\label{la-frecuencia-de-interuxe9s-determina-ciclos-observados}
\textbf{Libro guía:} John Semmlow, \emph{Circuits, Signals, and Systems
for Bioengineers}, 3.ª ed., 2018.

§4.6, \textbf{Time--Frequency Analysis} (página 205 del visor PDF).

\textbf{Correspondencia:} Actividad o interpretación biomédica de
elaboración docente apoyada en las secciones indicadas.

\textbf{Microcurrículo:} semana 10, sesión 1 (sesión global 19). Los
numerales del cronograma son identificadores curriculares; no son los
apartados del libro citados arriba.
\end{frame}

\begin{frame}{140. Solapamiento mejora seguimiento visual}
\phantomsection\label{solapamiento-mejora-seguimiento-visual}
\textbf{Libro guía:} John Semmlow, \emph{Circuits, Signals, and Systems
for Bioengineers}, 3.ª ed., 2018.

§4.6, \textbf{Time--Frequency Analysis} (página 205 del visor PDF).

\textbf{Correspondencia:} Actividad o interpretación biomédica de
elaboración docente apoyada en las secciones indicadas.

\textbf{Microcurrículo:} semana 10, sesión 1 (sesión global 19). Los
numerales del cronograma son identificadores curriculares; no son los
apartados del libro citados arriba.
\end{frame}

\begin{frame}{141. Zero padding facilita localizar máximos}
\phantomsection\label{zero-padding-facilita-localizar-muxe1ximos}
\textbf{Libro guía:} John Semmlow, \emph{Circuits, Signals, and Systems
for Bioengineers}, 3.ª ed., 2018.

§4.6, \textbf{Time--Frequency Analysis} (página 205 del visor PDF).

\textbf{Correspondencia:} Actividad o interpretación biomédica de
elaboración docente apoyada en las secciones indicadas.

\textbf{Microcurrículo:} semana 10, sesión 1 (sesión global 19). Los
numerales del cronograma son identificadores curriculares; no son los
apartados del libro citados arriba.
\end{frame}

\begin{frame}{142. Normalizar entre ventanas requiere coherencia}
\phantomsection\label{normalizar-entre-ventanas-requiere-coherencia}
\textbf{Libro guía:} John Semmlow, \emph{Circuits, Signals, and Systems
for Bioengineers}, 3.ª ed., 2018.

§4.6, \textbf{Time--Frequency Analysis} (página 205 del visor PDF).

\textbf{Correspondencia:} Actividad o interpretación biomédica de
elaboración docente apoyada en las secciones indicadas.

\textbf{Microcurrículo:} semana 10, sesión 1 (sesión global 19). Los
numerales del cronograma son identificadores curriculares; no son los
apartados del libro citados arriba.
\end{frame}

\begin{frame}{143. Aplicación: activación muscular durante marcha}
\phantomsection\label{aplicaciuxf3n-activaciuxf3n-muscular-durante-marcha}
\textbf{Libro guía:} John Semmlow, \emph{Circuits, Signals, and Systems
for Bioengineers}, 3.ª ed., 2018.

§4.6, \textbf{Time--Frequency Analysis} (página 205 del visor PDF).

\textbf{Correspondencia:} Actividad o interpretación biomédica de
elaboración docente apoyada en las secciones indicadas.

\textbf{Microcurrículo:} semana 10, sesión 1 (sesión global 19). Los
numerales del cronograma son identificadores curriculares; no son los
apartados del libro citados arriba.
\end{frame}

\begin{frame}{144. Aplicación: cambio de ritmo respiratorio}
\phantomsection\label{aplicaciuxf3n-cambio-de-ritmo-respiratorio}
\textbf{Libro guía:} John Semmlow, \emph{Circuits, Signals, and Systems
for Bioengineers}, 3.ª ed., 2018.

§4.6, \textbf{Time--Frequency Analysis} (página 205 del visor PDF).

\textbf{Correspondencia:} Actividad o interpretación biomédica de
elaboración docente apoyada en las secciones indicadas.

\textbf{Microcurrículo:} semana 10, sesión 1 (sesión global 19). Los
numerales del cronograma son identificadores curriculares; no son los
apartados del libro citados arriba.
\end{frame}

\begin{frame}{145. Flujo de selección documentada}
\phantomsection\label{flujo-de-selecciuxf3n-documentada}
\textbf{Libro guía:} John Semmlow, \emph{Circuits, Signals, and Systems
for Bioengineers}, 3.ª ed., 2018.

§4.6, \textbf{Time--Frequency Analysis} (página 205 del visor PDF).

\textbf{Correspondencia:} Actividad o interpretación biomédica de
elaboración docente apoyada en las secciones indicadas.

\textbf{Microcurrículo:} semana 10, sesión 1 (sesión global 19). Los
numerales del cronograma son identificadores curriculares; no son los
apartados del libro citados arriba.
\end{frame}

\begin{frame}{146. Verificación de la sesión 19}
\phantomsection\label{verificaciuxf3n-de-la-sesiuxf3n-19}
\textbf{Libro guía:} John Semmlow, \emph{Circuits, Signals, and Systems
for Bioengineers}, 3.ª ed., 2018.

§4.2.4, \textbf{Data Truncation---Window Functions} (página 188 del
visor PDF); §4.6, \textbf{Time--Frequency Analysis} (página 205 del
visor PDF).

\textbf{Correspondencia:} El tema se desarrolla en las secciones
indicadas. La redacción es una síntesis docente.

\textbf{Microcurrículo:} semana 10, sesión 1 (sesión global 19). Los
numerales del cronograma son identificadores curriculares; no son los
apartados del libro citados arriba.
\end{frame}

\begin{frame}{147. Sesión 20 · Un sistema transforma una señal}
\phantomsection\label{sesiuxf3n-20-un-sistema-transforma-una-seuxf1al}
\textbf{Libro guía:} John Semmlow, \emph{Circuits, Signals, and Systems
for Bioengineers}, 3.ª ed., 2018.

§5.2, \textbf{Linear Systems Analysis---An Overview} (página 212 del
visor PDF).

\textbf{Correspondencia:} El tema se desarrolla en las secciones
indicadas. La redacción es una síntesis docente.

\textbf{Microcurrículo:} semana 10, sesión 2 (sesión global 20). Los
numerales del cronograma son identificadores curriculares; no son los
apartados del libro citados arriba.
\end{frame}

\begin{frame}{148. Linealidad combina aditividad y homogeneidad}
\phantomsection\label{linealidad-combina-aditividad-y-homogeneidad}
\textbf{Libro guía:} John Semmlow, \emph{Circuits, Signals, and Systems
for Bioengineers}, 3.ª ed., 2018.

§5.2.1, \textbf{Superposition And Linearity} (página 213 del visor PDF).

\textbf{Correspondencia:} El tema se desarrolla en las secciones
indicadas. La redacción es una síntesis docente.

\textbf{Microcurrículo:} semana 10, sesión 2 (sesión global 20). Los
numerales del cronograma son identificadores curriculares; no son los
apartados del libro citados arriba.
\end{frame}

\begin{frame}{149. Saturación viola linealidad}
\phantomsection\label{saturaciuxf3n-viola-linealidad}
\textbf{Libro guía:} John Semmlow, \emph{Circuits, Signals, and Systems
for Bioengineers}, 3.ª ed., 2018.

§1.4.4, \textbf{Linear Versus Nonlinear Signals And Systems} (página 44
del visor PDF); §5.2.1, \textbf{Superposition And Linearity} (página 213
del visor PDF).

\textbf{Correspondencia:} El tema se desarrolla en las secciones
indicadas. La redacción es una síntesis docente.

\textbf{Microcurrículo:} semana 10, sesión 2 (sesión global 20). Los
numerales del cronograma son identificadores curriculares; no son los
apartados del libro citados arriba.
\end{frame}

\begin{frame}{150. Invariancia temporal conserva la regla}
\phantomsection\label{invariancia-temporal-conserva-la-regla}
\textbf{Libro guía:} John Semmlow, \emph{Circuits, Signals, and Systems
for Bioengineers}, 3.ª ed., 2018.

§1.4.4, \textbf{Linear Versus Nonlinear Signals And Systems} (página 44
del visor PDF); §5.2, \textbf{Linear Systems Analysis---An Overview}
(página 212 del visor PDF).

\textbf{Correspondencia:} Las secciones aportan el fundamento. La
diapositiva amplía la formulación o precisa condiciones que no se
presentan con idéntica notación en el libro. La invariancia temporal se
comprueba separadamente de la linealidad.

\textbf{Microcurrículo:} semana 10, sesión 2 (sesión global 20). Los
numerales del cronograma son identificadores curriculares; no son los
apartados del libro citados arriba.
\end{frame}

\begin{frame}{151. Adaptación fisiológica rompe invariancia}
\phantomsection\label{adaptaciuxf3n-fisioluxf3gica-rompe-invariancia}
\textbf{Libro guía:} John Semmlow, \emph{Circuits, Signals, and Systems
for Bioengineers}, 3.ª ed., 2018.

§1.4.4, \textbf{Linear Versus Nonlinear Signals And Systems} (página 44
del visor PDF); §5.2, \textbf{Linear Systems Analysis---An Overview}
(página 212 del visor PDF).

\textbf{Correspondencia:} Las secciones aportan el fundamento. La
diapositiva amplía la formulación o precisa condiciones que no se
presentan con idéntica notación en el libro. La invariancia temporal se
comprueba separadamente de la linealidad.

\textbf{Microcurrículo:} semana 10, sesión 2 (sesión global 20). Los
numerales del cronograma son identificadores curriculares; no son los
apartados del libro citados arriba.
\end{frame}

\begin{frame}{152. Memoria y causalidad son propiedades distintas}
\phantomsection\label{memoria-y-causalidad-son-propiedades-distintas}
\textbf{Libro guía:} John Semmlow, \emph{Circuits, Signals, and Systems
for Bioengineers}, 3.ª ed., 2018.

§1.4.3, \textbf{Causal Versus Noncausal Signals And Systems} (página 43
del visor PDF); §8.4.3, \textbf{Causal And Noncausal Filters} (página
365 del visor PDF).

\textbf{Correspondencia:} El tema se desarrolla en las secciones
indicadas. La redacción es una síntesis docente.

\textbf{Microcurrículo:} semana 10, sesión 2 (sesión global 20). Los
numerales del cronograma son identificadores curriculares; no son los
apartados del libro citados arriba.
\end{frame}

\begin{frame}{153. Estabilidad BIBO limita la salida}
\phantomsection\label{estabilidad-bibo-limita-la-salida}
\textbf{Libro guía:} John Semmlow, \emph{Circuits, Signals, and Systems
for Bioengineers}, 3.ª ed., 2018.

§5.4, \textbf{Using The Impulse Response To Find A Systems Output To Any
Input---Convolution} (página 219 del visor PDF); §7.5, \textbf{The
Laplace Domain, The Frequency Domain, And The Time Domain} (página 330
del visor PDF).

\textbf{Correspondencia:} Las secciones aportan el fundamento. La
diapositiva amplía la formulación o precisa condiciones que no se
presentan con idéntica notación en el libro. Se explicita el criterio
BIBO y el caso de término directo distribucional.

\textbf{Microcurrículo:} semana 10, sesión 2 (sesión global 20). Los
numerales del cronograma son identificadores curriculares; no son los
apartados del libro citados arriba.
\end{frame}

\begin{frame}{154. El impulso ideal concentra área unitaria}
\phantomsection\label{el-impulso-ideal-concentra-uxe1rea-unitaria}
\textbf{Libro guía:} John Semmlow, \emph{Circuits, Signals, and Systems
for Bioengineers}, 3.ª ed., 2018.

§5.3, \textbf{A Slice In Time: The Impulse Signal} (página 214 del visor
PDF); §5.3.1, \textbf{Real-World Impulse Signals} (página 215 del visor
PDF).

\textbf{Correspondencia:} El tema se desarrolla en las secciones
indicadas. La redacción es una síntesis docente.

\textbf{Microcurrículo:} semana 10, sesión 2 (sesión global 20). Los
numerales del cronograma son identificadores curriculares; no son los
apartados del libro citados arriba.
\end{frame}

\begin{frame}{155. En discreto, el impulso es una secuencia ordinaria}
\phantomsection\label{en-discreto-el-impulso-es-una-secuencia-ordinaria}
\textbf{Libro guía:} John Semmlow, \emph{Circuits, Signals, and Systems
for Bioengineers}, 3.ª ed., 2018.

§5.3, \textbf{A Slice In Time: The Impulse Signal} (página 214 del visor
PDF); §5.4, \textbf{Using The Impulse Response To Find A Systems Output
To Any Input---Convolution} (página 219 del visor PDF).

\textbf{Correspondencia:} El tema se desarrolla en las secciones
indicadas. La redacción es una síntesis docente.

\textbf{Microcurrículo:} semana 10, sesión 2 (sesión global 20). Los
numerales del cronograma son identificadores curriculares; no son los
apartados del libro citados arriba.
\end{frame}

\begin{frame}{156. La respuesta al impulso caracteriza un sistema LTI}
\phantomsection\label{la-respuesta-al-impulso-caracteriza-un-sistema-lti}
\textbf{Libro guía:} John Semmlow, \emph{Circuits, Signals, and Systems
for Bioengineers}, 3.ª ed., 2018.

§5.4, \textbf{Using The Impulse Response To Find A Systems Output To Any
Input---Convolution} (página 219 del visor PDF).

\textbf{Correspondencia:} El tema se desarrolla en las secciones
indicadas. La redacción es una síntesis docente.

\textbf{Microcurrículo:} semana 10, sesión 2 (sesión global 20). Los
numerales del cronograma son identificadores curriculares; no son los
apartados del libro citados arriba.
\end{frame}

\begin{frame}{157. La salida es una convolución}
\phantomsection\label{la-salida-es-una-convoluciuxf3n}
\textbf{Libro guía:} John Semmlow, \emph{Circuits, Signals, and Systems
for Bioengineers}, 3.ª ed., 2018.

§5.4, \textbf{Using The Impulse Response To Find A Systems Output To Any
Input---Convolution} (página 219 del visor PDF).

\textbf{Correspondencia:} El tema se desarrolla en las secciones
indicadas. La redacción es una síntesis docente.

\textbf{Microcurrículo:} semana 10, sesión 2 (sesión global 20). Los
numerales del cronograma son identificadores curriculares; no son los
apartados del libro citados arriba.
\end{frame}

\begin{frame}{158. Causalidad de un LTI se lee en \(h\)}
\phantomsection\label{causalidad-de-un-lti-se-lee-en-h}
\textbf{Libro guía:} John Semmlow, \emph{Circuits, Signals, and Systems
for Bioengineers}, 3.ª ed., 2018.

§1.4.3, \textbf{Causal Versus Noncausal Signals And Systems} (página 43
del visor PDF); §5.4, \textbf{Using The Impulse Response To Find A
Systems Output To Any Input---Convolution} (página 219 del visor PDF).

\textbf{Correspondencia:} El tema se desarrolla en las secciones
indicadas. La redacción es una síntesis docente.

\textbf{Microcurrículo:} semana 10, sesión 2 (sesión global 20). Los
numerales del cronograma son identificadores curriculares; no son los
apartados del libro citados arriba.
\end{frame}

\begin{frame}{159. Las exponenciales complejas son autofunciones LTI}
\phantomsection\label{las-exponenciales-complejas-son-autofunciones-lti}
\textbf{Libro guía:} John Semmlow, \emph{Circuits, Signals, and Systems
for Bioengineers}, 3.ª ed., 2018.

§6.3, \textbf{The Response Of System Elements To Sinusoidal Inputs:
Phasor Analysis} (página 251 del visor PDF).

\textbf{Correspondencia:} El tema se desarrolla en las secciones
indicadas. La redacción es una síntesis docente.

\textbf{Microcurrículo:} semana 10, sesión 2 (sesión global 20). Los
numerales del cronograma son identificadores curriculares; no son los
apartados del libro citados arriba.
\end{frame}

\begin{frame}{160. \(H(f)\) es la transformada de \(h(t)\)}
\phantomsection\label{hf-es-la-transformada-de-ht}
\textbf{Libro guía:} John Semmlow, \emph{Circuits, Signals, and Systems
for Bioengineers}, 3.ª ed., 2018.

§5.6, \textbf{Convolution In The Frequency Domain} (página 237 del visor
PDF); §6.7, \textbf{The Transfer Function And The Fourier Transform}
(página 287 del visor PDF).

\textbf{Correspondencia:} El tema se desarrolla en las secciones
indicadas. La redacción es una síntesis docente.

\textbf{Microcurrículo:} semana 10, sesión 2 (sesión global 20). Los
numerales del cronograma son identificadores curriculares; no son los
apartados del libro citados arriba.
\end{frame}

\begin{frame}{161. Ejemplos de sistemas de procesamiento}
\phantomsection\label{ejemplos-de-sistemas-de-procesamiento}
\textbf{Libro guía:} John Semmlow, \emph{Circuits, Signals, and Systems
for Bioengineers}, 3.ª ed., 2018.

§5.5, \textbf{Applied Convolution---Basic Filters} (página 231 del visor
PDF); §6.5, \textbf{The Spectrum Of System Elements: The Bode Plot}
(página 263 del visor PDF).

\textbf{Correspondencia:} El tema se desarrolla en las secciones
indicadas. La redacción es una síntesis docente.

\textbf{Microcurrículo:} semana 10, sesión 2 (sesión global 20). Los
numerales del cronograma son identificadores curriculares; no son los
apartados del libro citados arriba.
\end{frame}

\begin{frame}{162. Código: probar linealidad numéricamente}
\phantomsection\label{cuxf3digo-probar-linealidad-numuxe9ricamente}
\textbf{Libro guía:} John Semmlow, \emph{Circuits, Signals, and Systems
for Bioengineers}, 3.ª ed., 2018.

§5.2.1, \textbf{Superposition And Linearity} (página 213 del visor PDF).

\textbf{Correspondencia:} Ejemplo, figura o implementación de
elaboración docente a partir de estos fundamentos. No es una
transcripción del código MATLAB del libro.

\textbf{Microcurrículo:} semana 10, sesión 2 (sesión global 20). Los
numerales del cronograma son identificadores curriculares; no son los
apartados del libro citados arriba.
\end{frame}

\begin{frame}{163. Verificación de la sesión 20}
\phantomsection\label{verificaciuxf3n-de-la-sesiuxf3n-20}
\textbf{Libro guía:} John Semmlow, \emph{Circuits, Signals, and Systems
for Bioengineers}, 3.ª ed., 2018.

§5.4, \textbf{Using The Impulse Response To Find A Systems Output To Any
Input---Convolution} (página 219 del visor PDF); §5.5, \textbf{Applied
Convolution---Basic Filters} (página 231 del visor PDF).

\textbf{Correspondencia:} Actividad o interpretación biomédica de
elaboración docente apoyada en las secciones indicadas.

\textbf{Microcurrículo:} semana 10, sesión 2 (sesión global 20). Los
numerales del cronograma son identificadores curriculares; no son los
apartados del libro citados arriba.
\end{frame}

\begin{frame}{164. Cierre de la semana 10}
\phantomsection\label{cierre-de-la-semana-10}
\textbf{Libro guía:} John Semmlow, \emph{Circuits, Signals, and Systems
for Bioengineers}, 3.ª ed., 2018.

§4.6, \textbf{Time--Frequency Analysis} (página 205 del visor PDF);
§5.4, \textbf{Using The Impulse Response To Find A Systems Output To Any
Input---Convolution} (página 219 del visor PDF).

\textbf{Correspondencia:} Síntesis o encuadre que reúne varias
secciones; no corresponde a un único apartado del libro.

\textbf{Microcurrículo:} semana 10, sesión 2 (sesión global 20). Los
numerales del cronograma son identificadores curriculares; no son los
apartados del libro citados arriba.
\end{frame}

\begin{frame}{165. Síntesis y estudio autónomo}
\phantomsection\label{suxedntesis-y-estudio-autuxf3nomo-1}
\textbf{Libro guía:} John Semmlow, \emph{Circuits, Signals, and Systems
for Bioengineers}, 3.ª ed., 2018.

§3.4, \textbf{Time--Frequency Transformation In The Discrete Domain}
(página 142 del visor PDF); §4.2, \textbf{Data Acquisition And Storage}
(página 175 del visor PDF); §4.6, \textbf{Time--Frequency Analysis}
(página 205 del visor PDF); §5.4, \textbf{Using The Impulse Response To
Find A Systems Output To Any Input---Convolution} (página 219 del visor
PDF).

\textbf{Correspondencia:} Síntesis o encuadre que reúne varias
secciones; no corresponde a un único apartado del libro.

\textbf{Lectura:} use estas secciones como ruta de preparación y
consulta. La bibliografía aparece al final del bloque.
\end{frame}

\begin{frame}{166. La cadena completa de decisiones}
\phantomsection\label{la-cadena-completa-de-decisiones}
\textbf{Libro guía:} John Semmlow, \emph{Circuits, Signals, and Systems
for Bioengineers}, 3.ª ed., 2018.

§3.4, \textbf{Time--Frequency Transformation In The Discrete Domain}
(página 142 del visor PDF); §4.2, \textbf{Data Acquisition And Storage}
(página 175 del visor PDF); §4.6, \textbf{Time--Frequency Analysis}
(página 205 del visor PDF); §5.4, \textbf{Using The Impulse Response To
Find A Systems Output To Any Input---Convolution} (página 219 del visor
PDF).

\textbf{Correspondencia:} Síntesis o encuadre que reúne varias
secciones; no corresponde a un único apartado del libro.

\textbf{Lectura:} use estas secciones como ruta de preparación y
consulta. La bibliografía aparece al final del bloque.
\end{frame}

\begin{frame}{167. Errores que deben poder diagnosticarse}
\phantomsection\label{errores-que-deben-poder-diagnosticarse}
\textbf{Libro guía:} John Semmlow, \emph{Circuits, Signals, and Systems
for Bioengineers}, 3.ª ed., 2018.

§3.4, \textbf{Time--Frequency Transformation In The Discrete Domain}
(página 142 del visor PDF); §4.2, \textbf{Data Acquisition And Storage}
(página 175 del visor PDF); §4.6, \textbf{Time--Frequency Analysis}
(página 205 del visor PDF); §5.4, \textbf{Using The Impulse Response To
Find A Systems Output To Any Input---Convolution} (página 219 del visor
PDF).

\textbf{Correspondencia:} Síntesis o encuadre que reúne varias
secciones; no corresponde a un único apartado del libro.

\textbf{Lectura:} use estas secciones como ruta de preparación y
consulta. La bibliografía aparece al final del bloque.
\end{frame}

\begin{frame}{168. Autoevaluación acumulativa}
\phantomsection\label{autoevaluaciuxf3n-acumulativa-1}
\textbf{Libro guía:} John Semmlow, \emph{Circuits, Signals, and Systems
for Bioengineers}, 3.ª ed., 2018.

§3.4, \textbf{Time--Frequency Transformation In The Discrete Domain}
(página 142 del visor PDF); §4.2, \textbf{Data Acquisition And Storage}
(página 175 del visor PDF); §4.6, \textbf{Time--Frequency Analysis}
(página 205 del visor PDF); §5.4, \textbf{Using The Impulse Response To
Find A Systems Output To Any Input---Convolution} (página 219 del visor
PDF).

\textbf{Correspondencia:} Síntesis o encuadre que reúne varias
secciones; no corresponde a un único apartado del libro.

\textbf{Lectura:} use estas secciones como ruta de preparación y
consulta. La bibliografía aparece al final del bloque.
\end{frame}

\begin{frame}{169. Actividad integradora sugerida}
\phantomsection\label{actividad-integradora-sugerida-1}
\textbf{Libro guía:} John Semmlow, \emph{Circuits, Signals, and Systems
for Bioengineers}, 3.ª ed., 2018.

§3.4, \textbf{Time--Frequency Transformation In The Discrete Domain}
(página 142 del visor PDF); §4.2, \textbf{Data Acquisition And Storage}
(página 175 del visor PDF); §4.6, \textbf{Time--Frequency Analysis}
(página 205 del visor PDF); §5.4, \textbf{Using The Impulse Response To
Find A Systems Output To Any Input---Convolution} (página 219 del visor
PDF).

\textbf{Correspondencia:} Síntesis o encuadre que reúne varias
secciones; no corresponde a un único apartado del libro.

\textbf{Lectura:} use estas secciones como ruta de preparación y
consulta. La bibliografía aparece al final del bloque.
\end{frame}

\begin{frame}{170. Referencias y alcance}
\phantomsection\label{referencias-y-alcance-1}
\textbf{Libro guía:} John Semmlow, \emph{Circuits, Signals, and Systems
for Bioengineers}, 3.ª ed., 2018.

§3.4, \textbf{Time--Frequency Transformation In The Discrete Domain}
(página 142 del visor PDF); §4.2, \textbf{Data Acquisition And Storage}
(página 175 del visor PDF); §4.6, \textbf{Time--Frequency Analysis}
(página 205 del visor PDF); §5.4, \textbf{Using The Impulse Response To
Find A Systems Output To Any Input---Convolution} (página 219 del visor
PDF).

\textbf{Correspondencia:} Síntesis o encuadre que reúne varias
secciones; no corresponde a un único apartado del libro.

\textbf{Lectura:} use estas secciones como ruta de preparación y
consulta. La bibliografía aparece al final del bloque.
\end{frame}

\begin{frame}{171. Cierre: criterio de dominio}
\phantomsection\label{cierre-criterio-de-dominio-1}
\textbf{Libro guía:} John Semmlow, \emph{Circuits, Signals, and Systems
for Bioengineers}, 3.ª ed., 2018.

§3.4, \textbf{Time--Frequency Transformation In The Discrete Domain}
(página 142 del visor PDF); §4.2, \textbf{Data Acquisition And Storage}
(página 175 del visor PDF); §4.6, \textbf{Time--Frequency Analysis}
(página 205 del visor PDF); §5.4, \textbf{Using The Impulse Response To
Find A Systems Output To Any Input---Convolution} (página 219 del visor
PDF).

\textbf{Correspondencia:} Síntesis o encuadre que reúne varias
secciones; no corresponde a un único apartado del libro.

\textbf{Lectura:} use estas secciones como ruta de preparación y
consulta. La bibliografía aparece al final del bloque.
\end{frame}

\begin{frame}{172. Referencias}
\phantomsection\label{referencias-1}
\textbf{Libro guía:} John Semmlow, \emph{Circuits, Signals, and Systems
for Bioengineers}, 3.ª ed., 2018.

§3.4, \textbf{Time--Frequency Transformation In The Discrete Domain}
(página 142 del visor PDF); §4.2, \textbf{Data Acquisition And Storage}
(página 175 del visor PDF); §4.6, \textbf{Time--Frequency Analysis}
(página 205 del visor PDF); §5.4, \textbf{Using The Impulse Response To
Find A Systems Output To Any Input---Convolution} (página 219 del visor
PDF).

\textbf{Correspondencia:} Síntesis o encuadre que reúne varias
secciones; no corresponde a un único apartado del libro.

\textbf{Lectura:} use estas secciones como ruta de preparación y
consulta. La bibliografía aparece al final del bloque.
\end{frame}

\section{semanas-11-15.qmd}\label{semanas-11-15.qmd}

\begin{frame}{1. Sistemas y Señales Biomédicos}
\phantomsection\label{sistemas-y-seuxf1ales-biomuxe9dicos-2}
\textbf{Libro guía:} John Semmlow, \emph{Circuits, Signals, and Systems
for Bioengineers}, 3.ª ed., 2018.

§5.4, \textbf{Using The Impulse Response To Find A Systems Output To Any
Input---Convolution} (página 219 del visor PDF); §6.4, \textbf{The
Transfer Function} (página 255 del visor PDF); §7.4, \textbf{Nonzero
Initial Conditions---Initial And Final Value Theorems} (página 326 del
visor PDF); §8.4, \textbf{Linear Filters---Introduction} (página 358 del
visor PDF).

\textbf{Correspondencia:} Síntesis o encuadre que reúne varias
secciones; no corresponde a un único apartado del libro. Lectura
principal: Semmlow (2018), tercera edición. La correspondencia usa la
edición del archivo aportado, que difiere de la segunda edición de 2012
citada en el microcurrículo.

\textbf{Lectura:} use estas secciones como ruta de preparación y
consulta. La bibliografía aparece al final del bloque.
\end{frame}

\begin{frame}{2. Cómo estudiar las semanas 11--15}
\phantomsection\label{cuxf3mo-estudiar-las-semanas-1115}
\textbf{Libro guía:} John Semmlow, \emph{Circuits, Signals, and Systems
for Bioengineers}, 3.ª ed., 2018.

§5.4, \textbf{Using The Impulse Response To Find A Systems Output To Any
Input---Convolution} (página 219 del visor PDF); §6.4, \textbf{The
Transfer Function} (página 255 del visor PDF); §7.4, \textbf{Nonzero
Initial Conditions---Initial And Final Value Theorems} (página 326 del
visor PDF); §8.4, \textbf{Linear Filters---Introduction} (página 358 del
visor PDF).

\textbf{Correspondencia:} Síntesis o encuadre que reúne varias
secciones; no corresponde a un único apartado del libro.

\textbf{Lectura:} use estas secciones como ruta de preparación y
consulta. La bibliografía aparece al final del bloque.
\end{frame}

\begin{frame}{3. Del modelo LTI a representaciones para clasificación}
\phantomsection\label{del-modelo-lti-a-representaciones-para-clasificaciuxf3n}
\textbf{Libro guía:} John Semmlow, \emph{Circuits, Signals, and Systems
for Bioengineers}, 3.ª ed., 2018.

§5.4, \textbf{Using The Impulse Response To Find A Systems Output To Any
Input---Convolution} (página 219 del visor PDF); §6.4, \textbf{The
Transfer Function} (página 255 del visor PDF); §7.4, \textbf{Nonzero
Initial Conditions---Initial And Final Value Theorems} (página 326 del
visor PDF); §8.4, \textbf{Linear Filters---Introduction} (página 358 del
visor PDF).

\textbf{Correspondencia:} Síntesis o encuadre que reúne varias
secciones; no corresponde a un único apartado del libro.

\textbf{Lectura:} use estas secciones como ruta de preparación y
consulta. La bibliografía aparece al final del bloque.
\end{frame}

\begin{frame}{4. Mapa de las diez sesiones}
\phantomsection\label{mapa-de-las-diez-sesiones-2}
\textbf{Libro guía:} John Semmlow, \emph{Circuits, Signals, and Systems
for Bioengineers}, 3.ª ed., 2018.

§5.4, \textbf{Using The Impulse Response To Find A Systems Output To Any
Input---Convolution} (página 219 del visor PDF); §6.4, \textbf{The
Transfer Function} (página 255 del visor PDF); §7.4, \textbf{Nonzero
Initial Conditions---Initial And Final Value Theorems} (página 326 del
visor PDF); §8.4, \textbf{Linear Filters---Introduction} (página 358 del
visor PDF).

\textbf{Correspondencia:} Síntesis o encuadre que reúne varias
secciones; no corresponde a un único apartado del libro.

\textbf{Lectura:} use estas secciones como ruta de preparación y
consulta. La bibliografía aparece al final del bloque.
\end{frame}

\begin{frame}{5. Organización de cada sesión y trabajo independiente}
\phantomsection\label{organizaciuxf3n-de-cada-sesiuxf3n-y-trabajo-independiente-2}
\textbf{Libro guía:} John Semmlow, \emph{Circuits, Signals, and Systems
for Bioengineers}, 3.ª ed., 2018.

§1.1.1, \textbf{Goals Of This Book} (página 10 del visor PDF).

\textbf{Correspondencia:} Organización docente basada en las actividades
y RAE del microcurrículo. No corresponde a una sección temática del
libro. Microcurrículo: actividades de aprendizaje y RAE, PDF 2--3.
Distribución docente propuesta para las sesiones de 90 minutos.

\textbf{Lectura:} use estas secciones como ruta de preparación y
consulta. La bibliografía aparece al final del bloque.
\end{frame}

\begin{frame}{6. Evidencias de aprendizaje y criterios de revisión}
\phantomsection\label{evidencias-de-aprendizaje-y-criterios-de-revisiuxf3n-2}
\textbf{Libro guía:} John Semmlow, \emph{Circuits, Signals, and Systems
for Bioengineers}, 3.ª ed., 2018.

§1.1.1, \textbf{Goals Of This Book} (página 10 del visor PDF).

\textbf{Correspondencia:} Organización docente basada en las actividades
y RAE del microcurrículo. No corresponde a una sección temática del
libro. Microcurrículo: RAE y actividades de evaluación, PDF 2--3. No se
resuelven aquí las inconsistencias administrativas del documento.

\textbf{Lectura:} use estas secciones como ruta de preparación y
consulta. La bibliografía aparece al final del bloque.
\end{frame}

\begin{frame}{7. Convenciones de los ejemplos en Python}
\phantomsection\label{convenciones-de-los-ejemplos-en-python-2}
\textbf{Libro guía:} John Semmlow, \emph{Circuits, Signals, and Systems
for Bioengineers}, 3.ª ed., 2018.

§1.1.1, \textbf{Goals Of This Book} (página 10 del visor PDF).

\textbf{Correspondencia:} Ejemplo, figura o implementación de
elaboración docente a partir de estos fundamentos. No es una
transcripción del código MATLAB del libro. El lenguaje del
microcurrículo suministrado es Python. El libro guía utiliza MATLAB; se
conservan conceptos y se explicitan diferencias de convenciones.

\textbf{Lectura:} use estas secciones como ruta de preparación y
consulta. La bibliografía aparece al final del bloque.
\end{frame}

\begin{frame}{8. Continuidad con las semanas 6--10}
\phantomsection\label{continuidad-con-las-semanas-610}
\textbf{Libro guía:} John Semmlow, \emph{Circuits, Signals, and Systems
for Bioengineers}, 3.ª ed., 2018.

§5.4, \textbf{Using The Impulse Response To Find A Systems Output To Any
Input---Convolution} (página 219 del visor PDF); §6.4, \textbf{The
Transfer Function} (página 255 del visor PDF); §7.4, \textbf{Nonzero
Initial Conditions---Initial And Final Value Theorems} (página 326 del
visor PDF); §8.4, \textbf{Linear Filters---Introduction} (página 358 del
visor PDF).

\textbf{Correspondencia:} Síntesis o encuadre que reúne varias
secciones; no corresponde a un único apartado del libro.

\textbf{Lectura:} use estas secciones como ruta de preparación y
consulta. La bibliografía aparece al final del bloque.
\end{frame}

\begin{frame}{9. Ruta de comprobación}
\phantomsection\label{ruta-de-comprobaciuxf3n}
\textbf{Libro guía:} John Semmlow, \emph{Circuits, Signals, and Systems
for Bioengineers}, 3.ª ed., 2018.

§5.4, \textbf{Using The Impulse Response To Find A Systems Output To Any
Input---Convolution} (página 219 del visor PDF); §6.4, \textbf{The
Transfer Function} (página 255 del visor PDF); §7.4, \textbf{Nonzero
Initial Conditions---Initial And Final Value Theorems} (página 326 del
visor PDF); §8.4, \textbf{Linear Filters---Introduction} (página 358 del
visor PDF).

\textbf{Correspondencia:} Síntesis o encuadre que reúne varias
secciones; no corresponde a un único apartado del libro.

\textbf{Lectura:} use estas secciones como ruta de preparación y
consulta. La bibliografía aparece al final del bloque.
\end{frame}

\begin{frame}{10. Semana 11 · Convolución en tiempo y sistemas
aplicados}
\phantomsection\label{semana-11-convoluciuxf3n-en-tiempo-y-sistemas-aplicados}
\textbf{Libro guía:} John Semmlow, \emph{Circuits, Signals, and Systems
for Bioengineers}, 3.ª ed., 2018.

§5.4, \textbf{Using The Impulse Response To Find A Systems Output To Any
Input---Convolution} (página 219 del visor PDF); §6.4, \textbf{The
Transfer Function} (página 255 del visor PDF); §7.4, \textbf{Nonzero
Initial Conditions---Initial And Final Value Theorems} (página 326 del
visor PDF); §8.4, \textbf{Linear Filters---Introduction} (página 358 del
visor PDF).

\textbf{Correspondencia:} Síntesis o encuadre que reúne varias
secciones; no corresponde a un único apartado del libro.

\textbf{Lectura:} use estas secciones como ruta de preparación y
consulta. La bibliografía aparece al final del bloque.
\end{frame}

\begin{frame}{11. Sesión 21 · Cálculo de la convolución}
\phantomsection\label{sesiuxf3n-21-cuxe1lculo-de-la-convoluciuxf3n}
\textbf{Libro guía:} John Semmlow, \emph{Circuits, Signals, and Systems
for Bioengineers}, 3.ª ed., 2018.

§5.4, \textbf{Using The Impulse Response To Find A Systems Output To Any
Input---Convolution} (página 219 del visor PDF).

\textbf{Correspondencia:} El tema se desarrolla en las secciones
indicadas. La redacción es una síntesis docente.

\textbf{Microcurrículo:} semana 11, sesión 1 (sesión global 21). Los
numerales del cronograma son identificadores curriculares; no son los
apartados del libro citados arriba.
\end{frame}

\begin{frame}{12. Preguntas que orientan la sesión}
\phantomsection\label{preguntas-que-orientan-la-sesiuxf3n}
\textbf{Libro guía:} John Semmlow, \emph{Circuits, Signals, and Systems
for Bioengineers}, 3.ª ed., 2018.

§5.4, \textbf{Using The Impulse Response To Find A Systems Output To Any
Input---Convolution} (página 219 del visor PDF).

\textbf{Correspondencia:} El tema se desarrolla en las secciones
indicadas. La redacción es una síntesis docente.

\textbf{Microcurrículo:} semana 11, sesión 1 (sesión global 21). Los
numerales del cronograma son identificadores curriculares; no son los
apartados del libro citados arriba.
\end{frame}

\begin{frame}{13. Vocabulario mínimo}
\phantomsection\label{vocabulario-muxednimo}
\textbf{Libro guía:} John Semmlow, \emph{Circuits, Signals, and Systems
for Bioengineers}, 3.ª ed., 2018.

§5.4, \textbf{Using The Impulse Response To Find A Systems Output To Any
Input---Convolution} (página 219 del visor PDF).

\textbf{Correspondencia:} El tema se desarrolla en las secciones
indicadas. La redacción es una síntesis docente.

\textbf{Microcurrículo:} semana 11, sesión 1 (sesión global 21). Los
numerales del cronograma son identificadores curriculares; no son los
apartados del libro citados arriba.
\end{frame}

\begin{frame}{14. Relaciones matemáticas centrales}
\phantomsection\label{relaciones-matemuxe1ticas-centrales}
\textbf{Libro guía:} John Semmlow, \emph{Circuits, Signals, and Systems
for Bioengineers}, 3.ª ed., 2018.

§5.4, \textbf{Using The Impulse Response To Find A Systems Output To Any
Input---Convolution} (página 219 del visor PDF).

\textbf{Correspondencia:} El tema se desarrolla en las secciones
indicadas. La redacción es una síntesis docente.

\textbf{Microcurrículo:} semana 11, sesión 1 (sesión global 21). Los
numerales del cronograma son identificadores curriculares; no son los
apartados del libro citados arriba.
\end{frame}

\begin{frame}{15. Convolución continua: un caso por intervalos}
\phantomsection\label{convoluciuxf3n-continua-un-caso-por-intervalos}
\textbf{Libro guía:} John Semmlow, \emph{Circuits, Signals, and Systems
for Bioengineers}, 3.ª ed., 2018.

§5.4, \textbf{Using The Impulse Response To Find A Systems Output To Any
Input---Convolution} (página 219 del visor PDF).

\textbf{Correspondencia:} Ejemplo, figura o implementación de
elaboración docente a partir de estos fundamentos. No es una
transcripción del código MATLAB del libro. Desarrollo didáctico original
para el cálculo por intervalos.

\textbf{Microcurrículo:} semana 11, sesión 1 (sesión global 21). Los
numerales del cronograma son identificadores curriculares; no son los
apartados del libro citados arriba.
\end{frame}

\begin{frame}{16. Procedimiento de análisis}
\phantomsection\label{procedimiento-de-anuxe1lisis}
\textbf{Libro guía:} John Semmlow, \emph{Circuits, Signals, and Systems
for Bioengineers}, 3.ª ed., 2018.

§5.4, \textbf{Using The Impulse Response To Find A Systems Output To Any
Input---Convolution} (página 219 del visor PDF).

\textbf{Correspondencia:} El tema se desarrolla en las secciones
indicadas. La redacción es una síntesis docente.

\textbf{Microcurrículo:} semana 11, sesión 1 (sesión global 21). Los
numerales del cronograma son identificadores curriculares; no son los
apartados del libro citados arriba.
\end{frame}

\begin{frame}{17. Ejemplo biomédico}
\phantomsection\label{ejemplo-biomuxe9dico}
\textbf{Libro guía:} John Semmlow, \emph{Circuits, Signals, and Systems
for Bioengineers}, 3.ª ed., 2018.

§5.4, \textbf{Using The Impulse Response To Find A Systems Output To Any
Input---Convolution} (página 219 del visor PDF); §1.4.5,
\textbf{Biosystems Modeling} (página 46 del visor PDF).

\textbf{Correspondencia:} Actividad o interpretación biomédica de
elaboración docente apoyada en las secciones indicadas.

\textbf{Microcurrículo:} semana 11, sesión 1 (sesión global 21). Los
numerales del cronograma son identificadores curriculares; no son los
apartados del libro citados arriba.
\end{frame}

\begin{frame}{18. Implementación mínima en Python}
\phantomsection\label{implementaciuxf3n-muxednima-en-python}
\textbf{Libro guía:} John Semmlow, \emph{Circuits, Signals, and Systems
for Bioengineers}, 3.ª ed., 2018.

§5.4.2, \textbf{Matlab Implementation} (página 229 del visor PDF).

\textbf{Correspondencia:} Ejemplo, figura o implementación de
elaboración docente a partir de estos fundamentos. No es una
transcripción del código MATLAB del libro.

\textbf{Microcurrículo:} semana 11, sesión 1 (sesión global 21). Los
numerales del cronograma son identificadores curriculares; no son los
apartados del libro citados arriba.
\end{frame}

\begin{frame}{19. Convolución continua aproximada con muestras}
\phantomsection\label{convoluciuxf3n-continua-aproximada-con-muestras}
\textbf{Libro guía:} John Semmlow, \emph{Circuits, Signals, and Systems
for Bioengineers}, 3.ª ed., 2018.

§5.4, \textbf{Using The Impulse Response To Find A Systems Output To Any
Input---Convolution} (página 219 del visor PDF); §5.4.2, \textbf{Matlab
Implementation} (página 229 del visor PDF).

\textbf{Correspondencia:} Ejemplo, figura o implementación de
elaboración docente a partir de estos fundamentos. No es una
transcripción del código MATLAB del libro.

\textbf{Microcurrículo:} semana 11, sesión 1 (sesión global 21). Los
numerales del cronograma son identificadores curriculares; no son los
apartados del libro citados arriba.
\end{frame}

\begin{frame}{20. Convolución: solución y aproximación}
\phantomsection\label{convoluciuxf3n-soluciuxf3n-y-aproximaciuxf3n}
\textbf{Libro guía:} John Semmlow, \emph{Circuits, Signals, and Systems
for Bioengineers}, 3.ª ed., 2018.

§5.4, \textbf{Using The Impulse Response To Find A Systems Output To Any
Input---Convolution} (página 219 del visor PDF); §5.4.2, \textbf{Matlab
Implementation} (página 229 del visor PDF).

\textbf{Correspondencia:} Ejemplo, figura o implementación de
elaboración docente a partir de estos fundamentos. No es una
transcripción del código MATLAB del libro. Figura de cálculo propio con
señales sintéticas. Código reproducible:
codigo/ejemplos\_verificados.py.

\textbf{Microcurrículo:} semana 11, sesión 1 (sesión global 21). Los
numerales del cronograma son identificadores curriculares; no son los
apartados del libro citados arriba.
\end{frame}

\begin{frame}{21. Verificaciones antes de interpretar}
\phantomsection\label{verificaciones-antes-de-interpretar}
\textbf{Libro guía:} John Semmlow, \emph{Circuits, Signals, and Systems
for Bioengineers}, 3.ª ed., 2018.

§5.4, \textbf{Using The Impulse Response To Find A Systems Output To Any
Input---Convolution} (página 219 del visor PDF).

\textbf{Correspondencia:} El tema se desarrolla en las secciones
indicadas. La redacción es una síntesis docente.

\textbf{Microcurrículo:} semana 11, sesión 1 (sesión global 21). Los
numerales del cronograma son identificadores curriculares; no son los
apartados del libro citados arriba.
\end{frame}

\begin{frame}{22. Errores frecuentes}
\phantomsection\label{errores-frecuentes}
\textbf{Libro guía:} John Semmlow, \emph{Circuits, Signals, and Systems
for Bioengineers}, 3.ª ed., 2018.

§5.4, \textbf{Using The Impulse Response To Find A Systems Output To Any
Input---Convolution} (página 219 del visor PDF).

\textbf{Correspondencia:} El tema se desarrolla en las secciones
indicadas. La redacción es una síntesis docente.

\textbf{Microcurrículo:} semana 11, sesión 1 (sesión global 21). Los
numerales del cronograma son identificadores curriculares; no son los
apartados del libro citados arriba.
\end{frame}

\begin{frame}{23. Actividad guiada}
\phantomsection\label{actividad-guiada}
\textbf{Libro guía:} John Semmlow, \emph{Circuits, Signals, and Systems
for Bioengineers}, 3.ª ed., 2018.

§5.4, \textbf{Using The Impulse Response To Find A Systems Output To Any
Input---Convolution} (página 219 del visor PDF).

\textbf{Correspondencia:} El tema se desarrolla en las secciones
indicadas. La redacción es una síntesis docente.

\textbf{Microcurrículo:} semana 11, sesión 1 (sesión global 21). Los
numerales del cronograma son identificadores curriculares; no son los
apartados del libro citados arriba.
\end{frame}

\begin{frame}{24. Pregunta de comprobación}
\phantomsection\label{pregunta-de-comprobaciuxf3n}
\textbf{Libro guía:} John Semmlow, \emph{Circuits, Signals, and Systems
for Bioengineers}, 3.ª ed., 2018.

§5.4, \textbf{Using The Impulse Response To Find A Systems Output To Any
Input---Convolution} (página 219 del visor PDF).

\textbf{Correspondencia:} El tema se desarrolla en las secciones
indicadas. La redacción es una síntesis docente.

\textbf{Microcurrículo:} semana 11, sesión 1 (sesión global 21). Los
numerales del cronograma son identificadores curriculares; no son los
apartados del libro citados arriba.
\end{frame}

\begin{frame}{25. Síntesis de la sesión}
\phantomsection\label{suxedntesis-de-la-sesiuxf3n}
\textbf{Libro guía:} John Semmlow, \emph{Circuits, Signals, and Systems
for Bioengineers}, 3.ª ed., 2018.

§5.4, \textbf{Using The Impulse Response To Find A Systems Output To Any
Input---Convolution} (página 219 del visor PDF).

\textbf{Correspondencia:} El tema se desarrolla en las secciones
indicadas. La redacción es una síntesis docente.

\textbf{Microcurrículo:} semana 11, sesión 1 (sesión global 21). Los
numerales del cronograma son identificadores curriculares; no son los
apartados del libro citados arriba.
\end{frame}

\begin{frame}{26. Sesión 22 · Convolución aplicada: suavizado,
diferenciación y promedio móvil}
\phantomsection\label{sesiuxf3n-22-convoluciuxf3n-aplicada-suavizado-diferenciaciuxf3n-y-promedio-muxf3vil}
\textbf{Libro guía:} John Semmlow, \emph{Circuits, Signals, and Systems
for Bioengineers}, 3.ª ed., 2018.

§5.5, \textbf{Applied Convolution---Basic Filters} (página 231 del visor
PDF); §8.5.1, \textbf{Derivative Filters---The Two-Point Central
Difference Algorithm} (página 377 del visor PDF).

\textbf{Correspondencia:} El tema se desarrolla en las secciones
indicadas. La redacción es una síntesis docente.

\textbf{Microcurrículo:} semana 11, sesión 2 (sesión global 22). Los
numerales del cronograma son identificadores curriculares; no son los
apartados del libro citados arriba.
\end{frame}

\begin{frame}{27. Preguntas que orientan la sesión}
\phantomsection\label{preguntas-que-orientan-la-sesiuxf3n-1}
\textbf{Libro guía:} John Semmlow, \emph{Circuits, Signals, and Systems
for Bioengineers}, 3.ª ed., 2018.

§5.5, \textbf{Applied Convolution---Basic Filters} (página 231 del visor
PDF); §8.5.1, \textbf{Derivative Filters---The Two-Point Central
Difference Algorithm} (página 377 del visor PDF).

\textbf{Correspondencia:} El tema se desarrolla en las secciones
indicadas. La redacción es una síntesis docente.

\textbf{Microcurrículo:} semana 11, sesión 2 (sesión global 22). Los
numerales del cronograma son identificadores curriculares; no son los
apartados del libro citados arriba.
\end{frame}

\begin{frame}{28. Vocabulario mínimo}
\phantomsection\label{vocabulario-muxednimo-1}
\textbf{Libro guía:} John Semmlow, \emph{Circuits, Signals, and Systems
for Bioengineers}, 3.ª ed., 2018.

§5.5, \textbf{Applied Convolution---Basic Filters} (página 231 del visor
PDF); §8.5.1, \textbf{Derivative Filters---The Two-Point Central
Difference Algorithm} (página 377 del visor PDF).

\textbf{Correspondencia:} El tema se desarrolla en las secciones
indicadas. La redacción es una síntesis docente.

\textbf{Microcurrículo:} semana 11, sesión 2 (sesión global 22). Los
numerales del cronograma son identificadores curriculares; no son los
apartados del libro citados arriba.
\end{frame}

\begin{frame}{29. Relaciones matemáticas centrales}
\phantomsection\label{relaciones-matemuxe1ticas-centrales-1}
\textbf{Libro guía:} John Semmlow, \emph{Circuits, Signals, and Systems
for Bioengineers}, 3.ª ed., 2018.

§5.5, \textbf{Applied Convolution---Basic Filters} (página 231 del visor
PDF); §8.5.1, \textbf{Derivative Filters---The Two-Point Central
Difference Algorithm} (página 377 del visor PDF).

\textbf{Correspondencia:} El tema se desarrolla en las secciones
indicadas. La redacción es una síntesis docente.

\textbf{Microcurrículo:} semana 11, sesión 2 (sesión global 22). Los
numerales del cronograma son identificadores curriculares; no son los
apartados del libro citados arriba.
\end{frame}

\begin{frame}{30. Procedimiento de análisis}
\phantomsection\label{procedimiento-de-anuxe1lisis-1}
\textbf{Libro guía:} John Semmlow, \emph{Circuits, Signals, and Systems
for Bioengineers}, 3.ª ed., 2018.

§5.5, \textbf{Applied Convolution---Basic Filters} (página 231 del visor
PDF); §8.5.1, \textbf{Derivative Filters---The Two-Point Central
Difference Algorithm} (página 377 del visor PDF).

\textbf{Correspondencia:} El tema se desarrolla en las secciones
indicadas. La redacción es una síntesis docente.

\textbf{Microcurrículo:} semana 11, sesión 2 (sesión global 22). Los
numerales del cronograma son identificadores curriculares; no son los
apartados del libro citados arriba.
\end{frame}

\begin{frame}{31. Ejemplo biomédico}
\phantomsection\label{ejemplo-biomuxe9dico-1}
\textbf{Libro guía:} John Semmlow, \emph{Circuits, Signals, and Systems
for Bioengineers}, 3.ª ed., 2018.

§5.5, \textbf{Applied Convolution---Basic Filters} (página 231 del visor
PDF).

\textbf{Correspondencia:} Actividad o interpretación biomédica de
elaboración docente apoyada en las secciones indicadas.

\textbf{Microcurrículo:} semana 11, sesión 2 (sesión global 22). Los
numerales del cronograma son identificadores curriculares; no son los
apartados del libro citados arriba.
\end{frame}

\begin{frame}{32. Implementación mínima en Python}
\phantomsection\label{implementaciuxf3n-muxednima-en-python-1}
\textbf{Libro guía:} John Semmlow, \emph{Circuits, Signals, and Systems
for Bioengineers}, 3.ª ed., 2018.

§5.5, \textbf{Applied Convolution---Basic Filters} (página 231 del visor
PDF).

\textbf{Correspondencia:} Ejemplo, figura o implementación de
elaboración docente a partir de estos fundamentos. No es una
transcripción del código MATLAB del libro.

\textbf{Microcurrículo:} semana 11, sesión 2 (sesión global 22). Los
numerales del cronograma son identificadores curriculares; no son los
apartados del libro citados arriba.
\end{frame}

\begin{frame}{33. Verificaciones antes de interpretar}
\phantomsection\label{verificaciones-antes-de-interpretar-1}
\textbf{Libro guía:} John Semmlow, \emph{Circuits, Signals, and Systems
for Bioengineers}, 3.ª ed., 2018.

§5.5, \textbf{Applied Convolution---Basic Filters} (página 231 del visor
PDF); §8.5.1, \textbf{Derivative Filters---The Two-Point Central
Difference Algorithm} (página 377 del visor PDF).

\textbf{Correspondencia:} El tema se desarrolla en las secciones
indicadas. La redacción es una síntesis docente.

\textbf{Microcurrículo:} semana 11, sesión 2 (sesión global 22). Los
numerales del cronograma son identificadores curriculares; no son los
apartados del libro citados arriba.
\end{frame}

\begin{frame}{34. Errores frecuentes}
\phantomsection\label{errores-frecuentes-1}
\textbf{Libro guía:} John Semmlow, \emph{Circuits, Signals, and Systems
for Bioengineers}, 3.ª ed., 2018.

§5.5, \textbf{Applied Convolution---Basic Filters} (página 231 del visor
PDF); §8.5.1, \textbf{Derivative Filters---The Two-Point Central
Difference Algorithm} (página 377 del visor PDF).

\textbf{Correspondencia:} El tema se desarrolla en las secciones
indicadas. La redacción es una síntesis docente.

\textbf{Microcurrículo:} semana 11, sesión 2 (sesión global 22). Los
numerales del cronograma son identificadores curriculares; no son los
apartados del libro citados arriba.
\end{frame}

\begin{frame}{35. Actividad guiada}
\phantomsection\label{actividad-guiada-1}
\textbf{Libro guía:} John Semmlow, \emph{Circuits, Signals, and Systems
for Bioengineers}, 3.ª ed., 2018.

§5.5, \textbf{Applied Convolution---Basic Filters} (página 231 del visor
PDF); §8.5.1, \textbf{Derivative Filters---The Two-Point Central
Difference Algorithm} (página 377 del visor PDF).

\textbf{Correspondencia:} El tema se desarrolla en las secciones
indicadas. La redacción es una síntesis docente.

\textbf{Microcurrículo:} semana 11, sesión 2 (sesión global 22). Los
numerales del cronograma son identificadores curriculares; no son los
apartados del libro citados arriba.
\end{frame}

\begin{frame}{36. Pregunta de comprobación}
\phantomsection\label{pregunta-de-comprobaciuxf3n-1}
\textbf{Libro guía:} John Semmlow, \emph{Circuits, Signals, and Systems
for Bioengineers}, 3.ª ed., 2018.

§5.5, \textbf{Applied Convolution---Basic Filters} (página 231 del visor
PDF); §8.5.1, \textbf{Derivative Filters---The Two-Point Central
Difference Algorithm} (página 377 del visor PDF).

\textbf{Correspondencia:} El tema se desarrolla en las secciones
indicadas. La redacción es una síntesis docente.

\textbf{Microcurrículo:} semana 11, sesión 2 (sesión global 22). Los
numerales del cronograma son identificadores curriculares; no son los
apartados del libro citados arriba.
\end{frame}

\begin{frame}{37. Síntesis de la sesión}
\phantomsection\label{suxedntesis-de-la-sesiuxf3n-1}
\textbf{Libro guía:} John Semmlow, \emph{Circuits, Signals, and Systems
for Bioengineers}, 3.ª ed., 2018.

§5.5, \textbf{Applied Convolution---Basic Filters} (página 231 del visor
PDF); §8.5.1, \textbf{Derivative Filters---The Two-Point Central
Difference Algorithm} (página 377 del visor PDF).

\textbf{Correspondencia:} El tema se desarrolla en las secciones
indicadas. La redacción es una síntesis docente.

\textbf{Microcurrículo:} semana 11, sesión 2 (sesión global 22). Los
numerales del cronograma son identificadores curriculares; no son los
apartados del libro citados arriba.
\end{frame}

\begin{frame}{38. Semana 12 · Respuesta frecuencial y transformada de
Laplace}
\phantomsection\label{semana-12-respuesta-frecuencial-y-transformada-de-laplace}
\textbf{Libro guía:} John Semmlow, \emph{Circuits, Signals, and Systems
for Bioengineers}, 3.ª ed., 2018.

§5.4, \textbf{Using The Impulse Response To Find A Systems Output To Any
Input---Convolution} (página 219 del visor PDF); §6.4, \textbf{The
Transfer Function} (página 255 del visor PDF); §7.4, \textbf{Nonzero
Initial Conditions---Initial And Final Value Theorems} (página 326 del
visor PDF); §8.4, \textbf{Linear Filters---Introduction} (página 358 del
visor PDF).

\textbf{Correspondencia:} Síntesis o encuadre que reúne varias
secciones; no corresponde a un único apartado del libro.

\textbf{Lectura:} use estas secciones como ruta de preparación y
consulta. La bibliografía aparece al final del bloque.
\end{frame}

\begin{frame}{39. Sesión 23 · Convolución en frecuencia y respuesta del
sistema}
\phantomsection\label{sesiuxf3n-23-convoluciuxf3n-en-frecuencia-y-respuesta-del-sistema}
\textbf{Libro guía:} John Semmlow, \emph{Circuits, Signals, and Systems
for Bioengineers}, 3.ª ed., 2018.

§5.6, \textbf{Convolution In The Frequency Domain} (página 237 del visor
PDF); §6.4, \textbf{The Transfer Function} (página 255 del visor PDF);
§6.7, \textbf{The Transfer Function And The Fourier Transform} (página
287 del visor PDF).

\textbf{Correspondencia:} El tema se desarrolla en las secciones
indicadas. La redacción es una síntesis docente.

\textbf{Microcurrículo:} semana 12, sesión 1 (sesión global 23). Los
numerales del cronograma son identificadores curriculares; no son los
apartados del libro citados arriba.
\end{frame}

\begin{frame}{40. Preguntas que orientan la sesión}
\phantomsection\label{preguntas-que-orientan-la-sesiuxf3n-2}
\textbf{Libro guía:} John Semmlow, \emph{Circuits, Signals, and Systems
for Bioengineers}, 3.ª ed., 2018.

§5.6, \textbf{Convolution In The Frequency Domain} (página 237 del visor
PDF); §6.4, \textbf{The Transfer Function} (página 255 del visor PDF);
§6.7, \textbf{The Transfer Function And The Fourier Transform} (página
287 del visor PDF).

\textbf{Correspondencia:} El tema se desarrolla en las secciones
indicadas. La redacción es una síntesis docente.

\textbf{Microcurrículo:} semana 12, sesión 1 (sesión global 23). Los
numerales del cronograma son identificadores curriculares; no son los
apartados del libro citados arriba.
\end{frame}

\begin{frame}{41. Vocabulario mínimo}
\phantomsection\label{vocabulario-muxednimo-2}
\textbf{Libro guía:} John Semmlow, \emph{Circuits, Signals, and Systems
for Bioengineers}, 3.ª ed., 2018.

§6.4, \textbf{The Transfer Function} (página 255 del visor PDF); §8.2.3,
\textbf{Transformation From The Z-Domain To The Frequency Domain}
(página 353 del visor PDF).

\textbf{Correspondencia:} El tema se desarrolla en las secciones
indicadas. La redacción es una síntesis docente.

\textbf{Microcurrículo:} semana 12, sesión 1 (sesión global 23). Los
numerales del cronograma son identificadores curriculares; no son los
apartados del libro citados arriba.
\end{frame}

\begin{frame}{42. Relaciones matemáticas centrales}
\phantomsection\label{relaciones-matemuxe1ticas-centrales-2}
\textbf{Libro guía:} John Semmlow, \emph{Circuits, Signals, and Systems
for Bioengineers}, 3.ª ed., 2018.

§5.6, \textbf{Convolution In The Frequency Domain} (página 237 del visor
PDF); §6.4, \textbf{The Transfer Function} (página 255 del visor PDF);
§6.7, \textbf{The Transfer Function And The Fourier Transform} (página
287 del visor PDF).

\textbf{Correspondencia:} El tema se desarrolla en las secciones
indicadas. La redacción es una síntesis docente.

\textbf{Microcurrículo:} semana 12, sesión 1 (sesión global 23). Los
numerales del cronograma son identificadores curriculares; no son los
apartados del libro citados arriba.
\end{frame}

\begin{frame}{43. Modelos de sistemas y conexiones}
\phantomsection\label{modelos-de-sistemas-y-conexiones}
\textbf{Libro guía:} John Semmlow, \emph{Circuits, Signals, and Systems
for Bioengineers}, 3.ª ed., 2018.

§6.2, \textbf{Systems Analysis Models} (página 248 del visor PDF);
§1.4.5, \textbf{Biosystems Modeling} (página 46 del visor PDF).

\textbf{Correspondencia:} El tema se desarrolla en las secciones
indicadas. La redacción es una síntesis docente.

\textbf{Microcurrículo:} semana 12, sesión 1 (sesión global 23). Los
numerales del cronograma son identificadores curriculares; no son los
apartados del libro citados arriba.
\end{frame}

\begin{frame}{44. Respuesta sinusoidal y fasores}
\phantomsection\label{respuesta-sinusoidal-y-fasores}
\textbf{Libro guía:} John Semmlow, \emph{Circuits, Signals, and Systems
for Bioengineers}, 3.ª ed., 2018.

§6.3, \textbf{The Response Of System Elements To Sinusoidal Inputs:
Phasor Analysis} (página 251 del visor PDF); §6.4, \textbf{The Transfer
Function} (página 255 del visor PDF).

\textbf{Correspondencia:} El tema se desarrolla en las secciones
indicadas. La redacción es una síntesis docente.

\textbf{Microcurrículo:} semana 12, sesión 1 (sesión global 23). Los
numerales del cronograma son identificadores curriculares; no son los
apartados del libro citados arriba.
\end{frame}

\begin{frame}{45. Bode de un sistema de primer orden}
\phantomsection\label{bode-de-un-sistema-de-primer-orden}
\textbf{Libro guía:} John Semmlow, \emph{Circuits, Signals, and Systems
for Bioengineers}, 3.ª ed., 2018.

§6.5.4, \textbf{First-Order Element} (página 266 del visor PDF).

\textbf{Correspondencia:} El tema se desarrolla en las secciones
indicadas. La redacción es una síntesis docente.

\textbf{Microcurrículo:} semana 12, sesión 1 (sesión global 23). Los
numerales del cronograma son identificadores curriculares; no son los
apartados del libro citados arriba.
\end{frame}

\begin{frame}{46. Magnitud y fase de un primer orden}
\phantomsection\label{magnitud-y-fase-de-un-primer-orden}
\textbf{Libro guía:} John Semmlow, \emph{Circuits, Signals, and Systems
for Bioengineers}, 3.ª ed., 2018.

§6.5.4, \textbf{First-Order Element} (página 266 del visor PDF).

\textbf{Correspondencia:} Ejemplo, figura o implementación de
elaboración docente a partir de estos fundamentos. No es una
transcripción del código MATLAB del libro. Figura de cálculo propio con
señales sintéticas. Código reproducible:
codigo/ejemplos\_verificados.py.

\textbf{Microcurrículo:} semana 12, sesión 1 (sesión global 23). Los
numerales del cronograma son identificadores curriculares; no son los
apartados del libro citados arriba.
\end{frame}

\begin{frame}{47. Segundo orden, amortiguamiento y resonancia}
\phantomsection\label{segundo-orden-amortiguamiento-y-resonancia}
\textbf{Libro guía:} John Semmlow, \emph{Circuits, Signals, and Systems
for Bioengineers}, 3.ª ed., 2018.

§6.5.5, \textbf{Second-Order Element} (página 271 del visor PDF);
§7.3.6, \textbf{Second-Order Element} (página 314 del visor PDF).

\textbf{Correspondencia:} El tema se desarrolla en las secciones
indicadas. La redacción es una síntesis docente.

\textbf{Microcurrículo:} semana 12, sesión 1 (sesión global 23). Los
numerales del cronograma son identificadores curriculares; no son los
apartados del libro citados arriba.
\end{frame}

\begin{frame}{48. Composición de Bode y reconocimiento de dinámica}
\phantomsection\label{composiciuxf3n-de-bode-y-reconocimiento-de-dinuxe1mica}
\textbf{Libro guía:} John Semmlow, \emph{Circuits, Signals, and Systems
for Bioengineers}, 3.ª ed., 2018.

§6.5.1, \textbf{Constant Gain Element} (página 264 del visor PDF);
§6.5.2, \textbf{Derivative Element} (página 264 del visor PDF); §6.5.3,
\textbf{Integrator Element} (página 266 del visor PDF); §6.6,
\textbf{Bode Plots Combining Multiple Elements} (página 277 del visor
PDF); §6.6.1, \textbf{Constructing The Transfer Function From The System
Spectrum} (página 286 del visor PDF).

\textbf{Correspondencia:} El tema se desarrolla en las secciones
indicadas. La redacción es una síntesis docente.

\textbf{Microcurrículo:} semana 12, sesión 1 (sesión global 23). Los
numerales del cronograma son identificadores curriculares; no son los
apartados del libro citados arriba.
\end{frame}

\begin{frame}{49. Procedimiento de análisis}
\phantomsection\label{procedimiento-de-anuxe1lisis-2}
\textbf{Libro guía:} John Semmlow, \emph{Circuits, Signals, and Systems
for Bioengineers}, 3.ª ed., 2018.

§5.6, \textbf{Convolution In The Frequency Domain} (página 237 del visor
PDF); §6.4, \textbf{The Transfer Function} (página 255 del visor PDF);
§6.7, \textbf{The Transfer Function And The Fourier Transform} (página
287 del visor PDF).

\textbf{Correspondencia:} El tema se desarrolla en las secciones
indicadas. La redacción es una síntesis docente.

\textbf{Microcurrículo:} semana 12, sesión 1 (sesión global 23). Los
numerales del cronograma son identificadores curriculares; no son los
apartados del libro citados arriba.
\end{frame}

\begin{frame}{50. Ejemplo biomédico}
\phantomsection\label{ejemplo-biomuxe9dico-2}
\textbf{Libro guía:} John Semmlow, \emph{Circuits, Signals, and Systems
for Bioengineers}, 3.ª ed., 2018.

§5.6, \textbf{Convolution In The Frequency Domain} (página 237 del visor
PDF); §6.4, \textbf{The Transfer Function} (página 255 del visor PDF);
§6.7, \textbf{The Transfer Function And The Fourier Transform} (página
287 del visor PDF).

\textbf{Correspondencia:} El tema se desarrolla en las secciones
indicadas. La redacción es una síntesis docente.

\textbf{Microcurrículo:} semana 12, sesión 1 (sesión global 23). Los
numerales del cronograma son identificadores curriculares; no son los
apartados del libro citados arriba.
\end{frame}

\begin{frame}{51. Implementación mínima en Python}
\phantomsection\label{implementaciuxf3n-muxednima-en-python-2}
\textbf{Libro guía:} John Semmlow, \emph{Circuits, Signals, and Systems
for Bioengineers}, 3.ª ed., 2018.

§8.4.1, \textbf{Filter Properties} (página 360 del visor PDF).

\textbf{Correspondencia:} Ejemplo, figura o implementación de
elaboración docente a partir de estos fundamentos. No es una
transcripción del código MATLAB del libro.

\textbf{Microcurrículo:} semana 12, sesión 1 (sesión global 23). Los
numerales del cronograma son identificadores curriculares; no son los
apartados del libro citados arriba.
\end{frame}

\begin{frame}{52. Verificaciones antes de interpretar}
\phantomsection\label{verificaciones-antes-de-interpretar-2}
\textbf{Libro guía:} John Semmlow, \emph{Circuits, Signals, and Systems
for Bioengineers}, 3.ª ed., 2018.

§5.6, \textbf{Convolution In The Frequency Domain} (página 237 del visor
PDF); §6.4, \textbf{The Transfer Function} (página 255 del visor PDF);
§6.7, \textbf{The Transfer Function And The Fourier Transform} (página
287 del visor PDF).

\textbf{Correspondencia:} El tema se desarrolla en las secciones
indicadas. La redacción es una síntesis docente.

\textbf{Microcurrículo:} semana 12, sesión 1 (sesión global 23). Los
numerales del cronograma son identificadores curriculares; no son los
apartados del libro citados arriba.
\end{frame}

\begin{frame}{53. Errores frecuentes}
\phantomsection\label{errores-frecuentes-2}
\textbf{Libro guía:} John Semmlow, \emph{Circuits, Signals, and Systems
for Bioengineers}, 3.ª ed., 2018.

§5.6, \textbf{Convolution In The Frequency Domain} (página 237 del visor
PDF); §6.4, \textbf{The Transfer Function} (página 255 del visor PDF);
§6.7, \textbf{The Transfer Function And The Fourier Transform} (página
287 del visor PDF).

\textbf{Correspondencia:} El tema se desarrolla en las secciones
indicadas. La redacción es una síntesis docente.

\textbf{Microcurrículo:} semana 12, sesión 1 (sesión global 23). Los
numerales del cronograma son identificadores curriculares; no son los
apartados del libro citados arriba.
\end{frame}

\begin{frame}{54. Actividad guiada}
\phantomsection\label{actividad-guiada-2}
\textbf{Libro guía:} John Semmlow, \emph{Circuits, Signals, and Systems
for Bioengineers}, 3.ª ed., 2018.

§5.6, \textbf{Convolution In The Frequency Domain} (página 237 del visor
PDF); §6.4, \textbf{The Transfer Function} (página 255 del visor PDF);
§6.7, \textbf{The Transfer Function And The Fourier Transform} (página
287 del visor PDF).

\textbf{Correspondencia:} El tema se desarrolla en las secciones
indicadas. La redacción es una síntesis docente.

\textbf{Microcurrículo:} semana 12, sesión 1 (sesión global 23). Los
numerales del cronograma son identificadores curriculares; no son los
apartados del libro citados arriba.
\end{frame}

\begin{frame}{55. Pregunta de comprobación}
\phantomsection\label{pregunta-de-comprobaciuxf3n-2}
\textbf{Libro guía:} John Semmlow, \emph{Circuits, Signals, and Systems
for Bioengineers}, 3.ª ed., 2018.

§5.6, \textbf{Convolution In The Frequency Domain} (página 237 del visor
PDF); §6.4, \textbf{The Transfer Function} (página 255 del visor PDF);
§6.7, \textbf{The Transfer Function And The Fourier Transform} (página
287 del visor PDF).

\textbf{Correspondencia:} El tema se desarrolla en las secciones
indicadas. La redacción es una síntesis docente.

\textbf{Microcurrículo:} semana 12, sesión 1 (sesión global 23). Los
numerales del cronograma son identificadores curriculares; no son los
apartados del libro citados arriba.
\end{frame}

\begin{frame}{56. Síntesis de la sesión}
\phantomsection\label{suxedntesis-de-la-sesiuxf3n-2}
\textbf{Libro guía:} John Semmlow, \emph{Circuits, Signals, and Systems
for Bioengineers}, 3.ª ed., 2018.

§5.6, \textbf{Convolution In The Frequency Domain} (página 237 del visor
PDF); §6.4, \textbf{The Transfer Function} (página 255 del visor PDF);
§6.7, \textbf{The Transfer Function And The Fourier Transform} (página
287 del visor PDF).

\textbf{Correspondencia:} El tema se desarrolla en las secciones
indicadas. La redacción es una síntesis docente.

\textbf{Microcurrículo:} semana 12, sesión 1 (sesión global 23). Los
numerales del cronograma son identificadores curriculares; no son los
apartados del libro citados arriba.
\end{frame}

\begin{frame}{57. Sesión 24 · Transformada de Laplace, región de
convergencia y función de transferencia}
\phantomsection\label{sesiuxf3n-24-transformada-de-laplace-regiuxf3n-de-convergencia-y-funciuxf3n-de-transferencia}
\textbf{Libro guía:} John Semmlow, \emph{Circuits, Signals, and Systems
for Bioengineers}, 3.ª ed., 2018.

§7.2, \textbf{The Laplace Transform} (página 299 del visor PDF); §7.3,
\textbf{The Laplace Domain Transfer Function} (página 304 del visor
PDF).

\textbf{Correspondencia:} El tema se desarrolla en las secciones
indicadas. La redacción es una síntesis docente.

\textbf{Microcurrículo:} semana 12, sesión 2 (sesión global 24). Los
numerales del cronograma son identificadores curriculares; no son los
apartados del libro citados arriba.
\end{frame}

\begin{frame}{58. Preguntas que orientan la sesión}
\phantomsection\label{preguntas-que-orientan-la-sesiuxf3n-3}
\textbf{Libro guía:} John Semmlow, \emph{Circuits, Signals, and Systems
for Bioengineers}, 3.ª ed., 2018.

§7.2, \textbf{The Laplace Transform} (página 299 del visor PDF); §7.3,
\textbf{The Laplace Domain Transfer Function} (página 304 del visor
PDF).

\textbf{Correspondencia:} El tema se desarrolla en las secciones
indicadas. La redacción es una síntesis docente.

\textbf{Microcurrículo:} semana 12, sesión 2 (sesión global 24). Los
numerales del cronograma son identificadores curriculares; no son los
apartados del libro citados arriba.
\end{frame}

\begin{frame}{59. Vocabulario mínimo}
\phantomsection\label{vocabulario-muxednimo-3}
\textbf{Libro guía:} John Semmlow, \emph{Circuits, Signals, and Systems
for Bioengineers}, 3.ª ed., 2018.

§7.2.1, \textbf{Definition Of The Laplace Transform} (página 299 del
visor PDF); §7.3, \textbf{The Laplace Domain Transfer Function} (página
304 del visor PDF).

\textbf{Correspondencia:} Las secciones aportan el fundamento. La
diapositiva amplía la formulación o precisa condiciones que no se
presentan con idéntica notación en el libro.

\textbf{Microcurrículo:} semana 12, sesión 2 (sesión global 24). Los
numerales del cronograma son identificadores curriculares; no son los
apartados del libro citados arriba.
\end{frame}

\begin{frame}{60. Relaciones matemáticas centrales}
\phantomsection\label{relaciones-matemuxe1ticas-centrales-3}
\textbf{Libro guía:} John Semmlow, \emph{Circuits, Signals, and Systems
for Bioengineers}, 3.ª ed., 2018.

§7.2.1, \textbf{Definition Of The Laplace Transform} (página 299 del
visor PDF); §7.3, \textbf{The Laplace Domain Transfer Function} (página
304 del visor PDF).

\textbf{Correspondencia:} Las secciones aportan el fundamento. La
diapositiva amplía la formulación o precisa condiciones que no se
presentan con idéntica notación en el libro. La definición bilateral y
la ROC amplían el tratamiento unilateral del libro. Véase la distinción
explícita añadida en esta sesión.

\textbf{Microcurrículo:} semana 12, sesión 2 (sesión global 24). Los
numerales del cronograma son identificadores curriculares; no son los
apartados del libro citados arriba.
\end{frame}

\begin{frame}{61. Laplace bilateral y unilateral}
\phantomsection\label{laplace-bilateral-y-unilateral}
\textbf{Libro guía:} John Semmlow, \emph{Circuits, Signals, and Systems
for Bioengineers}, 3.ª ed., 2018.

§7.2.1, \textbf{Definition Of The Laplace Transform} (página 299 del
visor PDF); §7.2.2, \textbf{Calculus Operations In The Laplace Domain}
(página 301 del visor PDF); §7.4.1, \textbf{Nonzero Initial Conditions}
(página 326 del visor PDF).

\textbf{Correspondencia:} Las secciones aportan el fundamento. La
diapositiva amplía la formulación o precisa condiciones que no se
presentan con idéntica notación en el libro.

\textbf{Microcurrículo:} semana 12, sesión 2 (sesión global 24). Los
numerales del cronograma son identificadores curriculares; no son los
apartados del libro citados arriba.
\end{frame}

\begin{frame}{62. La misma expresión, dos regiones de convergencia}
\phantomsection\label{la-misma-expresiuxf3n-dos-regiones-de-convergencia}
\textbf{Libro guía:} John Semmlow, \emph{Circuits, Signals, and Systems
for Bioengineers}, 3.ª ed., 2018.

§7.2.1, \textbf{Definition Of The Laplace Transform} (página 299 del
visor PDF); §7.2.4, \textbf{Converting The Laplace Transform To The
Frequency Domain} (página 303 del visor PDF); §7.5, \textbf{The Laplace
Domain, The Frequency Domain, And The Time Domain} (página 330 del visor
PDF).

\textbf{Correspondencia:} Las secciones aportan el fundamento. La
diapositiva amplía la formulación o precisa condiciones que no se
presentan con idéntica notación en el libro. La ROC bilateral es una
extensión matemática respecto al tratamiento unilateral del libro.

\textbf{Microcurrículo:} semana 12, sesión 2 (sesión global 24). Los
numerales del cronograma son identificadores curriculares; no son los
apartados del libro citados arriba.
\end{frame}

\begin{frame}{63. Procedimiento de análisis}
\phantomsection\label{procedimiento-de-anuxe1lisis-3}
\textbf{Libro guía:} John Semmlow, \emph{Circuits, Signals, and Systems
for Bioengineers}, 3.ª ed., 2018.

§7.2, \textbf{The Laplace Transform} (página 299 del visor PDF); §7.3,
\textbf{The Laplace Domain Transfer Function} (página 304 del visor
PDF).

\textbf{Correspondencia:} El tema se desarrolla en las secciones
indicadas. La redacción es una síntesis docente.

\textbf{Microcurrículo:} semana 12, sesión 2 (sesión global 24). Los
numerales del cronograma son identificadores curriculares; no son los
apartados del libro citados arriba.
\end{frame}

\begin{frame}{64. Ejemplo biomédico}
\phantomsection\label{ejemplo-biomuxe9dico-3}
\textbf{Libro guía:} John Semmlow, \emph{Circuits, Signals, and Systems
for Bioengineers}, 3.ª ed., 2018.

§7.3.5, \textbf{First-Order Element} (página 310 del visor PDF);
§12.6.1, \textbf{Electrical Elements With Zero Initial Conditions}
(página 552 del visor PDF).

\textbf{Correspondencia:} Actividad o interpretación biomédica de
elaboración docente apoyada en las secciones indicadas.

\textbf{Microcurrículo:} semana 12, sesión 2 (sesión global 24). Los
numerales del cronograma son identificadores curriculares; no son los
apartados del libro citados arriba.
\end{frame}

\begin{frame}{65. Implementación mínima en Python}
\phantomsection\label{implementaciuxf3n-muxednima-en-python-3}
\textbf{Libro guía:} John Semmlow, \emph{Circuits, Signals, and Systems
for Bioengineers}, 3.ª ed., 2018.

§7.3.5, \textbf{First-Order Element} (página 310 del visor PDF).

\textbf{Correspondencia:} Ejemplo, figura o implementación de
elaboración docente a partir de estos fundamentos. No es una
transcripción del código MATLAB del libro.

\textbf{Microcurrículo:} semana 12, sesión 2 (sesión global 24). Los
numerales del cronograma son identificadores curriculares; no son los
apartados del libro citados arriba.
\end{frame}

\begin{frame}{66. Verificaciones antes de interpretar}
\phantomsection\label{verificaciones-antes-de-interpretar-3}
\textbf{Libro guía:} John Semmlow, \emph{Circuits, Signals, and Systems
for Bioengineers}, 3.ª ed., 2018.

§7.5, \textbf{The Laplace Domain, The Frequency Domain, And The Time
Domain} (página 330 del visor PDF).

\textbf{Correspondencia:} Las secciones aportan el fundamento. La
diapositiva amplía la formulación o precisa condiciones que no se
presentan con idéntica notación en el libro.

\textbf{Microcurrículo:} semana 12, sesión 2 (sesión global 24). Los
numerales del cronograma son identificadores curriculares; no son los
apartados del libro citados arriba.
\end{frame}

\begin{frame}{67. Errores frecuentes}
\phantomsection\label{errores-frecuentes-3}
\textbf{Libro guía:} John Semmlow, \emph{Circuits, Signals, and Systems
for Bioengineers}, 3.ª ed., 2018.

§7.2, \textbf{The Laplace Transform} (página 299 del visor PDF); §7.3,
\textbf{The Laplace Domain Transfer Function} (página 304 del visor
PDF).

\textbf{Correspondencia:} El tema se desarrolla en las secciones
indicadas. La redacción es una síntesis docente.

\textbf{Microcurrículo:} semana 12, sesión 2 (sesión global 24). Los
numerales del cronograma son identificadores curriculares; no son los
apartados del libro citados arriba.
\end{frame}

\begin{frame}{68. Actividad guiada}
\phantomsection\label{actividad-guiada-3}
\textbf{Libro guía:} John Semmlow, \emph{Circuits, Signals, and Systems
for Bioengineers}, 3.ª ed., 2018.

§7.2, \textbf{The Laplace Transform} (página 299 del visor PDF); §7.3,
\textbf{The Laplace Domain Transfer Function} (página 304 del visor
PDF).

\textbf{Correspondencia:} El tema se desarrolla en las secciones
indicadas. La redacción es una síntesis docente.

\textbf{Microcurrículo:} semana 12, sesión 2 (sesión global 24). Los
numerales del cronograma son identificadores curriculares; no son los
apartados del libro citados arriba.
\end{frame}

\begin{frame}{69. Pregunta de comprobación}
\phantomsection\label{pregunta-de-comprobaciuxf3n-3}
\textbf{Libro guía:} John Semmlow, \emph{Circuits, Signals, and Systems
for Bioengineers}, 3.ª ed., 2018.

§7.2, \textbf{The Laplace Transform} (página 299 del visor PDF); §7.3,
\textbf{The Laplace Domain Transfer Function} (página 304 del visor
PDF).

\textbf{Correspondencia:} El tema se desarrolla en las secciones
indicadas. La redacción es una síntesis docente.

\textbf{Microcurrículo:} semana 12, sesión 2 (sesión global 24). Los
numerales del cronograma son identificadores curriculares; no son los
apartados del libro citados arriba.
\end{frame}

\begin{frame}{70. Síntesis de la sesión}
\phantomsection\label{suxedntesis-de-la-sesiuxf3n-3}
\textbf{Libro guía:} John Semmlow, \emph{Circuits, Signals, and Systems
for Bioengineers}, 3.ª ed., 2018.

§7.2, \textbf{The Laplace Transform} (página 299 del visor PDF); §7.3,
\textbf{The Laplace Domain Transfer Function} (página 304 del visor
PDF).

\textbf{Correspondencia:} El tema se desarrolla en las secciones
indicadas. La redacción es una síntesis docente.

\textbf{Microcurrículo:} semana 12, sesión 2 (sesión global 24). Los
numerales del cronograma son identificadores curriculares; no son los
apartados del libro citados arriba.
\end{frame}

\begin{frame}{71. Semana 13 · Dinámica, transformada Z y diseño FIR}
\phantomsection\label{semana-13-dinuxe1mica-transformada-z-y-diseuxf1o-fir}
\textbf{Libro guía:} John Semmlow, \emph{Circuits, Signals, and Systems
for Bioengineers}, 3.ª ed., 2018.

§5.4, \textbf{Using The Impulse Response To Find A Systems Output To Any
Input---Convolution} (página 219 del visor PDF); §6.4, \textbf{The
Transfer Function} (página 255 del visor PDF); §7.4, \textbf{Nonzero
Initial Conditions---Initial And Final Value Theorems} (página 326 del
visor PDF); §8.4, \textbf{Linear Filters---Introduction} (página 358 del
visor PDF).

\textbf{Correspondencia:} Síntesis o encuadre que reúne varias
secciones; no corresponde a un único apartado del libro.

\textbf{Lectura:} use estas secciones como ruta de preparación y
consulta. La bibliografía aparece al final del bloque.
\end{frame}

\begin{frame}{72. Sesión 25 · Condiciones iniciales, valores límite e
identificación de sistemas}
\phantomsection\label{sesiuxf3n-25-condiciones-iniciales-valores-luxedmite-e-identificaciuxf3n-de-sistemas}
\textbf{Libro guía:} John Semmlow, \emph{Circuits, Signals, and Systems
for Bioengineers}, 3.ª ed., 2018.

§7.4.1, \textbf{Nonzero Initial Conditions} (página 326 del visor PDF);
§7.4.2, \textbf{Initial And Final Value Theorems} (página 328 del visor
PDF).

\textbf{Correspondencia:} El tema se desarrolla en las secciones
indicadas. La redacción es una síntesis docente.

\textbf{Microcurrículo:} semana 13, sesión 1 (sesión global 25). Los
numerales del cronograma son identificadores curriculares; no son los
apartados del libro citados arriba.
\end{frame}

\begin{frame}{73. Preguntas que orientan la sesión}
\phantomsection\label{preguntas-que-orientan-la-sesiuxf3n-4}
\textbf{Libro guía:} John Semmlow, \emph{Circuits, Signals, and Systems
for Bioengineers}, 3.ª ed., 2018.

§7.4.1, \textbf{Nonzero Initial Conditions} (página 326 del visor PDF);
§7.4.2, \textbf{Initial And Final Value Theorems} (página 328 del visor
PDF).

\textbf{Correspondencia:} El tema se desarrolla en las secciones
indicadas. La redacción es una síntesis docente.

\textbf{Microcurrículo:} semana 13, sesión 1 (sesión global 25). Los
numerales del cronograma son identificadores curriculares; no son los
apartados del libro citados arriba.
\end{frame}

\begin{frame}{74. Vocabulario mínimo}
\phantomsection\label{vocabulario-muxednimo-4}
\textbf{Libro guía:} John Semmlow, \emph{Circuits, Signals, and Systems
for Bioengineers}, 3.ª ed., 2018.

§7.4.1, \textbf{Nonzero Initial Conditions} (página 326 del visor PDF);
§7.4.2, \textbf{Initial And Final Value Theorems} (página 328 del visor
PDF).

\textbf{Correspondencia:} El tema se desarrolla en las secciones
indicadas. La redacción es una síntesis docente.

\textbf{Microcurrículo:} semana 13, sesión 1 (sesión global 25). Los
numerales del cronograma son identificadores curriculares; no son los
apartados del libro citados arriba.
\end{frame}

\begin{frame}{75. Relaciones matemáticas centrales}
\phantomsection\label{relaciones-matemuxe1ticas-centrales-4}
\textbf{Libro guía:} John Semmlow, \emph{Circuits, Signals, and Systems
for Bioengineers}, 3.ª ed., 2018.

§7.4.1, \textbf{Nonzero Initial Conditions} (página 326 del visor PDF);
§7.4.2, \textbf{Initial And Final Value Theorems} (página 328 del visor
PDF).

\textbf{Correspondencia:} El tema se desarrolla en las secciones
indicadas. La redacción es una síntesis docente.

\textbf{Microcurrículo:} semana 13, sesión 1 (sesión global 25). Los
numerales del cronograma son identificadores curriculares; no son los
apartados del libro citados arriba.
\end{frame}

\begin{frame}{76. Teorema del valor final y contraejemplo}
\phantomsection\label{teorema-del-valor-final-y-contraejemplo}
\textbf{Libro guía:} John Semmlow, \emph{Circuits, Signals, and Systems
for Bioengineers}, 3.ª ed., 2018.

§7.4.2, \textbf{Initial And Final Value Theorems} (página 328 del visor
PDF).

\textbf{Correspondencia:} El tema se desarrolla en las secciones
indicadas. La redacción es una síntesis docente.

\textbf{Microcurrículo:} semana 13, sesión 1 (sesión global 25). Los
numerales del cronograma son identificadores curriculares; no son los
apartados del libro citados arriba.
\end{frame}

\begin{frame}{77. Respuesta con estado inicial: desarrollo completo}
\phantomsection\label{respuesta-con-estado-inicial-desarrollo-completo}
\textbf{Libro guía:} John Semmlow, \emph{Circuits, Signals, and Systems
for Bioengineers}, 3.ª ed., 2018.

§7.3.5, \textbf{First-Order Element} (página 310 del visor PDF); §7.4.1,
\textbf{Nonzero Initial Conditions} (página 326 del visor PDF); §7.4.2,
\textbf{Initial And Final Value Theorems} (página 328 del visor PDF).

\textbf{Correspondencia:} Ejemplo, figura o implementación de
elaboración docente a partir de estos fundamentos. No es una
transcripción del código MATLAB del libro.

\textbf{Microcurrículo:} semana 13, sesión 1 (sesión global 25). Los
numerales del cronograma son identificadores curriculares; no son los
apartados del libro citados arriba.
\end{frame}

\begin{frame}{78. Procedimiento de análisis}
\phantomsection\label{procedimiento-de-anuxe1lisis-4}
\textbf{Libro guía:} John Semmlow, \emph{Circuits, Signals, and Systems
for Bioengineers}, 3.ª ed., 2018.

§7.6, \textbf{System Identification} (página 333 del visor PDF).

\textbf{Correspondencia:} El tema se desarrolla en las secciones
indicadas. La redacción es una síntesis docente.

\textbf{Microcurrículo:} semana 13, sesión 1 (sesión global 25). Los
numerales del cronograma son identificadores curriculares; no son los
apartados del libro citados arriba.
\end{frame}

\begin{frame}{79. Ejemplo biomédico}
\phantomsection\label{ejemplo-biomuxe9dico-4}
\textbf{Libro guía:} John Semmlow, \emph{Circuits, Signals, and Systems
for Bioengineers}, 3.ª ed., 2018.

§7.6, \textbf{System Identification} (página 333 del visor PDF).

\textbf{Correspondencia:} El tema se desarrolla en las secciones
indicadas. La redacción es una síntesis docente.

\textbf{Microcurrículo:} semana 13, sesión 1 (sesión global 25). Los
numerales del cronograma son identificadores curriculares; no son los
apartados del libro citados arriba.
\end{frame}

\begin{frame}{80. Implementación mínima en Python}
\phantomsection\label{implementaciuxf3n-muxednima-en-python-4}
\textbf{Libro guía:} John Semmlow, \emph{Circuits, Signals, and Systems
for Bioengineers}, 3.ª ed., 2018.

§7.6, \textbf{System Identification} (página 333 del visor PDF).

\textbf{Correspondencia:} Ejemplo, figura o implementación de
elaboración docente a partir de estos fundamentos. No es una
transcripción del código MATLAB del libro.

\textbf{Microcurrículo:} semana 13, sesión 1 (sesión global 25). Los
numerales del cronograma son identificadores curriculares; no son los
apartados del libro citados arriba.
\end{frame}

\begin{frame}{81. Verificaciones antes de interpretar}
\phantomsection\label{verificaciones-antes-de-interpretar-4}
\textbf{Libro guía:} John Semmlow, \emph{Circuits, Signals, and Systems
for Bioengineers}, 3.ª ed., 2018.

§7.6, \textbf{System Identification} (página 333 del visor PDF).

\textbf{Correspondencia:} El tema se desarrolla en las secciones
indicadas. La redacción es una síntesis docente.

\textbf{Microcurrículo:} semana 13, sesión 1 (sesión global 25). Los
numerales del cronograma son identificadores curriculares; no son los
apartados del libro citados arriba.
\end{frame}

\begin{frame}{82. Errores frecuentes}
\phantomsection\label{errores-frecuentes-4}
\textbf{Libro guía:} John Semmlow, \emph{Circuits, Signals, and Systems
for Bioengineers}, 3.ª ed., 2018.

§7.6, \textbf{System Identification} (página 333 del visor PDF).

\textbf{Correspondencia:} El tema se desarrolla en las secciones
indicadas. La redacción es una síntesis docente.

\textbf{Microcurrículo:} semana 13, sesión 1 (sesión global 25). Los
numerales del cronograma son identificadores curriculares; no son los
apartados del libro citados arriba.
\end{frame}

\begin{frame}{83. Actividad guiada}
\phantomsection\label{actividad-guiada-4}
\textbf{Libro guía:} John Semmlow, \emph{Circuits, Signals, and Systems
for Bioengineers}, 3.ª ed., 2018.

§7.4, \textbf{Nonzero Initial Conditions---Initial And Final Value
Theorems} (página 326 del visor PDF); §7.6, \textbf{System
Identification} (página 333 del visor PDF).

\textbf{Correspondencia:} Actividad o interpretación biomédica de
elaboración docente apoyada en las secciones indicadas.

\textbf{Microcurrículo:} semana 13, sesión 1 (sesión global 25). Los
numerales del cronograma son identificadores curriculares; no son los
apartados del libro citados arriba.
\end{frame}

\begin{frame}{84. Pregunta de comprobación}
\phantomsection\label{pregunta-de-comprobaciuxf3n-4}
\textbf{Libro guía:} John Semmlow, \emph{Circuits, Signals, and Systems
for Bioengineers}, 3.ª ed., 2018.

§7.4.2, \textbf{Initial And Final Value Theorems} (página 328 del visor
PDF); §7.6, \textbf{System Identification} (página 333 del visor PDF).

\textbf{Correspondencia:} El tema se desarrolla en las secciones
indicadas. La redacción es una síntesis docente.

\textbf{Microcurrículo:} semana 13, sesión 1 (sesión global 25). Los
numerales del cronograma son identificadores curriculares; no son los
apartados del libro citados arriba.
\end{frame}

\begin{frame}{85. Síntesis de la sesión}
\phantomsection\label{suxedntesis-de-la-sesiuxf3n-4}
\textbf{Libro guía:} John Semmlow, \emph{Circuits, Signals, and Systems
for Bioengineers}, 3.ª ed., 2018.

§7.4.2, \textbf{Initial And Final Value Theorems} (página 328 del visor
PDF); §7.6, \textbf{System Identification} (página 333 del visor PDF).

\textbf{Correspondencia:} El tema se desarrolla en las secciones
indicadas. La redacción es una síntesis docente.

\textbf{Microcurrículo:} semana 13, sesión 1 (sesión global 25). Los
numerales del cronograma son identificadores curriculares; no son los
apartados del libro citados arriba.
\end{frame}

\begin{frame}{86. Sesión 26 · Transformada Z, ecuaciones en diferencias
y filtros FIR}
\phantomsection\label{sesiuxf3n-26-transformada-z-ecuaciones-en-diferencias-y-filtros-fir}
\textbf{Libro guía:} John Semmlow, \emph{Circuits, Signals, and Systems
for Bioengineers}, 3.ª ed., 2018.

§8.2, \textbf{The Z-Transform} (página 349 del visor PDF); §8.3,
\textbf{Difference Equations} (página 355 del visor PDF); §8.5,
\textbf{Design Of Finite Impulse Response Filters} (página 366 del visor
PDF).

\textbf{Correspondencia:} El tema se desarrolla en las secciones
indicadas. La redacción es una síntesis docente.

\textbf{Microcurrículo:} semana 13, sesión 2 (sesión global 26). Los
numerales del cronograma son identificadores curriculares; no son los
apartados del libro citados arriba.
\end{frame}

\begin{frame}{87. Preguntas que orientan la sesión}
\phantomsection\label{preguntas-que-orientan-la-sesiuxf3n-5}
\textbf{Libro guía:} John Semmlow, \emph{Circuits, Signals, and Systems
for Bioengineers}, 3.ª ed., 2018.

§8.2, \textbf{The Z-Transform} (página 349 del visor PDF); §8.3,
\textbf{Difference Equations} (página 355 del visor PDF); §8.5,
\textbf{Design Of Finite Impulse Response Filters} (página 366 del visor
PDF).

\textbf{Correspondencia:} El tema se desarrolla en las secciones
indicadas. La redacción es una síntesis docente.

\textbf{Microcurrículo:} semana 13, sesión 2 (sesión global 26). Los
numerales del cronograma son identificadores curriculares; no son los
apartados del libro citados arriba.
\end{frame}

\begin{frame}{88. Vocabulario mínimo}
\phantomsection\label{vocabulario-muxednimo-5}
\textbf{Libro guía:} John Semmlow, \emph{Circuits, Signals, and Systems
for Bioengineers}, 3.ª ed., 2018.

§8.2, \textbf{The Z-Transform} (página 349 del visor PDF); §8.4.2,
\textbf{Finite Impulse Response Versus Infinite Impulse Response Filter
Characteristics} (página 363 del visor PDF).

\textbf{Correspondencia:} El tema se desarrolla en las secciones
indicadas. La redacción es una síntesis docente.

\textbf{Microcurrículo:} semana 13, sesión 2 (sesión global 26). Los
numerales del cronograma son identificadores curriculares; no son los
apartados del libro citados arriba.
\end{frame}

\begin{frame}{89. Relaciones matemáticas centrales}
\phantomsection\label{relaciones-matemuxe1ticas-centrales-5}
\textbf{Libro guía:} John Semmlow, \emph{Circuits, Signals, and Systems
for Bioengineers}, 3.ª ed., 2018.

§8.2.2, \textbf{The Digital Transfer Function} (página 351 del visor
PDF); §8.2.3, \textbf{Transformation From The Z-Domain To The Frequency
Domain} (página 353 del visor PDF); §8.3, \textbf{Difference Equations}
(página 355 del visor PDF).

\textbf{Correspondencia:} Las secciones aportan el fundamento. La
diapositiva amplía la formulación o precisa condiciones que no se
presentan con idéntica notación en el libro.

\textbf{Microcurrículo:} semana 13, sesión 2 (sesión global 26). Los
numerales del cronograma son identificadores curriculares; no son los
apartados del libro citados arriba.
\end{frame}

\begin{frame}{90. Ecuación en diferencias y transferencia}
\phantomsection\label{ecuaciuxf3n-en-diferencias-y-transferencia}
\textbf{Libro guía:} John Semmlow, \emph{Circuits, Signals, and Systems
for Bioengineers}, 3.ª ed., 2018.

§8.2.2, \textbf{The Digital Transfer Function} (página 351 del visor
PDF); §8.2.3, \textbf{Transformation From The Z-Domain To The Frequency
Domain} (página 353 del visor PDF); §8.3, \textbf{Difference Equations}
(página 355 del visor PDF).

\textbf{Correspondencia:} Ejemplo, figura o implementación de
elaboración docente a partir de estos fundamentos. No es una
transcripción del código MATLAB del libro.

\textbf{Microcurrículo:} semana 13, sesión 2 (sesión global 26). Los
numerales del cronograma son identificadores curriculares; no son los
apartados del libro citados arriba.
\end{frame}

\begin{frame}{91. Procedimiento de análisis}
\phantomsection\label{procedimiento-de-anuxe1lisis-5}
\textbf{Libro guía:} John Semmlow, \emph{Circuits, Signals, and Systems
for Bioengineers}, 3.ª ed., 2018.

§8.5, \textbf{Design Of Finite Impulse Response Filters} (página 366 del
visor PDF); §8.6.1, \textbf{Finite Impulse Response Filter Design}
(página 384 del visor PDF).

\textbf{Correspondencia:} El tema se desarrolla en las secciones
indicadas. La redacción es una síntesis docente.

\textbf{Microcurrículo:} semana 13, sesión 2 (sesión global 26). Los
numerales del cronograma son identificadores curriculares; no son los
apartados del libro citados arriba.
\end{frame}

\begin{frame}{92. Ejemplo biomédico}
\phantomsection\label{ejemplo-biomuxe9dico-5}
\textbf{Libro guía:} John Semmlow, \emph{Circuits, Signals, and Systems
for Bioengineers}, 3.ª ed., 2018.

§8.5, \textbf{Design Of Finite Impulse Response Filters} (página 366 del
visor PDF); §8.6.1, \textbf{Finite Impulse Response Filter Design}
(página 384 del visor PDF).

\textbf{Correspondencia:} El tema se desarrolla en las secciones
indicadas. La redacción es una síntesis docente.

\textbf{Microcurrículo:} semana 13, sesión 2 (sesión global 26). Los
numerales del cronograma son identificadores curriculares; no son los
apartados del libro citados arriba.
\end{frame}

\begin{frame}{93. Implementación mínima en Python}
\phantomsection\label{implementaciuxf3n-muxednima-en-python-5}
\textbf{Libro guía:} John Semmlow, \emph{Circuits, Signals, and Systems
for Bioengineers}, 3.ª ed., 2018.

§8.6.1, \textbf{Finite Impulse Response Filter Design} (página 384 del
visor PDF).

\textbf{Correspondencia:} Ejemplo, figura o implementación de
elaboración docente a partir de estos fundamentos. No es una
transcripción del código MATLAB del libro.

\textbf{Microcurrículo:} semana 13, sesión 2 (sesión global 26). Los
numerales del cronograma son identificadores curriculares; no son los
apartados del libro citados arriba.

\textbf{Implementación:}
\href{https://docs.scipy.org/doc/scipy/reference/generated/scipy.signal.firwin.html}{SciPy:
firwin y frecuencia de media amplitud}.
\end{frame}

\begin{frame}{94. Especificaciones de un FIR: más que un corte}
\phantomsection\label{especificaciones-de-un-fir-muxe1s-que-un-corte}
\textbf{Libro guía:} John Semmlow, \emph{Circuits, Signals, and Systems
for Bioengineers}, 3.ª ed., 2018.

§8.5, \textbf{Design Of Finite Impulse Response Filters} (página 366 del
visor PDF); §8.6.1, \textbf{Finite Impulse Response Filter Design}
(página 384 del visor PDF).

\textbf{Correspondencia:} Ejemplo, figura o implementación de
elaboración docente a partir de estos fundamentos. No es una
transcripción del código MATLAB del libro.

\textbf{Microcurrículo:} semana 13, sesión 2 (sesión global 26). Los
numerales del cronograma son identificadores curriculares; no son los
apartados del libro citados arriba.

\textbf{Implementación:}
\href{https://docs.scipy.org/doc/scipy/reference/generated/scipy.signal.firwin.html}{SciPy:
firwin y frecuencia de media amplitud}.
\end{frame}

\begin{frame}{95. Verificaciones antes de interpretar}
\phantomsection\label{verificaciones-antes-de-interpretar-5}
\textbf{Libro guía:} John Semmlow, \emph{Circuits, Signals, and Systems
for Bioengineers}, 3.ª ed., 2018.

§8.4.1, \textbf{Filter Properties} (página 360 del visor PDF); §8.5,
\textbf{Design Of Finite Impulse Response Filters} (página 366 del visor
PDF).

\textbf{Correspondencia:} Actividad o interpretación biomédica de
elaboración docente apoyada en las secciones indicadas.

\textbf{Microcurrículo:} semana 13, sesión 2 (sesión global 26). Los
numerales del cronograma son identificadores curriculares; no son los
apartados del libro citados arriba.
\end{frame}

\begin{frame}{96. Errores frecuentes}
\phantomsection\label{errores-frecuentes-5}
\textbf{Libro guía:} John Semmlow, \emph{Circuits, Signals, and Systems
for Bioengineers}, 3.ª ed., 2018.

§8.4.1, \textbf{Filter Properties} (página 360 del visor PDF); §8.5,
\textbf{Design Of Finite Impulse Response Filters} (página 366 del visor
PDF).

\textbf{Correspondencia:} Actividad o interpretación biomédica de
elaboración docente apoyada en las secciones indicadas.

\textbf{Microcurrículo:} semana 13, sesión 2 (sesión global 26). Los
numerales del cronograma son identificadores curriculares; no son los
apartados del libro citados arriba.
\end{frame}

\begin{frame}{97. Actividad guiada}
\phantomsection\label{actividad-guiada-5}
\textbf{Libro guía:} John Semmlow, \emph{Circuits, Signals, and Systems
for Bioengineers}, 3.ª ed., 2018.

§8.4.1, \textbf{Filter Properties} (página 360 del visor PDF); §8.5,
\textbf{Design Of Finite Impulse Response Filters} (página 366 del visor
PDF).

\textbf{Correspondencia:} Actividad o interpretación biomédica de
elaboración docente apoyada en las secciones indicadas.

\textbf{Microcurrículo:} semana 13, sesión 2 (sesión global 26). Los
numerales del cronograma son identificadores curriculares; no son los
apartados del libro citados arriba.
\end{frame}

\begin{frame}{98. Pregunta de comprobación}
\phantomsection\label{pregunta-de-comprobaciuxf3n-5}
\textbf{Libro guía:} John Semmlow, \emph{Circuits, Signals, and Systems
for Bioengineers}, 3.ª ed., 2018.

§8.4.1, \textbf{Filter Properties} (página 360 del visor PDF); §8.5,
\textbf{Design Of Finite Impulse Response Filters} (página 366 del visor
PDF).

\textbf{Correspondencia:} Actividad o interpretación biomédica de
elaboración docente apoyada en las secciones indicadas.

\textbf{Microcurrículo:} semana 13, sesión 2 (sesión global 26). Los
numerales del cronograma son identificadores curriculares; no son los
apartados del libro citados arriba.
\end{frame}

\begin{frame}{99. Síntesis de la sesión}
\phantomsection\label{suxedntesis-de-la-sesiuxf3n-5}
\textbf{Libro guía:} John Semmlow, \emph{Circuits, Signals, and Systems
for Bioengineers}, 3.ª ed., 2018.

§8.2, \textbf{The Z-Transform} (página 349 del visor PDF); §8.5,
\textbf{Design Of Finite Impulse Response Filters} (página 366 del visor
PDF).

\textbf{Correspondencia:} Síntesis o encuadre que reúne varias
secciones; no corresponde a un único apartado del libro.

\textbf{Microcurrículo:} semana 13, sesión 2 (sesión global 26). Los
numerales del cronograma son identificadores curriculares; no son los
apartados del libro citados arriba.
\end{frame}

\begin{frame}{100. Semana 14 · Filtros IIR y fundamentos de wavelets}
\phantomsection\label{semana-14-filtros-iir-y-fundamentos-de-wavelets}
\textbf{Libro guía:} John Semmlow, \emph{Circuits, Signals, and Systems
for Bioengineers}, 3.ª ed., 2018.

§5.4, \textbf{Using The Impulse Response To Find A Systems Output To Any
Input---Convolution} (página 219 del visor PDF); §6.4, \textbf{The
Transfer Function} (página 255 del visor PDF); §7.4, \textbf{Nonzero
Initial Conditions---Initial And Final Value Theorems} (página 326 del
visor PDF); §8.4, \textbf{Linear Filters---Introduction} (página 358 del
visor PDF).

\textbf{Correspondencia:} Síntesis o encuadre que reúne varias
secciones; no corresponde a un único apartado del libro.

\textbf{Lectura:} use estas secciones como ruta de preparación y
consulta. La bibliografía aparece al final del bloque.
\end{frame}

\begin{frame}{101. Sesión 27 · Filtros IIR: diseño, estabilidad, fase y
aplicaciones}
\phantomsection\label{sesiuxf3n-27-filtros-iir-diseuxf1o-estabilidad-fase-y-aplicaciones}
\textbf{Libro guía:} John Semmlow, \emph{Circuits, Signals, and Systems
for Bioengineers}, 3.ª ed., 2018.

§8.4.2, \textbf{Finite Impulse Response Versus Infinite Impulse Response
Filter Characteristics} (página 363 del visor PDF); §8.4.3,
\textbf{Causal And Noncausal Filters} (página 365 del visor PDF);
§8.6.2, \textbf{Designing Infinite Impulse Response Filters} (página 387
del visor PDF).

\textbf{Correspondencia:} El tema se desarrolla en las secciones
indicadas. La redacción es una síntesis docente.

\textbf{Microcurrículo:} semana 14, sesión 1 (sesión global 27). Los
numerales del cronograma son identificadores curriculares; no son los
apartados del libro citados arriba.
\end{frame}

\begin{frame}{102. Preguntas que orientan la sesión}
\phantomsection\label{preguntas-que-orientan-la-sesiuxf3n-6}
\textbf{Libro guía:} John Semmlow, \emph{Circuits, Signals, and Systems
for Bioengineers}, 3.ª ed., 2018.

§8.4.2, \textbf{Finite Impulse Response Versus Infinite Impulse Response
Filter Characteristics} (página 363 del visor PDF); §8.4.3,
\textbf{Causal And Noncausal Filters} (página 365 del visor PDF);
§8.6.2, \textbf{Designing Infinite Impulse Response Filters} (página 387
del visor PDF).

\textbf{Correspondencia:} El tema se desarrolla en las secciones
indicadas. La redacción es una síntesis docente.

\textbf{Microcurrículo:} semana 14, sesión 1 (sesión global 27). Los
numerales del cronograma son identificadores curriculares; no son los
apartados del libro citados arriba.
\end{frame}

\begin{frame}{103. Vocabulario mínimo}
\phantomsection\label{vocabulario-muxednimo-6}
\textbf{Libro guía:} John Semmlow, \emph{Circuits, Signals, and Systems
for Bioengineers}, 3.ª ed., 2018.

§8.4.2, \textbf{Finite Impulse Response Versus Infinite Impulse Response
Filter Characteristics} (página 363 del visor PDF); §8.4.3,
\textbf{Causal And Noncausal Filters} (página 365 del visor PDF);
§8.6.2, \textbf{Designing Infinite Impulse Response Filters} (página 387
del visor PDF).

\textbf{Correspondencia:} El tema se desarrolla en las secciones
indicadas. La redacción es una síntesis docente.

\textbf{Microcurrículo:} semana 14, sesión 1 (sesión global 27). Los
numerales del cronograma son identificadores curriculares; no son los
apartados del libro citados arriba.
\end{frame}

\begin{frame}{104. Relaciones matemáticas centrales}
\phantomsection\label{relaciones-matemuxe1ticas-centrales-6}
\textbf{Libro guía:} John Semmlow, \emph{Circuits, Signals, and Systems
for Bioengineers}, 3.ª ed., 2018.

§8.3, \textbf{Difference Equations} (página 355 del visor PDF); §8.4.3,
\textbf{Causal And Noncausal Filters} (página 365 del visor PDF).

\textbf{Correspondencia:} Las secciones aportan el fundamento. La
diapositiva amplía la formulación o precisa condiciones que no se
presentan con idéntica notación en el libro.

\textbf{Microcurrículo:} semana 14, sesión 1 (sesión global 27). Los
numerales del cronograma son identificadores curriculares; no son los
apartados del libro citados arriba.

\textbf{Implementación:}
\href{https://docs.scipy.org/doc/scipy/reference/generated/scipy.signal.filtfilt.html}{SciPy:
filtfilt}.
\end{frame}

\begin{frame}{105. Procedimiento de análisis}
\phantomsection\label{procedimiento-de-anuxe1lisis-6}
\textbf{Libro guía:} John Semmlow, \emph{Circuits, Signals, and Systems
for Bioengineers}, 3.ª ed., 2018.

§8.4.2, \textbf{Finite Impulse Response Versus Infinite Impulse Response
Filter Characteristics} (página 363 del visor PDF); §8.4.3,
\textbf{Causal And Noncausal Filters} (página 365 del visor PDF);
§8.6.2, \textbf{Designing Infinite Impulse Response Filters} (página 387
del visor PDF).

\textbf{Correspondencia:} El tema se desarrolla en las secciones
indicadas. La redacción es una síntesis docente.

\textbf{Microcurrículo:} semana 14, sesión 1 (sesión global 27). Los
numerales del cronograma son identificadores curriculares; no son los
apartados del libro citados arriba.
\end{frame}

\begin{frame}{106. Ejemplo biomédico}
\phantomsection\label{ejemplo-biomuxe9dico-6}
\textbf{Libro guía:} John Semmlow, \emph{Circuits, Signals, and Systems
for Bioengineers}, 3.ª ed., 2018.

§8.4.2, \textbf{Finite Impulse Response Versus Infinite Impulse Response
Filter Characteristics} (página 363 del visor PDF); §8.4.3,
\textbf{Causal And Noncausal Filters} (página 365 del visor PDF);
§8.6.2, \textbf{Designing Infinite Impulse Response Filters} (página 387
del visor PDF).

\textbf{Correspondencia:} El tema se desarrolla en las secciones
indicadas. La redacción es una síntesis docente.

\textbf{Microcurrículo:} semana 14, sesión 1 (sesión global 27). Los
numerales del cronograma son identificadores curriculares; no son los
apartados del libro citados arriba.
\end{frame}

\begin{frame}{107. Implementación mínima en Python}
\phantomsection\label{implementaciuxf3n-muxednima-en-python-6}
\textbf{Libro guía:} John Semmlow, \emph{Circuits, Signals, and Systems
for Bioengineers}, 3.ª ed., 2018.

§8.4.3, \textbf{Causal And Noncausal Filters} (página 365 del visor
PDF); §8.6.2, \textbf{Designing Infinite Impulse Response Filters}
(página 387 del visor PDF).

\textbf{Correspondencia:} Ejemplo, figura o implementación de
elaboración docente a partir de estos fundamentos. No es una
transcripción del código MATLAB del libro.

\textbf{Microcurrículo:} semana 14, sesión 1 (sesión global 27). Los
numerales del cronograma son identificadores curriculares; no son los
apartados del libro citados arriba.

\textbf{Implementación:}
\href{https://docs.scipy.org/doc/scipy/reference/generated/scipy.signal.filtfilt.html}{SciPy:
filtfilt}.
\end{frame}

\begin{frame}{108. Verificaciones antes de interpretar}
\phantomsection\label{verificaciones-antes-de-interpretar-6}
\textbf{Libro guía:} John Semmlow, \emph{Circuits, Signals, and Systems
for Bioengineers}, 3.ª ed., 2018.

§8.4.1, \textbf{Filter Properties} (página 360 del visor PDF); §8.6.2,
\textbf{Designing Infinite Impulse Response Filters} (página 387 del
visor PDF).

\textbf{Correspondencia:} Actividad o interpretación biomédica de
elaboración docente apoyada en las secciones indicadas.

\textbf{Microcurrículo:} semana 14, sesión 1 (sesión global 27). Los
numerales del cronograma son identificadores curriculares; no son los
apartados del libro citados arriba.
\end{frame}

\begin{frame}{109. Errores frecuentes}
\phantomsection\label{errores-frecuentes-6}
\textbf{Libro guía:} John Semmlow, \emph{Circuits, Signals, and Systems
for Bioengineers}, 3.ª ed., 2018.

§8.4.2, \textbf{Finite Impulse Response Versus Infinite Impulse Response
Filter Characteristics} (página 363 del visor PDF); §8.4.3,
\textbf{Causal And Noncausal Filters} (página 365 del visor PDF);
§8.6.2, \textbf{Designing Infinite Impulse Response Filters} (página 387
del visor PDF).

\textbf{Correspondencia:} El tema se desarrolla en las secciones
indicadas. La redacción es una síntesis docente.

\textbf{Microcurrículo:} semana 14, sesión 1 (sesión global 27). Los
numerales del cronograma son identificadores curriculares; no son los
apartados del libro citados arriba.
\end{frame}

\begin{frame}{110. Actividad guiada}
\phantomsection\label{actividad-guiada-6}
\textbf{Libro guía:} John Semmlow, \emph{Circuits, Signals, and Systems
for Bioengineers}, 3.ª ed., 2018.

§8.4.2, \textbf{Finite Impulse Response Versus Infinite Impulse Response
Filter Characteristics} (página 363 del visor PDF); §8.4.3,
\textbf{Causal And Noncausal Filters} (página 365 del visor PDF);
§8.6.2, \textbf{Designing Infinite Impulse Response Filters} (página 387
del visor PDF).

\textbf{Correspondencia:} El tema se desarrolla en las secciones
indicadas. La redacción es una síntesis docente.

\textbf{Microcurrículo:} semana 14, sesión 1 (sesión global 27). Los
numerales del cronograma son identificadores curriculares; no son los
apartados del libro citados arriba.
\end{frame}

\begin{frame}{111. Pregunta de comprobación}
\phantomsection\label{pregunta-de-comprobaciuxf3n-6}
\textbf{Libro guía:} John Semmlow, \emph{Circuits, Signals, and Systems
for Bioengineers}, 3.ª ed., 2018.

§8.4.2, \textbf{Finite Impulse Response Versus Infinite Impulse Response
Filter Characteristics} (página 363 del visor PDF); §8.4.3,
\textbf{Causal And Noncausal Filters} (página 365 del visor PDF);
§8.6.2, \textbf{Designing Infinite Impulse Response Filters} (página 387
del visor PDF).

\textbf{Correspondencia:} El tema se desarrolla en las secciones
indicadas. La redacción es una síntesis docente.

\textbf{Microcurrículo:} semana 14, sesión 1 (sesión global 27). Los
numerales del cronograma son identificadores curriculares; no son los
apartados del libro citados arriba.

\textbf{Implementación:}
\href{https://docs.scipy.org/doc/scipy/reference/generated/scipy.signal.filtfilt.html}{SciPy:
filtfilt}.
\end{frame}

\begin{frame}{112. Síntesis de la sesión}
\phantomsection\label{suxedntesis-de-la-sesiuxf3n-6}
\textbf{Libro guía:} John Semmlow, \emph{Circuits, Signals, and Systems
for Bioengineers}, 3.ª ed., 2018.

§8.4.2, \textbf{Finite Impulse Response Versus Infinite Impulse Response
Filter Characteristics} (página 363 del visor PDF); §8.4.3,
\textbf{Causal And Noncausal Filters} (página 365 del visor PDF);
§8.6.2, \textbf{Designing Infinite Impulse Response Filters} (página 387
del visor PDF).

\textbf{Correspondencia:} El tema se desarrolla en las secciones
indicadas. La redacción es una síntesis docente.

\textbf{Microcurrículo:} semana 14, sesión 1 (sesión global 27). Los
numerales del cronograma son identificadores curriculares; no son los
apartados del libro citados arriba.
\end{frame}

\begin{frame}{113. Sesión 28 · Fundamentos de wavelets}
\phantomsection\label{sesiuxf3n-28-fundamentos-de-wavelets}
\textbf{Libro guía:} John Semmlow, \emph{Circuits, Signals, and Systems
for Bioengineers}, 3.ª ed., 2018.

§4.6, \textbf{Time--Frequency Analysis} (página 205 del visor PDF).

\textbf{Correspondencia:} No existe una sección específica del libro
guía para el tema completo de esta diapositiva. Las secciones anteriores
son antecedentes, no referencias directas. Complemento directo: Semmlow
y Griffel (2014), §7.1--7.2.2 (Wavelet Analysis; Continuous Wavelet
Transform y características tiempo--frecuencia). Python: Esakkirajan et
al.~(2024), §5.7.

\textbf{Microcurrículo:} semana 14, sesión 2 (sesión global 28). Los
numerales del cronograma son identificadores curriculares; no son los
apartados del libro citados arriba.
\end{frame}

\begin{frame}{114. Preguntas que orientan la sesión}
\phantomsection\label{preguntas-que-orientan-la-sesiuxf3n-7}
\textbf{Libro guía:} John Semmlow, \emph{Circuits, Signals, and Systems
for Bioengineers}, 3.ª ed., 2018.

§4.6, \textbf{Time--Frequency Analysis} (página 205 del visor PDF).

\textbf{Correspondencia:} No existe una sección específica del libro
guía para el tema completo de esta diapositiva. Las secciones anteriores
son antecedentes, no referencias directas. Complemento directo: Semmlow
y Griffel (2014), §7.1--7.2.2 (Wavelet Analysis; Continuous Wavelet
Transform y características tiempo--frecuencia). Python: Esakkirajan et
al.~(2024), §5.7.

\textbf{Microcurrículo:} semana 14, sesión 2 (sesión global 28). Los
numerales del cronograma son identificadores curriculares; no son los
apartados del libro citados arriba.
\end{frame}

\begin{frame}{115. Vocabulario mínimo}
\phantomsection\label{vocabulario-muxednimo-7}
\textbf{Libro guía:} John Semmlow, \emph{Circuits, Signals, and Systems
for Bioengineers}, 3.ª ed., 2018.

§4.6, \textbf{Time--Frequency Analysis} (página 205 del visor PDF).

\textbf{Correspondencia:} No existe una sección específica del libro
guía para el tema completo de esta diapositiva. Las secciones anteriores
son antecedentes, no referencias directas. Complemento directo: Semmlow
y Griffel (2014), §7.1--7.2.2 (Wavelet Analysis; Continuous Wavelet
Transform y características tiempo--frecuencia). Python: Esakkirajan et
al.~(2024), §5.7.

\textbf{Microcurrículo:} semana 14, sesión 2 (sesión global 28). Los
numerales del cronograma son identificadores curriculares; no son los
apartados del libro citados arriba.
\end{frame}

\begin{frame}{116. Relaciones matemáticas centrales}
\phantomsection\label{relaciones-matemuxe1ticas-centrales-7}
\textbf{Libro guía:} John Semmlow, \emph{Circuits, Signals, and Systems
for Bioengineers}, 3.ª ed., 2018.

§4.6, \textbf{Time--Frequency Analysis} (página 205 del visor PDF).

\textbf{Correspondencia:} No existe una sección específica del libro
guía para el tema completo de esta diapositiva. Las secciones anteriores
son antecedentes, no referencias directas. Complemento directo: Semmlow
y Griffel (2014), §7.1--7.2.2 (Wavelet Analysis; Continuous Wavelet
Transform y características tiempo--frecuencia). Python: Esakkirajan et
al.~(2024), §5.7.

\textbf{Microcurrículo:} semana 14, sesión 2 (sesión global 28). Los
numerales del cronograma son identificadores curriculares; no son los
apartados del libro citados arriba.
\end{frame}

\begin{frame}{117. Procedimiento de análisis}
\phantomsection\label{procedimiento-de-anuxe1lisis-7}
\textbf{Libro guía:} John Semmlow, \emph{Circuits, Signals, and Systems
for Bioengineers}, 3.ª ed., 2018.

§4.6, \textbf{Time--Frequency Analysis} (página 205 del visor PDF).

\textbf{Correspondencia:} No existe una sección específica del libro
guía para el tema completo de esta diapositiva. Las secciones anteriores
son antecedentes, no referencias directas. Complemento directo: Semmlow
y Griffel (2014), §7.1--7.2.2 (Wavelet Analysis; Continuous Wavelet
Transform y características tiempo--frecuencia). Python: Esakkirajan et
al.~(2024), §5.7.

\textbf{Microcurrículo:} semana 14, sesión 2 (sesión global 28). Los
numerales del cronograma son identificadores curriculares; no son los
apartados del libro citados arriba.
\end{frame}

\begin{frame}{118. Ejemplo biomédico}
\phantomsection\label{ejemplo-biomuxe9dico-7}
\textbf{Libro guía:} John Semmlow, \emph{Circuits, Signals, and Systems
for Bioengineers}, 3.ª ed., 2018.

§4.6, \textbf{Time--Frequency Analysis} (página 205 del visor PDF).

\textbf{Correspondencia:} No existe una sección específica del libro
guía para el tema completo de esta diapositiva. Las secciones anteriores
son antecedentes, no referencias directas. Complemento directo: Semmlow
y Griffel (2014), §7.1--7.2.2 (Wavelet Analysis; Continuous Wavelet
Transform y características tiempo--frecuencia). Python: Esakkirajan et
al.~(2024), §5.7.

\textbf{Microcurrículo:} semana 14, sesión 2 (sesión global 28). Los
numerales del cronograma son identificadores curriculares; no son los
apartados del libro citados arriba.
\end{frame}

\begin{frame}{119. Implementación mínima en Python}
\phantomsection\label{implementaciuxf3n-muxednima-en-python-7}
\textbf{Libro guía:} John Semmlow, \emph{Circuits, Signals, and Systems
for Bioengineers}, 3.ª ed., 2018.

§4.6, \textbf{Time--Frequency Analysis} (página 205 del visor PDF).

\textbf{Correspondencia:} No existe una sección específica del libro
guía para el tema completo de esta diapositiva. Las secciones anteriores
son antecedentes, no referencias directas. Complemento directo: Semmlow
y Griffel (2014), §7.1--7.2.2 (Wavelet Analysis; Continuous Wavelet
Transform y características tiempo--frecuencia). Python: Esakkirajan et
al.~(2024), §5.7.

\textbf{Microcurrículo:} semana 14, sesión 2 (sesión global 28). Los
numerales del cronograma son identificadores curriculares; no son los
apartados del libro citados arriba.

\textbf{Implementación:}
\href{https://pywavelets.readthedocs.io/en/latest/ref/index.html}{PyWavelets:
transformadas y convenciones}.
\end{frame}

\begin{frame}{120. Admisibilidad y escala física}
\phantomsection\label{admisibilidad-y-escala-fuxedsica}
\textbf{Libro guía:} John Semmlow, \emph{Circuits, Signals, and Systems
for Bioengineers}, 3.ª ed., 2018.

§4.6, \textbf{Time--Frequency Analysis} (página 205 del visor PDF).

\textbf{Correspondencia:} No existe una sección específica del libro
guía para el tema completo de esta diapositiva. Las secciones anteriores
son antecedentes, no referencias directas. Complemento directo: Semmlow
y Griffel (2014), §7.2 y §7.2.1, PDF 242 y 244. Implementación:
documentación CWT de PyWavelets. Complemento directo: Semmlow y Griffel
(2014), §7.1--7.2.2 (Wavelet Analysis; Continuous Wavelet Transform y
características tiempo--frecuencia). Python: Esakkirajan et al.~(2024),
§5.7.

\textbf{Microcurrículo:} semana 14, sesión 2 (sesión global 28). Los
numerales del cronograma son identificadores curriculares; no son los
apartados del libro citados arriba.
\end{frame}

\begin{frame}{121. Verificaciones antes de interpretar}
\phantomsection\label{verificaciones-antes-de-interpretar-7}
\textbf{Libro guía:} John Semmlow, \emph{Circuits, Signals, and Systems
for Bioengineers}, 3.ª ed., 2018.

§4.6, \textbf{Time--Frequency Analysis} (página 205 del visor PDF).

\textbf{Correspondencia:} No existe una sección específica del libro
guía para el tema completo de esta diapositiva. Las secciones anteriores
son antecedentes, no referencias directas. Complemento directo: Semmlow
y Griffel (2014), §7.1--7.2.2 (Wavelet Analysis; Continuous Wavelet
Transform y características tiempo--frecuencia). Python: Esakkirajan et
al.~(2024), §5.7.

\textbf{Microcurrículo:} semana 14, sesión 2 (sesión global 28). Los
numerales del cronograma son identificadores curriculares; no son los
apartados del libro citados arriba.
\end{frame}

\begin{frame}{122. Errores frecuentes}
\phantomsection\label{errores-frecuentes-7}
\textbf{Libro guía:} John Semmlow, \emph{Circuits, Signals, and Systems
for Bioengineers}, 3.ª ed., 2018.

§4.6, \textbf{Time--Frequency Analysis} (página 205 del visor PDF).

\textbf{Correspondencia:} No existe una sección específica del libro
guía para el tema completo de esta diapositiva. Las secciones anteriores
son antecedentes, no referencias directas. Complemento directo: Semmlow
y Griffel (2014), §7.1--7.2.2 (Wavelet Analysis; Continuous Wavelet
Transform y características tiempo--frecuencia). Python: Esakkirajan et
al.~(2024), §5.7.

\textbf{Microcurrículo:} semana 14, sesión 2 (sesión global 28). Los
numerales del cronograma son identificadores curriculares; no son los
apartados del libro citados arriba.
\end{frame}

\begin{frame}{123. Actividad guiada}
\phantomsection\label{actividad-guiada-7}
\textbf{Libro guía:} John Semmlow, \emph{Circuits, Signals, and Systems
for Bioengineers}, 3.ª ed., 2018.

§4.6, \textbf{Time--Frequency Analysis} (página 205 del visor PDF).

\textbf{Correspondencia:} No existe una sección específica del libro
guía para el tema completo de esta diapositiva. Las secciones anteriores
son antecedentes, no referencias directas. Complemento directo: Semmlow
y Griffel (2014), §7.1--7.2.2 (Wavelet Analysis; Continuous Wavelet
Transform y características tiempo--frecuencia). Python: Esakkirajan et
al.~(2024), §5.7.

\textbf{Microcurrículo:} semana 14, sesión 2 (sesión global 28). Los
numerales del cronograma son identificadores curriculares; no son los
apartados del libro citados arriba.
\end{frame}

\begin{frame}{124. Pregunta de comprobación}
\phantomsection\label{pregunta-de-comprobaciuxf3n-7}
\textbf{Libro guía:} John Semmlow, \emph{Circuits, Signals, and Systems
for Bioengineers}, 3.ª ed., 2018.

§4.6, \textbf{Time--Frequency Analysis} (página 205 del visor PDF).

\textbf{Correspondencia:} No existe una sección específica del libro
guía para el tema completo de esta diapositiva. Las secciones anteriores
son antecedentes, no referencias directas. Complemento directo: Semmlow
y Griffel (2014), §7.1--7.2.2 (Wavelet Analysis; Continuous Wavelet
Transform y características tiempo--frecuencia). Python: Esakkirajan et
al.~(2024), §5.7.

\textbf{Microcurrículo:} semana 14, sesión 2 (sesión global 28). Los
numerales del cronograma son identificadores curriculares; no son los
apartados del libro citados arriba.
\end{frame}

\begin{frame}{125. Síntesis de la sesión}
\phantomsection\label{suxedntesis-de-la-sesiuxf3n-7}
\textbf{Libro guía:} John Semmlow, \emph{Circuits, Signals, and Systems
for Bioengineers}, 3.ª ed., 2018.

§4.6, \textbf{Time--Frequency Analysis} (página 205 del visor PDF).

\textbf{Correspondencia:} No existe una sección específica del libro
guía para el tema completo de esta diapositiva. Las secciones anteriores
son antecedentes, no referencias directas. Complemento directo: Semmlow
y Griffel (2014), §7.1--7.2.2 (Wavelet Analysis; Continuous Wavelet
Transform y características tiempo--frecuencia). Python: Esakkirajan et
al.~(2024), §5.7.

\textbf{Microcurrículo:} semana 14, sesión 2 (sesión global 28). Los
numerales del cronograma son identificadores curriculares; no son los
apartados del libro citados arriba.
\end{frame}

\begin{frame}{126. Semana 15 · Multirresolución y clasificación}
\phantomsection\label{semana-15-multirresoluciuxf3n-y-clasificaciuxf3n}
\textbf{Libro guía:} John Semmlow, \emph{Circuits, Signals, and Systems
for Bioengineers}, 3.ª ed., 2018.

§5.4, \textbf{Using The Impulse Response To Find A Systems Output To Any
Input---Convolution} (página 219 del visor PDF); §6.4, \textbf{The
Transfer Function} (página 255 del visor PDF); §7.4, \textbf{Nonzero
Initial Conditions---Initial And Final Value Theorems} (página 326 del
visor PDF); §8.4, \textbf{Linear Filters---Introduction} (página 358 del
visor PDF).

\textbf{Correspondencia:} Síntesis o encuadre que reúne varias
secciones; no corresponde a un único apartado del libro.

\textbf{Lectura:} use estas secciones como ruta de preparación y
consulta. La bibliografía aparece al final del bloque.
\end{frame}

\begin{frame}{127. Sesión 29 · Análisis multirresolución y por bandas}
\phantomsection\label{sesiuxf3n-29-anuxe1lisis-multirresoluciuxf3n-y-por-bandas}
\textbf{Libro guía:} John Semmlow, \emph{Circuits, Signals, and Systems
for Bioengineers}, 3.ª ed., 2018.

§8.4, \textbf{Linear Filters---Introduction} (página 358 del visor PDF);
§8.5, \textbf{Design Of Finite Impulse Response Filters} (página 366 del
visor PDF).

\textbf{Correspondencia:} No existe una sección específica del libro
guía para el tema completo de esta diapositiva. Las secciones anteriores
son antecedentes, no referencias directas. Complemento directo: Semmlow
y Griffel (2014), §7.3--7.3.2 (Discrete Wavelet Transform; Filter
Banks). Python: Esakkirajan et al.~(2024), §5.8.

\textbf{Microcurrículo:} semana 15, sesión 1 (sesión global 29). Los
numerales del cronograma son identificadores curriculares; no son los
apartados del libro citados arriba.
\end{frame}

\begin{frame}{128. Preguntas que orientan la sesión}
\phantomsection\label{preguntas-que-orientan-la-sesiuxf3n-8}
\textbf{Libro guía:} John Semmlow, \emph{Circuits, Signals, and Systems
for Bioengineers}, 3.ª ed., 2018.

§8.4, \textbf{Linear Filters---Introduction} (página 358 del visor PDF);
§8.5, \textbf{Design Of Finite Impulse Response Filters} (página 366 del
visor PDF).

\textbf{Correspondencia:} No existe una sección específica del libro
guía para el tema completo de esta diapositiva. Las secciones anteriores
son antecedentes, no referencias directas. Complemento directo: Semmlow
y Griffel (2014), §7.3--7.3.2 (Discrete Wavelet Transform; Filter
Banks). Python: Esakkirajan et al.~(2024), §5.8.

\textbf{Microcurrículo:} semana 15, sesión 1 (sesión global 29). Los
numerales del cronograma son identificadores curriculares; no son los
apartados del libro citados arriba.
\end{frame}

\begin{frame}{129. Vocabulario mínimo}
\phantomsection\label{vocabulario-muxednimo-8}
\textbf{Libro guía:} John Semmlow, \emph{Circuits, Signals, and Systems
for Bioengineers}, 3.ª ed., 2018.

§8.4, \textbf{Linear Filters---Introduction} (página 358 del visor PDF);
§8.5, \textbf{Design Of Finite Impulse Response Filters} (página 366 del
visor PDF).

\textbf{Correspondencia:} No existe una sección específica del libro
guía para el tema completo de esta diapositiva. Las secciones anteriores
son antecedentes, no referencias directas. Complemento directo: Semmlow
y Griffel (2014), §7.3--7.3.2 (Discrete Wavelet Transform; Filter
Banks). Python: Esakkirajan et al.~(2024), §5.8.

\textbf{Microcurrículo:} semana 15, sesión 1 (sesión global 29). Los
numerales del cronograma son identificadores curriculares; no son los
apartados del libro citados arriba.
\end{frame}

\begin{frame}{130. Relaciones matemáticas centrales}
\phantomsection\label{relaciones-matemuxe1ticas-centrales-8}
\textbf{Libro guía:} John Semmlow, \emph{Circuits, Signals, and Systems
for Bioengineers}, 3.ª ed., 2018.

§8.4, \textbf{Linear Filters---Introduction} (página 358 del visor PDF);
§8.5, \textbf{Design Of Finite Impulse Response Filters} (página 366 del
visor PDF).

\textbf{Correspondencia:} No existe una sección específica del libro
guía para el tema completo de esta diapositiva. Las secciones anteriores
son antecedentes, no referencias directas. Complemento directo: Semmlow
y Griffel (2014), §7.3--7.3.2 (Discrete Wavelet Transform; Filter
Banks). Python: Esakkirajan et al.~(2024), §5.8.

\textbf{Microcurrículo:} semana 15, sesión 1 (sesión global 29). Los
numerales del cronograma son identificadores curriculares; no son los
apartados del libro citados arriba.
\end{frame}

\begin{frame}{131. Procedimiento de análisis}
\phantomsection\label{procedimiento-de-anuxe1lisis-8}
\textbf{Libro guía:} John Semmlow, \emph{Circuits, Signals, and Systems
for Bioengineers}, 3.ª ed., 2018.

§8.4, \textbf{Linear Filters---Introduction} (página 358 del visor PDF);
§8.5, \textbf{Design Of Finite Impulse Response Filters} (página 366 del
visor PDF).

\textbf{Correspondencia:} No existe una sección específica del libro
guía para el tema completo de esta diapositiva. Las secciones anteriores
son antecedentes, no referencias directas. Complemento directo: Semmlow
y Griffel (2014), §7.3--7.3.2 (Discrete Wavelet Transform; Filter
Banks). Python: Esakkirajan et al.~(2024), §5.8.

\textbf{Microcurrículo:} semana 15, sesión 1 (sesión global 29). Los
numerales del cronograma son identificadores curriculares; no son los
apartados del libro citados arriba.
\end{frame}

\begin{frame}{132. Ejemplo biomédico}
\phantomsection\label{ejemplo-biomuxe9dico-8}
\textbf{Libro guía:} John Semmlow, \emph{Circuits, Signals, and Systems
for Bioengineers}, 3.ª ed., 2018.

§8.4, \textbf{Linear Filters---Introduction} (página 358 del visor PDF);
§8.5, \textbf{Design Of Finite Impulse Response Filters} (página 366 del
visor PDF).

\textbf{Correspondencia:} No existe una sección específica del libro
guía para el tema completo de esta diapositiva. Las secciones anteriores
son antecedentes, no referencias directas. Complemento directo: Semmlow
y Griffel (2014), §7.3--7.3.2 (Discrete Wavelet Transform; Filter
Banks). Python: Esakkirajan et al.~(2024), §5.8.

\textbf{Microcurrículo:} semana 15, sesión 1 (sesión global 29). Los
numerales del cronograma son identificadores curriculares; no son los
apartados del libro citados arriba.
\end{frame}

\begin{frame}{133. Implementación mínima en Python}
\phantomsection\label{implementaciuxf3n-muxednima-en-python-8}
\textbf{Libro guía:} John Semmlow, \emph{Circuits, Signals, and Systems
for Bioengineers}, 3.ª ed., 2018.

§8.4, \textbf{Linear Filters---Introduction} (página 358 del visor PDF);
§8.5, \textbf{Design Of Finite Impulse Response Filters} (página 366 del
visor PDF).

\textbf{Correspondencia:} No existe una sección específica del libro
guía para el tema completo de esta diapositiva. Las secciones anteriores
son antecedentes, no referencias directas. Complemento directo: Semmlow
y Griffel (2014), §7.3--7.3.2 (Discrete Wavelet Transform; Filter
Banks). Python: Esakkirajan et al.~(2024), §5.8.

\textbf{Microcurrículo:} semana 15, sesión 1 (sesión global 29). Los
numerales del cronograma son identificadores curriculares; no son los
apartados del libro citados arriba.

\textbf{Implementación:}
\href{https://pywavelets.readthedocs.io/en/latest/ref/index.html}{PyWavelets:
transformadas y convenciones}.
\end{frame}

\begin{frame}{134. Reconstrucción de detalles a la tasa original}
\phantomsection\label{reconstrucciuxf3n-de-detalles-a-la-tasa-original}
\textbf{Libro guía:} John Semmlow, \emph{Circuits, Signals, and Systems
for Bioengineers}, 3.ª ed., 2018.

§8.3, \textbf{Difference Equations} (página 355 del visor PDF); §8.4,
\textbf{Linear Filters---Introduction} (página 358 del visor PDF).

\textbf{Correspondencia:} No existe una sección específica del libro
guía para el tema completo de esta diapositiva. Las secciones anteriores
son antecedentes, no referencias directas. Complemento directo: Semmlow
y Griffel (2014), §7.3.1 y §7.3.2, PDF 249 y 254. Complemento directo:
Semmlow y Griffel (2014), §7.3--7.3.2 (Discrete Wavelet Transform;
Filter Banks). Python: Esakkirajan et al.~(2024), §5.8.

\textbf{Microcurrículo:} semana 15, sesión 1 (sesión global 29). Los
numerales del cronograma son identificadores curriculares; no son los
apartados del libro citados arriba.

\textbf{Implementación:}
\href{https://pywavelets.readthedocs.io/en/latest/ref/index.html}{PyWavelets:
transformadas y convenciones}.
\end{frame}

\begin{frame}{135. Verificaciones antes de interpretar}
\phantomsection\label{verificaciones-antes-de-interpretar-8}
\textbf{Libro guía:} John Semmlow, \emph{Circuits, Signals, and Systems
for Bioengineers}, 3.ª ed., 2018.

§8.4, \textbf{Linear Filters---Introduction} (página 358 del visor PDF);
§8.5, \textbf{Design Of Finite Impulse Response Filters} (página 366 del
visor PDF).

\textbf{Correspondencia:} No existe una sección específica del libro
guía para el tema completo de esta diapositiva. Las secciones anteriores
son antecedentes, no referencias directas. Complemento directo: Semmlow
y Griffel (2014), §7.3--7.3.2 (Discrete Wavelet Transform; Filter
Banks). Python: Esakkirajan et al.~(2024), §5.8.

\textbf{Microcurrículo:} semana 15, sesión 1 (sesión global 29). Los
numerales del cronograma son identificadores curriculares; no son los
apartados del libro citados arriba.
\end{frame}

\begin{frame}{136. Errores frecuentes}
\phantomsection\label{errores-frecuentes-8}
\textbf{Libro guía:} John Semmlow, \emph{Circuits, Signals, and Systems
for Bioengineers}, 3.ª ed., 2018.

§8.4, \textbf{Linear Filters---Introduction} (página 358 del visor PDF);
§8.5, \textbf{Design Of Finite Impulse Response Filters} (página 366 del
visor PDF).

\textbf{Correspondencia:} No existe una sección específica del libro
guía para el tema completo de esta diapositiva. Las secciones anteriores
son antecedentes, no referencias directas. Complemento directo: Semmlow
y Griffel (2014), §7.3--7.3.2 (Discrete Wavelet Transform; Filter
Banks). Python: Esakkirajan et al.~(2024), §5.8.

\textbf{Microcurrículo:} semana 15, sesión 1 (sesión global 29). Los
numerales del cronograma son identificadores curriculares; no son los
apartados del libro citados arriba.
\end{frame}

\begin{frame}{137. Actividad guiada}
\phantomsection\label{actividad-guiada-8}
\textbf{Libro guía:} John Semmlow, \emph{Circuits, Signals, and Systems
for Bioengineers}, 3.ª ed., 2018.

§8.4, \textbf{Linear Filters---Introduction} (página 358 del visor PDF);
§8.5, \textbf{Design Of Finite Impulse Response Filters} (página 366 del
visor PDF).

\textbf{Correspondencia:} No existe una sección específica del libro
guía para el tema completo de esta diapositiva. Las secciones anteriores
son antecedentes, no referencias directas. Complemento directo: Semmlow
y Griffel (2014), §7.3--7.3.2 (Discrete Wavelet Transform; Filter
Banks). Python: Esakkirajan et al.~(2024), §5.8.

\textbf{Microcurrículo:} semana 15, sesión 1 (sesión global 29). Los
numerales del cronograma son identificadores curriculares; no son los
apartados del libro citados arriba.
\end{frame}

\begin{frame}{138. Pregunta de comprobación}
\phantomsection\label{pregunta-de-comprobaciuxf3n-8}
\textbf{Libro guía:} John Semmlow, \emph{Circuits, Signals, and Systems
for Bioengineers}, 3.ª ed., 2018.

§8.4, \textbf{Linear Filters---Introduction} (página 358 del visor PDF);
§8.5, \textbf{Design Of Finite Impulse Response Filters} (página 366 del
visor PDF).

\textbf{Correspondencia:} No existe una sección específica del libro
guía para el tema completo de esta diapositiva. Las secciones anteriores
son antecedentes, no referencias directas. Complemento directo: Semmlow
y Griffel (2014), §7.3--7.3.2 (Discrete Wavelet Transform; Filter
Banks). Python: Esakkirajan et al.~(2024), §5.8.

\textbf{Microcurrículo:} semana 15, sesión 1 (sesión global 29). Los
numerales del cronograma son identificadores curriculares; no son los
apartados del libro citados arriba.
\end{frame}

\begin{frame}{139. Síntesis de la sesión}
\phantomsection\label{suxedntesis-de-la-sesiuxf3n-8}
\textbf{Libro guía:} John Semmlow, \emph{Circuits, Signals, and Systems
for Bioengineers}, 3.ª ed., 2018.

§8.4, \textbf{Linear Filters---Introduction} (página 358 del visor PDF);
§8.5, \textbf{Design Of Finite Impulse Response Filters} (página 366 del
visor PDF).

\textbf{Correspondencia:} No existe una sección específica del libro
guía para el tema completo de esta diapositiva. Las secciones anteriores
son antecedentes, no referencias directas. Complemento directo: Semmlow
y Griffel (2014), §7.3--7.3.2 (Discrete Wavelet Transform; Filter
Banks). Python: Esakkirajan et al.~(2024), §5.8.

\textbf{Microcurrículo:} semana 15, sesión 1 (sesión global 29). Los
numerales del cronograma son identificadores curriculares; no son los
apartados del libro citados arriba.
\end{frame}

\begin{frame}{140. Sesión 30 · Sistemas lineales para clasificación}
\phantomsection\label{sesiuxf3n-30-sistemas-lineales-para-clasificaciuxf3n}
\textbf{Libro guía:} John Semmlow, \emph{Circuits, Signals, and Systems
for Bioengineers}, 3.ª ed., 2018.

§1.4.4, \textbf{Linear Versus Nonlinear Signals And Systems} (página 44
del visor PDF); §8.4, \textbf{Linear Filters---Introduction} (página 358
del visor PDF).

\textbf{Correspondencia:} No existe una sección específica del libro
guía para el tema completo de esta diapositiva. Las secciones anteriores
son antecedentes, no referencias directas. Complemento: Semmlow y
Griffel (2014), §16.2 (Linear Discriminators), §16.3 (Evaluating
Classifier Performance) y §16.6 (capacidad y sobreajuste). La partición
por sujeto es una ampliación metodológica docente.

\textbf{Microcurrículo:} semana 15, sesión 2 (sesión global 30). Los
numerales del cronograma son identificadores curriculares; no son los
apartados del libro citados arriba.
\end{frame}

\begin{frame}{141. Preguntas que orientan la sesión}
\phantomsection\label{preguntas-que-orientan-la-sesiuxf3n-9}
\textbf{Libro guía:} John Semmlow, \emph{Circuits, Signals, and Systems
for Bioengineers}, 3.ª ed., 2018.

§1.4.4, \textbf{Linear Versus Nonlinear Signals And Systems} (página 44
del visor PDF); §8.4, \textbf{Linear Filters---Introduction} (página 358
del visor PDF).

\textbf{Correspondencia:} No existe una sección específica del libro
guía para el tema completo de esta diapositiva. Las secciones anteriores
son antecedentes, no referencias directas. Complemento: Semmlow y
Griffel (2014), §16.2 (Linear Discriminators), §16.3 (Evaluating
Classifier Performance) y §16.6 (capacidad y sobreajuste). La partición
por sujeto es una ampliación metodológica docente.

\textbf{Microcurrículo:} semana 15, sesión 2 (sesión global 30). Los
numerales del cronograma son identificadores curriculares; no son los
apartados del libro citados arriba.
\end{frame}

\begin{frame}{142. Vocabulario mínimo}
\phantomsection\label{vocabulario-muxednimo-9}
\textbf{Libro guía:} John Semmlow, \emph{Circuits, Signals, and Systems
for Bioengineers}, 3.ª ed., 2018.

§1.4.4, \textbf{Linear Versus Nonlinear Signals And Systems} (página 44
del visor PDF); §8.4, \textbf{Linear Filters---Introduction} (página 358
del visor PDF).

\textbf{Correspondencia:} No existe una sección específica del libro
guía para el tema completo de esta diapositiva. Las secciones anteriores
son antecedentes, no referencias directas. Complemento: Semmlow y
Griffel (2014), §16.2 (Linear Discriminators), §16.3 (Evaluating
Classifier Performance) y §16.6 (capacidad y sobreajuste). La partición
por sujeto es una ampliación metodológica docente.

\textbf{Microcurrículo:} semana 15, sesión 2 (sesión global 30). Los
numerales del cronograma son identificadores curriculares; no son los
apartados del libro citados arriba.
\end{frame}

\begin{frame}{143. Relaciones matemáticas centrales}
\phantomsection\label{relaciones-matemuxe1ticas-centrales-9}
\textbf{Libro guía:} John Semmlow, \emph{Circuits, Signals, and Systems
for Bioengineers}, 3.ª ed., 2018.

§1.4.4, \textbf{Linear Versus Nonlinear Signals And Systems} (página 44
del visor PDF); §8.4, \textbf{Linear Filters---Introduction} (página 358
del visor PDF).

\textbf{Correspondencia:} No existe una sección específica del libro
guía para el tema completo de esta diapositiva. Las secciones anteriores
son antecedentes, no referencias directas. Complemento: Semmlow y
Griffel (2014), §16.2 (Linear Discriminators), §16.3 (Evaluating
Classifier Performance) y §16.6 (capacidad y sobreajuste). La partición
por sujeto es una ampliación metodológica docente.

\textbf{Microcurrículo:} semana 15, sesión 2 (sesión global 30). Los
numerales del cronograma son identificadores curriculares; no son los
apartados del libro citados arriba.
\end{frame}

\begin{frame}{144. Discriminante lineal: alcance del modelo}
\phantomsection\label{discriminante-lineal-alcance-del-modelo}
\textbf{Libro guía:} John Semmlow, \emph{Circuits, Signals, and Systems
for Bioengineers}, 3.ª ed., 2018.

§8.4, \textbf{Linear Filters---Introduction} (página 358 del visor PDF).

\textbf{Correspondencia:} No existe una sección específica del libro
guía para el tema completo de esta diapositiva. Las secciones anteriores
son antecedentes, no referencias directas. Complemento directo: Semmlow
y Griffel (2014), §16.2. La aplicación y la partición por sujeto son
desarrollo docente. Complemento: Semmlow y Griffel (2014), §16.2 (Linear
Discriminators), §16.3 (Evaluating Classifier Performance) y §16.6
(capacidad y sobreajuste). La partición por sujeto es una ampliación
metodológica docente.

\textbf{Microcurrículo:} semana 15, sesión 2 (sesión global 30). Los
numerales del cronograma son identificadores curriculares; no son los
apartados del libro citados arriba.

\textbf{Ampliación específica de LDA:}
\href{https://scikit-learn.org/stable/modules/lda_qda.html}{scikit-learn:
derivación y supuestos de LDA}. El §16.2 del complemento introduce
discriminadores lineales y desarrolla otro ajuste; no se atribuye a ese
apartado esta derivación generativa exacta.
\end{frame}

\begin{frame}{145. Procedimiento de análisis}
\phantomsection\label{procedimiento-de-anuxe1lisis-9}
\textbf{Libro guía:} John Semmlow, \emph{Circuits, Signals, and Systems
for Bioengineers}, 3.ª ed., 2018.

§1.4.4, \textbf{Linear Versus Nonlinear Signals And Systems} (página 44
del visor PDF); §8.4, \textbf{Linear Filters---Introduction} (página 358
del visor PDF).

\textbf{Correspondencia:} No existe una sección específica del libro
guía para el tema completo de esta diapositiva. Las secciones anteriores
son antecedentes, no referencias directas. Complemento: Semmlow y
Griffel (2014), §16.2 (Linear Discriminators), §16.3 (Evaluating
Classifier Performance) y §16.6 (capacidad y sobreajuste). La partición
por sujeto es una ampliación metodológica docente.

\textbf{Microcurrículo:} semana 15, sesión 2 (sesión global 30). Los
numerales del cronograma son identificadores curriculares; no son los
apartados del libro citados arriba.
\end{frame}

\begin{frame}{146. Ejemplo biomédico}
\phantomsection\label{ejemplo-biomuxe9dico-9}
\textbf{Libro guía:} John Semmlow, \emph{Circuits, Signals, and Systems
for Bioengineers}, 3.ª ed., 2018.

§1.4.4, \textbf{Linear Versus Nonlinear Signals And Systems} (página 44
del visor PDF); §8.4, \textbf{Linear Filters---Introduction} (página 358
del visor PDF).

\textbf{Correspondencia:} No existe una sección específica del libro
guía para el tema completo de esta diapositiva. Las secciones anteriores
son antecedentes, no referencias directas. Complemento: Semmlow y
Griffel (2014), §16.2 (Linear Discriminators), §16.3 (Evaluating
Classifier Performance) y §16.6 (capacidad y sobreajuste). La partición
por sujeto es una ampliación metodológica docente.

\textbf{Microcurrículo:} semana 15, sesión 2 (sesión global 30). Los
numerales del cronograma son identificadores curriculares; no son los
apartados del libro citados arriba.
\end{frame}

\begin{frame}{147. Implementación mínima en Python}
\phantomsection\label{implementaciuxf3n-muxednima-en-python-9}
\textbf{Libro guía:} John Semmlow, \emph{Circuits, Signals, and Systems
for Bioengineers}, 3.ª ed., 2018.

§1.4.4, \textbf{Linear Versus Nonlinear Signals And Systems} (página 44
del visor PDF); §8.4, \textbf{Linear Filters---Introduction} (página 358
del visor PDF).

\textbf{Correspondencia:} No existe una sección específica del libro
guía para el tema completo de esta diapositiva. Las secciones anteriores
son antecedentes, no referencias directas. Complemento: Semmlow y
Griffel (2014), §16.2 (Linear Discriminators), §16.3 (Evaluating
Classifier Performance) y §16.6 (capacidad y sobreajuste). La partición
por sujeto es una ampliación metodológica docente.

\textbf{Microcurrículo:} semana 15, sesión 2 (sesión global 30). Los
numerales del cronograma son identificadores curriculares; no son los
apartados del libro citados arriba.

\textbf{Ampliación específica de LDA:}
\href{https://scikit-learn.org/stable/modules/lda_qda.html}{scikit-learn:
derivación y supuestos de LDA}. El §16.2 del complemento introduce
discriminadores lineales y desarrolla otro ajuste; no se atribuye a ese
apartado esta derivación generativa exacta.
\end{frame}

\begin{frame}{148. Verificaciones antes de interpretar}
\phantomsection\label{verificaciones-antes-de-interpretar-9}
\textbf{Libro guía:} John Semmlow, \emph{Circuits, Signals, and Systems
for Bioengineers}, 3.ª ed., 2018.

§1.4.4, \textbf{Linear Versus Nonlinear Signals And Systems} (página 44
del visor PDF); §8.4, \textbf{Linear Filters---Introduction} (página 358
del visor PDF).

\textbf{Correspondencia:} No existe una sección específica del libro
guía para el tema completo de esta diapositiva. Las secciones anteriores
son antecedentes, no referencias directas. Complemento: Semmlow y
Griffel (2014), §16.2 (Linear Discriminators), §16.3 (Evaluating
Classifier Performance) y §16.6 (capacidad y sobreajuste). La partición
por sujeto es una ampliación metodológica docente.

\textbf{Microcurrículo:} semana 15, sesión 2 (sesión global 30). Los
numerales del cronograma son identificadores curriculares; no son los
apartados del libro citados arriba.
\end{frame}

\begin{frame}{149. Errores frecuentes}
\phantomsection\label{errores-frecuentes-9}
\textbf{Libro guía:} John Semmlow, \emph{Circuits, Signals, and Systems
for Bioengineers}, 3.ª ed., 2018.

§1.4.4, \textbf{Linear Versus Nonlinear Signals And Systems} (página 44
del visor PDF); §8.4, \textbf{Linear Filters---Introduction} (página 358
del visor PDF).

\textbf{Correspondencia:} No existe una sección específica del libro
guía para el tema completo de esta diapositiva. Las secciones anteriores
son antecedentes, no referencias directas. Complemento: Semmlow y
Griffel (2014), §16.2 (Linear Discriminators), §16.3 (Evaluating
Classifier Performance) y §16.6 (capacidad y sobreajuste). La partición
por sujeto es una ampliación metodológica docente.

\textbf{Microcurrículo:} semana 15, sesión 2 (sesión global 30). Los
numerales del cronograma son identificadores curriculares; no son los
apartados del libro citados arriba.
\end{frame}

\begin{frame}{150. Actividad guiada}
\phantomsection\label{actividad-guiada-9}
\textbf{Libro guía:} John Semmlow, \emph{Circuits, Signals, and Systems
for Bioengineers}, 3.ª ed., 2018.

§1.4.4, \textbf{Linear Versus Nonlinear Signals And Systems} (página 44
del visor PDF); §8.4, \textbf{Linear Filters---Introduction} (página 358
del visor PDF).

\textbf{Correspondencia:} No existe una sección específica del libro
guía para el tema completo de esta diapositiva. Las secciones anteriores
son antecedentes, no referencias directas. Complemento: Semmlow y
Griffel (2014), §16.2 (Linear Discriminators), §16.3 (Evaluating
Classifier Performance) y §16.6 (capacidad y sobreajuste). La partición
por sujeto es una ampliación metodológica docente.

\textbf{Microcurrículo:} semana 15, sesión 2 (sesión global 30). Los
numerales del cronograma son identificadores curriculares; no son los
apartados del libro citados arriba.
\end{frame}

\begin{frame}{151. Pregunta de comprobación}
\phantomsection\label{pregunta-de-comprobaciuxf3n-9}
\textbf{Libro guía:} John Semmlow, \emph{Circuits, Signals, and Systems
for Bioengineers}, 3.ª ed., 2018.

§1.4.4, \textbf{Linear Versus Nonlinear Signals And Systems} (página 44
del visor PDF); §8.4, \textbf{Linear Filters---Introduction} (página 358
del visor PDF).

\textbf{Correspondencia:} No existe una sección específica del libro
guía para el tema completo de esta diapositiva. Las secciones anteriores
son antecedentes, no referencias directas. Complemento: Semmlow y
Griffel (2014), §16.2 (Linear Discriminators), §16.3 (Evaluating
Classifier Performance) y §16.6 (capacidad y sobreajuste). La partición
por sujeto es una ampliación metodológica docente.

\textbf{Microcurrículo:} semana 15, sesión 2 (sesión global 30). Los
numerales del cronograma son identificadores curriculares; no son los
apartados del libro citados arriba.
\end{frame}

\begin{frame}{152. Síntesis de la sesión}
\phantomsection\label{suxedntesis-de-la-sesiuxf3n-9}
\textbf{Libro guía:} John Semmlow, \emph{Circuits, Signals, and Systems
for Bioengineers}, 3.ª ed., 2018.

§1.4.4, \textbf{Linear Versus Nonlinear Signals And Systems} (página 44
del visor PDF); §8.4, \textbf{Linear Filters---Introduction} (página 358
del visor PDF).

\textbf{Correspondencia:} No existe una sección específica del libro
guía para el tema completo de esta diapositiva. Las secciones anteriores
son antecedentes, no referencias directas. Complemento: Semmlow y
Griffel (2014), §16.2 (Linear Discriminators), §16.3 (Evaluating
Classifier Performance) y §16.6 (capacidad y sobreajuste). La partición
por sujeto es una ampliación metodológica docente.

\textbf{Microcurrículo:} semana 15, sesión 2 (sesión global 30). Los
numerales del cronograma son identificadores curriculares; no son los
apartados del libro citados arriba.
\end{frame}

\begin{frame}{153. Síntesis y estudio autónomo}
\phantomsection\label{suxedntesis-y-estudio-autuxf3nomo-2}
\textbf{Libro guía:} John Semmlow, \emph{Circuits, Signals, and Systems
for Bioengineers}, 3.ª ed., 2018.

§5.4, \textbf{Using The Impulse Response To Find A Systems Output To Any
Input---Convolution} (página 219 del visor PDF); §6.4, \textbf{The
Transfer Function} (página 255 del visor PDF); §7.4, \textbf{Nonzero
Initial Conditions---Initial And Final Value Theorems} (página 326 del
visor PDF); §8.4, \textbf{Linear Filters---Introduction} (página 358 del
visor PDF).

\textbf{Correspondencia:} Síntesis o encuadre que reúne varias
secciones; no corresponde a un único apartado del libro.

\textbf{Lectura:} use estas secciones como ruta de preparación y
consulta. La bibliografía aparece al final del bloque.
\end{frame}

\begin{frame}{154. Un mismo sistema, cuatro representaciones}
\phantomsection\label{un-mismo-sistema-cuatro-representaciones}
\textbf{Libro guía:} John Semmlow, \emph{Circuits, Signals, and Systems
for Bioengineers}, 3.ª ed., 2018.

§5.4, \textbf{Using The Impulse Response To Find A Systems Output To Any
Input---Convolution} (página 219 del visor PDF); §6.7, \textbf{The
Transfer Function And The Fourier Transform} (página 287 del visor PDF);
§7.5, \textbf{The Laplace Domain, The Frequency Domain, And The Time
Domain} (página 330 del visor PDF); §8.2.3, \textbf{Transformation From
The Z-Domain To The Frequency Domain} (página 353 del visor PDF).

\textbf{Correspondencia:} Síntesis o encuadre que reúne varias
secciones; no corresponde a un único apartado del libro.

\textbf{Lectura:} use estas secciones como ruta de preparación y
consulta. La bibliografía aparece al final del bloque.
\end{frame}

\begin{frame}{155. Matriz de selección}
\phantomsection\label{matriz-de-selecciuxf3n}
\textbf{Libro guía:} John Semmlow, \emph{Circuits, Signals, and Systems
for Bioengineers}, 3.ª ed., 2018.

§5.4, \textbf{Using The Impulse Response To Find A Systems Output To Any
Input---Convolution} (página 219 del visor PDF); §7.5, \textbf{The
Laplace Domain, The Frequency Domain, And The Time Domain} (página 330
del visor PDF); §8.4, \textbf{Linear Filters---Introduction} (página 358
del visor PDF).

\textbf{Correspondencia:} Síntesis o encuadre que reúne varias
secciones; no corresponde a un único apartado del libro.

\textbf{Lectura:} use estas secciones como ruta de preparación y
consulta. La bibliografía aparece al final del bloque.
\end{frame}

\begin{frame}{156. Lista de control para filtros}
\phantomsection\label{lista-de-control-para-filtros}
\textbf{Libro guía:} John Semmlow, \emph{Circuits, Signals, and Systems
for Bioengineers}, 3.ª ed., 2018.

§8.4.1, \textbf{Filter Properties} (página 360 del visor PDF); §8.4.2,
\textbf{Finite Impulse Response Versus Infinite Impulse Response Filter
Characteristics} (página 363 del visor PDF).

\textbf{Correspondencia:} Síntesis o encuadre que reúne varias
secciones; no corresponde a un único apartado del libro.

\textbf{Lectura:} use estas secciones como ruta de preparación y
consulta. La bibliografía aparece al final del bloque.
\end{frame}

\begin{frame}{157. Lista de control para wavelets}
\phantomsection\label{lista-de-control-para-wavelets}
\textbf{Libro guía:} John Semmlow, \emph{Circuits, Signals, and Systems
for Bioengineers}, 3.ª ed., 2018.

§4.6, \textbf{Time--Frequency Analysis} (página 205 del visor PDF);
§8.4, \textbf{Linear Filters---Introduction} (página 358 del visor PDF).

\textbf{Correspondencia:} No existe una sección específica del libro
guía para el tema completo de esta diapositiva. Las secciones anteriores
son antecedentes, no referencias directas. Complemento directo: Semmlow
y Griffel (2014), §7.2 y §7.3.

\textbf{Lectura:} use estas secciones como ruta de preparación y
consulta. La bibliografía aparece al final del bloque.
\end{frame}

\begin{frame}{158. Lista de control para clasificación}
\phantomsection\label{lista-de-control-para-clasificaciuxf3n}
\textbf{Libro guía:} John Semmlow, \emph{Circuits, Signals, and Systems
for Bioengineers}, 3.ª ed., 2018.

§1.4.4, \textbf{Linear Versus Nonlinear Signals And Systems} (página 44
del visor PDF); §8.4, \textbf{Linear Filters---Introduction} (página 358
del visor PDF).

\textbf{Correspondencia:} No existe una sección específica del libro
guía para el tema completo de esta diapositiva. Las secciones anteriores
son antecedentes, no referencias directas. Complemento: Semmlow y
Griffel (2014), §16.2--16.3 y §16.6. La independencia por sujeto es una
ampliación docente.

\textbf{Lectura:} use estas secciones como ruta de preparación y
consulta. La bibliografía aparece al final del bloque.
\end{frame}

\begin{frame}{159. Problema integrador}
\phantomsection\label{problema-integrador}
\textbf{Libro guía:} John Semmlow, \emph{Circuits, Signals, and Systems
for Bioengineers}, 3.ª ed., 2018.

§5.5, \textbf{Applied Convolution---Basic Filters} (página 231 del visor
PDF); §8.4, \textbf{Linear Filters---Introduction} (página 358 del visor
PDF); §8.6, \textbf{Finite Impulse Response And Infinite Impulse
Response Filter Design Using The Signal Processing Toolbox} (página 382
del visor PDF).

\textbf{Correspondencia:} Síntesis o encuadre que reúne varias
secciones; no corresponde a un único apartado del libro. La parte
wavelet se apoya en Semmlow y Griffel (2014), §7.2--7.3; clasificación,
§16.2--16.3.

\textbf{Lectura:} use estas secciones como ruta de preparación y
consulta. La bibliografía aparece al final del bloque.
\end{frame}

\begin{frame}{160. Referencias y alcance}
\phantomsection\label{referencias-y-alcance-2}
\textbf{Libro guía:} John Semmlow, \emph{Circuits, Signals, and Systems
for Bioengineers}, 3.ª ed., 2018.

§5.4, \textbf{Using The Impulse Response To Find A Systems Output To Any
Input---Convolution} (página 219 del visor PDF); §6.4, \textbf{The
Transfer Function} (página 255 del visor PDF); §7.4, \textbf{Nonzero
Initial Conditions---Initial And Final Value Theorems} (página 326 del
visor PDF); §8.4, \textbf{Linear Filters---Introduction} (página 358 del
visor PDF).

\textbf{Correspondencia:} Síntesis o encuadre que reúne varias
secciones; no corresponde a un único apartado del libro.

\textbf{Lectura:} use estas secciones como ruta de preparación y
consulta. La bibliografía aparece al final del bloque.
\end{frame}

\begin{frame}{161. Cierre: criterio de dominio}
\phantomsection\label{cierre-criterio-de-dominio-2}
\textbf{Libro guía:} John Semmlow, \emph{Circuits, Signals, and Systems
for Bioengineers}, 3.ª ed., 2018.

§5.4, \textbf{Using The Impulse Response To Find A Systems Output To Any
Input---Convolution} (página 219 del visor PDF); §6.4, \textbf{The
Transfer Function} (página 255 del visor PDF); §7.4, \textbf{Nonzero
Initial Conditions---Initial And Final Value Theorems} (página 326 del
visor PDF); §8.4, \textbf{Linear Filters---Introduction} (página 358 del
visor PDF).

\textbf{Correspondencia:} Síntesis o encuadre que reúne varias
secciones; no corresponde a un único apartado del libro.

\textbf{Lectura:} use estas secciones como ruta de preparación y
consulta. La bibliografía aparece al final del bloque.
\end{frame}

\begin{frame}{162. Referencias}
\phantomsection\label{referencias-2}
\textbf{Libro guía:} John Semmlow, \emph{Circuits, Signals, and Systems
for Bioengineers}, 3.ª ed., 2018.

§5.4, \textbf{Using The Impulse Response To Find A Systems Output To Any
Input---Convolution} (página 219 del visor PDF); §6.4, \textbf{The
Transfer Function} (página 255 del visor PDF); §7.4, \textbf{Nonzero
Initial Conditions---Initial And Final Value Theorems} (página 326 del
visor PDF); §8.4, \textbf{Linear Filters---Introduction} (página 358 del
visor PDF).

\textbf{Correspondencia:} Síntesis o encuadre que reúne varias
secciones; no corresponde a un único apartado del libro.

\textbf{Lectura:} use estas secciones como ruta de preparación y
consulta. La bibliografía aparece al final del bloque.
\end{frame}




\end{document}
